% Native statements and proofs with reviewed registry prose.
\expandafter\def\csname statadic@statement@STATADIC-MMD-002\endcsname{For finite families of vectors $x_i,y_j$ in a real inner
  product space, real weights $w_i,v_j$ of equal total mass, and any reference
  $r$, the linear-kernel V-statistic satisfies
  \[
  \sum_{i,k}w_iw_k\langle x_i,x_k\rangle+
  \sum_{j,l}v_jv_l\langle y_j,y_l\rangle-
   2\sum_{i,j}w_iv_j\langle x_i,y_j\rangle=
  \left\|\sum_i w_i(x_i-r)-\sum_jv_j(y_j-r)\right\|^2.
  \]}
\expandafter\def\csname statadic@proof@STATADIC-MMD-002\endcsname{Equal total mass cancels the reference. Expand the inner
  product of the difference of weighted sums and distribute both finite sums.}
\expandafter\def\csname statadic@assumptions@STATADIC-MMD-002\endcsname{\begin{itemize}
\item Finite index types and real weights with equal total mass
\item Vectors in a real inner product space and any shared reference
\end{itemize}}
\expandafter\def\csname statadic@diagnostics@STATADIC-MMD-002\endcsname{\begin{itemize}
\item For weighted\_mmd\_squared\_cross, select its linear kernel and check nonnegative masses, equal total mass, and a common centering reference across Gram blocks. The implementation masks nonpositive entries. Signed or differently normalized weights need a separate comparison.
\end{itemize}}
\expandafter\def\csname statadic@statement@STATADIC-MMD-001\endcsname{On a finite probability law, let real pair arrays represent
  groups of sizes $m,n\geq2$. Assume common within-diagonal means
  $\delta_X,\delta_Y$, within off-diagonal means
  $\mu_X,\mu_Y$, and cross mean $\mu_{XY}$. The off-diagonal U estimator with
  all cross terms has expectation $\mu_X+\mu_Y-2\mu_{XY}$.}
\expandafter\def\csname statadic@proof@STATADIC-MMD-001\endcsname{Linearity of finite expectation and the $m(m-1)$ and $n(n-1)$
  off-diagonal counts give the result; see Gretton et al., Lemma 6, eq. (3),
  p. 728, and its empirical-estimator proof step, p. 729.
  The sampling assumptions that imply these moments are not derived here.}
\expandafter\def\csname statadic@assumptions@STATADIC-MMD-001\endcsname{\begin{itemize}
\item Normalized nonnegative finite probability law and exact real-valued pair arrays
\item Both sample index types have at least two elements
\item Within each group, supplied common diagonal and off-diagonal pair expectations
\item Supplied common expectation for every cross-group pair
\end{itemize}}
\expandafter\def\csname statadic@diagnostics@STATADIC-MMD-001\endcsname{\begin{itemize}
\item Identify the U-statistic path in mmd\_squared or mmd\_squared\_matrix and ensure at least two observations in each sample. Record the sampling design and whether within-group self-pairs were excluded.
\item The theorem supplies its within- and cross-pair expectations as hypotheses. If observations are clustered, paired across groups or the kernel was selected on these observations, derive the required pair expectations for that design before reporting unbiasedness.
\end{itemize}}
\expandafter\def\csname statadic@statement@STATADIC-MOMENTS-001\endcsname{Let $\Omega$ be finite and let $P$ have nonnegative weights
  $p_\omega$ summing to one. For a real estimator $f:\Omega\to\mathbb R$ and
  target $\theta\in\mathbb R$, write
  \[
  \mathbb E_P[g]=\sum_{\omega\in\Omega}p_\omega g(\omega),\quad
  m=\mathbb E_P[f],\quad \operatorname{Bias}=m-\theta,\quad
  \operatorname{Var}=\mathbb E_P[(f-m)^2].
  \]
  Then $\operatorname{MSE}=\mathbb E_P[(f-\theta)^2]$ satisfies
  \[
  \operatorname{MSE}=\operatorname{Var}+\operatorname{Bias}^2.
  \]}
\expandafter\def\csname statadic@proof@STATADIC-MOMENTS-001\endcsname{Following proof P5, equations MSE-ref1 and MSE-ref2, insert and
  subtract $m$ and expand:
  \[
  (f-\theta)^2=(f-m)^2+2(m-\theta)(f-m)+(m-\theta)^2.
  \]
  Finite expectation is linear, $\mathbb E_P[f-m]=m-m=0$, and normalized
  weights leave constants unchanged. Taking expectations therefore eliminates
  the cross term and yields $\operatorname{MSE}=\operatorname{Var}+
  \operatorname{Bias}^2$.}
\expandafter\def\csname statadic@assumptions@STATADIC-MOMENTS-001\endcsname{\begin{itemize}
\item Normalized nonnegative finite probability law
\item Real estimator and target
\end{itemize}}
\expandafter\def\csname statadic@diagnostics@STATADIC-MOMENTS-001\endcsname{\begin{itemize}
\item This identity has no registered Python fit. When using it in a simulation, specify the target, sampling law and estimator, verify nonnegative weights summing to one, and compute mean, variance and squared error with the same weights. Normalize or reconcile the denominators before comparing the decomposition.
\end{itemize}}
\expandafter\def\csname statadic@statement@STATADIC-LINEAR-004\endcsname{Let $X\in\mathbb R^{n\times p}$, $t,\beta\in\mathbb R^p$,
  and $Y=X\beta+\varepsilon$ on a finite probability space. Assume
  $\mathbb E\varepsilon=0$, $\mathbb E(\varepsilon\varepsilon^T)=\sigma^2I_n$
  with $\sigma^2>0$, and choose $a,b\in\mathbb R^n$ with
  \[ X^Ta=X^Tb=t,\qquad b\in\operatorname{col}(X). \]
  Then the supplied contrast weights $b$ satisfy
  \[ \mathbb E(b^TY)=t^T\beta,\qquad
     \operatorname{Var}(b^TY)\leq\operatorname{Var}(a^TY). \]}
\expandafter\def\csname statadic@proof@STATADIC-LINEAR-004\endcsname{Stępniak (2010), Theorem 4.1, pp. 4--5, decomposes a
  competitor into the least-squares component and an orthogonal difference.
  Here $X^T(a-b)=0$ and $b\in\operatorname{col}(X)$ give $b^T(a-b)=0$;
  spherical moments therefore yield the variance gap
  \[ \operatorname{Var}(a^TY)-\operatorname{Var}(b^TY)
     =\sigma^2\sum_i(a_i-b_i)^2\geq0. \]
  This conditional statement assumes the displayed weights; their existence
  and least-squares construction are separate results.}
\expandafter\def\csname statadic@assumptions@STATADIC-LINEAR-004\endcsname{\begin{itemize}
\item Finite probability law
\item Fixed design, zero-mean errors and spherical second moments
\item Strictly positive variance
\item Competitor reproduces the scalar target for every parameter vector
\item The candidate weights are supplied in the design column space and reproduce the chosen contrast.
\end{itemize}}
\expandafter\def\csname statadic@diagnostics@STATADIC-LINEAR-004\endcsname{\begin{itemize}
\item For the claimed contrast t, check X\ensuremath{{}^T}a=t for every competing linear estimator and verify that the claimed optimal weights lie in the design column space. Document zero-mean spherical errors on the finite law before using the variance ordering; a least-squares residual orthogonality check is only algebra.
\end{itemize}}
\expandafter\def\csname statadic@statement@STATADIC-LINEAR-005\endcsname{Let $X\in\mathbb R^{n\times p}$ have invertible $X^TX$,
  $t,\beta\in\mathbb R^p$, and $Y=X\beta+\varepsilon$ on a finite probability
  space, with $\mathbb E\varepsilon=0$ and
  $\mathbb E(\varepsilon\varepsilon^T)=\sigma^2 I_n$, $\sigma^2>0$. Put
  \[ \hat\beta=(X^TX)^{-1}X^TY. \]
  For any $a\in\mathbb R^n$ satisfying $X^Ta=t$, the exact-real OLS contrast
  satisfies
  \[ \mathbb E(t^T\hat\beta)=t^T\beta,\qquad
     \operatorname{Var}(t^T\hat\beta)\leq\operatorname{Var}(a^TY). \]}
\expandafter\def\csname statadic@proof@STATADIC-LINEAR-005\endcsname{As in Stępniak (2010), Theorem 4.1, pp. 4--5,
  set $b=X(X^TX)^{-1}t$. The Gram construction gives $X^Tb=t$ and
  $b\in\operatorname{col}(X)$.
  Hence $b^T(a-b)=0$, so the spherical moment assumptions imply
  \[ \operatorname{Var}(a^TY)-\operatorname{Var}(b^TY)
     =\sigma^2\sum_i(a_i-b_i)^2\geq0. \]
  The normal equations identify $b^TY=t^T\hat\beta$.}
\expandafter\def\csname statadic@assumptions@STATADIC-LINEAR-005\endcsname{\begin{itemize}
\item Finite probability law
\item Fixed design, zero-mean errors and spherical second moments
\item Strictly positive variance
\item Competitor reproduces the scalar target for every parameter vector
\item The Gram matrix X\ensuremath{{}^T}X has an invertible determinant.
\end{itemize}}
\expandafter\def\csname statadic@diagnostics@STATADIC-LINEAR-005\endcsname{\begin{itemize}
\item Inspect LinearFit.input\_valid, rank, identified and valid. Require the fitted design to retain all coefficient directions used in the contrast; obtain an estimable contrast or revise the design on rank failure.
\item Compare X\ensuremath{{}^T}(y\ensuremath{-}X\ensuremath{\hat\beta}) with zero at the chosen numerical precision. This checks the solved normal equations; spherical error second moments and conditional mean zero still require a sampling argument.
\end{itemize}}
\expandafter\def\csname statadic@statement@STATADIC-ENERGY-001\endcsname{Let $c_i$ be real weights on a finite set with
  $\sum_i c_i=0$. Suppose the real arrays $d$ and $k$ satisfy
  \[
    d_{ij}=k_{ii}+k_{jj}-2k_{ij}.
  \]
  Then the signed quadratic forms satisfy
  \[
    -\sum_i\sum_j c_i c_j d_{ij}
      =2\sum_i\sum_j c_i c_j k_{ij}.
  \]
  This algebraic identity does not assume kernel positivity.}
\expandafter\def\csname statadic@proof@STATADIC-ENERGY-001\endcsname{Substitute the generating relation. The two diagonal terms vanish
  because $\sum_i c_i=0$, leaving
  \[
    -\sum_{i,j}c_i c_j(k_{ii}+k_{jj}-2k_{ij})
      =2\sum_{i,j}c_i c_j k_{ij}.
  \]}
\expandafter\def\csname statadic@assumptions@STATADIC-ENERGY-001\endcsname{\begin{itemize}
\item Finite index type
\item Signed real coefficients with total zero
\item d(i,j) = k(i,i) + k(j,j) \ensuremath{-} 2k(i,j)
\end{itemize}}
\expandafter\def\csname statadic@diagnostics@STATADIC-ENERGY-001\endcsname{\begin{itemize}
\item Record the actual ground passed to energy\_distance and verify nonnegative normalized empirical masses on the same support convention. A small computed value alone does not identify a distribution.
\item For an energy–kernel comparison, check on the pooled support that d(i,j)=k(i,i)+k(j,j)\ensuremath{-}2k(i,j) and that the signed law difference sums to zero. Without that generating relation the cancellation does not identify the selected kernel discrepancy.
\end{itemize}}
\expandafter\def\csname statadic@statement@STATADIC-ENERGY-002\endcsname{Let $p_i,q_i$ be nonnegative weights on a finite set, each
  summing to one. Let $k$ be symmetric and positive semidefinite, and suppose
  $d_{ij}=k_{ii}+k_{jj}-2k_{ij}$. Define the unrooted raw energy and kernel form by
  \[
    E_d(p,q)=2\sum_{i,j}p_iq_jd_{ij}
      -\sum_{i,j}p_ip_jd_{ij}-\sum_{i,j}q_iq_jd_{ij},
    \qquad Q_k(c)=\sum_{i,j}c_ic_jk_{ij}.
  \]
  Then
  \[
    E_d(p,q)=2Q_k(p-q)\geq 0.
  \]
  The equality uses the stated generating convention; no square root is taken.}
\expandafter\def\csname statadic@proof@STATADIC-ENERGY-002\endcsname{Set $c_i=p_i-q_i$. The probability weights have equal mass,
  so $\sum_i c_i=0$. Symmetry gives $E_d(p,q)=-\sum_{i,j}c_ic_jd_{ij}$.
  Apply the signed energy--kernel identity and positive semidefiniteness:
  \[
    E_d(p,q)=-\sum_{i,j}c_ic_jd_{ij}=2Q_k(c)\geq 0.
  \]}
\expandafter\def\csname statadic@assumptions@STATADIC-ENERGY-002\endcsname{\begin{itemize}
\item Finite index type
\item Two finite probability laws with nonnegative weights summing to one
\item Symmetric kernel whose quadratic form is nonnegative for every real coefficient vector
\item Ground generated by that kernel
\end{itemize}}
\expandafter\def\csname statadic@diagnostics@STATADIC-ENERGY-002\endcsname{\begin{itemize}
\item Record the actual ground passed to energy\_distance and verify nonnegative normalized empirical masses on the same support convention. A small computed value alone does not identify a distribution.
\item Check the shared generating relation and PSD of the pooled kernel on the stated support; if PSD is assumed through a domain law, record its justification. Squaring, rooting or clamping the result cannot supply a missing PSD premise.
\end{itemize}}
\expandafter\def\csname statadic@statement@STATADIC-RANDOMIZATION-001\endcsname{Let a finite group $G$ act on a finite sample space with
  probability law $P$ satisfying $P(gx)=P(x)$ for all $g,x$. For a real statistic
  $T$, count upper-tail ties in
  \[
    R(x)=\sum_{g\in G}\mathbf{1}\{T(gx)\ge T(x)\}.
  \]
  For $\alpha\ge0$, rejection when $R(x)\le\alpha |G|$ satisfies
  \[
    P\{R(X)\le\alpha |G|\}\le\alpha.
  \]
  This is a finite-law rank formulation; equivalence to the paper's
  order-statistic rule remains to be formalized.}
\expandafter\def\csname statadic@proof@STATADIC-RANDOMIZATION-001\endcsname{Group multiplication permutes the transformation indices. On
  each orbit, the number of rejected transformations is at most $\alpha|G|$:
  choose a rejected statistic of minimum value and include all rejected
  transformations in its tie-inclusive upper tail. Invariance and finite-sum
  exchange then give
  \[
    P\{R(X)\le\alpha|G|\}
      =\mathbb{E}\left[\frac{1}{|G|}\sum_{g\in G}
        \mathbf{1}\{R(gX)\le\alpha|G|\}\right]\le\alpha.
  \]}
\expandafter\def\csname statadic@assumptions@STATADIC-RANDOMIZATION-001\endcsname{\begin{itemize}
\item Finite probability-weighted data space
\item Finite group action preserving data-law weights
\item Real statistic; upper-tail ranks include ties
\item Nonnegative alpha
\end{itemize}}
\expandafter\def\csname statadic@diagnostics@STATADIC-RANDOMIZATION-001\endcsname{\begin{itemize}
\item dispersion\_permutation\_null generates a requested number of residual-permutation draws. For an exact full-group claim, verify the complete admissible transformation group, its null invariance and finite statistics including ties. A generated Monte Carlo schedule does not satisfy full enumeration simply by containing many draws.
\end{itemize}}
\expandafter\def\csname statadic@statement@STATADIC-RANDOMIZATION-002\endcsname{Let a finite group $G$ act on a finite space $\Omega$ with
  nonnegative weights $w$ summing to one and invariant under the action. Let
  $B\in\mathbb N$, $T:\Omega\to\mathbb R$, and $\alpha\ge0$. For each tuple
  $g\in G^B$, put
  \[
  p_T(\omega,g)=\frac{1+\#\{k<B:T(\omega)\le T(g_k\omega)\}}{B+1}.
  \]
  Then, for independent uniform draws $g$,
  \[
  \sum_{\omega\in\Omega}\sum_{g\in G^B}
  \frac{w(\omega)}{|G|^B}\mathbf1\{p_T(\omega,g)\le\alpha\}\le\alpha.
  \]}
\expandafter\def\csname statadic@proof@STATADIC-RANDOMIZATION-002\endcsname{Include the identity at index zero. For each index $j$, swap
  zero and $j$ and right-multiply every entry by the inverse of the original
  $j$th entry, while acting by that entry on the observation. This is an
  involution preserving the product weights. It changes the observed statistic
  to the $j$th statistic and permutes the sampled values. The inclusive
  upper-rank counting bound limits the number of rejecting indices to
  $\alpha(B+1)$. Averaging over indices yields the bound.}
\expandafter\def\csname statadic@assumptions@STATADIC-RANDOMIZATION-002\endcsname{\begin{itemize}
\item Finite data space with nonnegative weights summing to one
\item Finite group action preserving each data-law weight
\item Independent uniformly distributed group draws with replacement, independent of the data
\item Real-valued statistic, inclusive ties and nonnegative alpha
\end{itemize}}
\expandafter\def\csname statadic@diagnostics@STATADIC-RANDOMIZATION-002\endcsname{\begin{itemize}
\item Record the transformation group, inclusion of the observed statistic, the inclusive tie rule and independent uniform draws with replacement for the sampled theorem. permutation\_test\_kernel returns a NaN p-value when the observed or any required draw statistic is nonfinite. Resolve those undefined statistics before calibration.
\end{itemize}}
\expandafter\def\csname statadic@statement@STATADIC-MMD-003\endcsname{For finite vector families $x_i,y_j$ in a real inner product
  space and real weights $w_i,v_j$ of equal mass, the linear V expression
  $\sum_{i,k}w_iw_k\langle x_i,x_k\rangle+\sum_{j,l}v_jv_l\langle y_j,y_l\rangle-2\sum_{i,j}w_iv_j\langle x_i,y_j\rangle$ is nonnegative.}
\expandafter\def\csname statadic@proof@STATADIC-MMD-003\endcsname{Apply the weighted mean reduction and nonnegativity of a square.}
\expandafter\def\csname statadic@assumptions@STATADIC-MMD-003\endcsname{\begin{itemize}
\item Finite index types and real weights with equal total mass
\item Vectors in a real inner product space and any shared reference
\end{itemize}}
\expandafter\def\csname statadic@diagnostics@STATADIC-MMD-003\endcsname{\begin{itemize}
\item For the linear-kernel weighted V statistic, check equal mass and compare the result with the squared norm of the difference of weighted feature means. Check positive-mass features for nonfinite entries before interpreting a NaN or a material negative output.
\end{itemize}}
\expandafter\def\csname statadic@statement@STATADIC-MMD-004\endcsname{For a symmetric positive semidefinite real kernel $k$ on
  the finite disjoint union $\iota\sqcup\kappa$ and real weights $w_i,v_j$,
  $\sum_{i,j}w_iw_jk(i,j)+\sum_{i,j}v_iv_jk(i,j)-2\sum_{i,j}w_iv_jk(i,j)\geq0$,
  where each index is embedded in its respective component of the union.}
\expandafter\def\csname statadic@proof@STATADIC-MMD-004\endcsname{Split the quadratic form of the signed pooled weights into four
  blocks, use symmetry on the reverse cross block, and apply PSD.}
\expandafter\def\csname statadic@assumptions@STATADIC-MMD-004\endcsname{\begin{itemize}
\item Finite disjoint union of sample indices
\item A symmetric real kernel that satisfies KernelPSD on the pooled support
\item Arbitrary real weights
\end{itemize}}
\expandafter\def\csname statadic@diagnostics@STATADIC-MMD-004\endcsname{\begin{itemize}
\item Check that the same symmetric PSD kernel generated both within blocks and the cross block. Positive eigenvalues of the separate within blocks do not certify the pooled block. A finite pooled Gram diagnostic concerns this sample; record the kernel law separately for population use.
\end{itemize}}
\expandafter\def\csname statadic@statement@STATADIC-MMD-005\endcsname{For each of the two samples $(1,-1)$, the linear kernel
  $k(a,b)=ab$ and the unique outcome in $\mathrm{Unit}$, the finite U statistic
  $\operatorname{mmdU}$ equals $-2$.}
\expandafter\def\csname statadic@proof@STATADIC-MMD-005\endcsname{Expand the two off-diagonal sums and the cross sum.}
\expandafter\def\csname statadic@assumptions@STATADIC-MMD-005\endcsname{\begin{itemize}
\item Two identical samples (1, \ensuremath{-}1), real linear kernel, and the unique Unit outcome
\end{itemize}}
\expandafter\def\csname statadic@diagnostics@STATADIC-MMD-005\endcsname{\begin{itemize}
\item Identify the U-statistic path in mmd\_squared or mmd\_squared\_matrix and ensure at least two observations in each sample. Record the sampling design and whether within-group self-pairs were excluded.
\item Keep signed outputs from the U estimator, including negative realizations. If the downstream analysis requires a rooted nonnegative quantity, choose an estimand for which that operation is justified.
\end{itemize}}
\expandafter\def\csname statadic@statement@STATADIC-MMD-006\endcsname{Given finite index sets $\iota,\kappa,\alpha,\beta$, a real
  kernel $k$ on $(\iota\sqcup\alpha)\sqcup(\kappa\sqcup\beta)$ and real weights
  $w,v$, let $\tilde w=(w,0)$ and $\tilde v=(v,0)$. Then the V expression
  $\sum_{i,j}\tilde w_i\tilde w_jk(i,j)+\sum_{i,j}\tilde v_i\tilde v_jk(i,j)-2\sum_{i,j}\tilde w_i\tilde v_jk(i,j)$
  equals the same expression restricted to $\iota$ and $\kappa$ using $w,v$.}
\expandafter\def\csname statadic@proof@STATADIC-MMD-006\endcsname{Split each finite sum over its disjoint union; padded terms vanish.}
\expandafter\def\csname statadic@assumptions@STATADIC-MMD-006\endcsname{\begin{itemize}
\item Finite original and padding index types
\item An arbitrary real kernel on the padded pooled support
\item Original real weights and exactly zero weights on both padding blocks
\end{itemize}}
\expandafter\def\csname statadic@diagnostics@STATADIC-MMD-006\endcsname{\begin{itemize}
\item Review packed lengths and masses on the weighted V path. Padded slots must have exactly zero mass, and active coordinates and kernel entries must be finite. Compare the result with the corresponding unpadded groups if a padding mask or pair-chunk schedule changes.
\end{itemize}}
\expandafter\def\csname statadic@statement@STATADIC-LINEAR-006\endcsname{Let $W=\operatorname{diag}(w_i)$ with $w_i>0$ and let
  $X^TWX$ be invertible. On a finite probability space suppose
  $Y=X\beta+\varepsilon$, $\mathbb E\varepsilon=0$, and
  $\mathbb E(\varepsilon_i\varepsilon_j)=\mathbf 1_{i=j}\sigma^2/w_i$
  for $\sigma^2>0$. Put $\hat\beta_W=(X^TWX)^{-1}X^TWY$,
  $\widetilde X=W^{1/2}X$, and $\widetilde Y=W^{1/2}Y$.
  For any $a$ with $\widetilde X^Ta=t$,
  \[\mathbb E(t^T\hat\beta_W)=t^T\beta,\qquad
  \operatorname{Var}(t^T\hat\beta_W)\leq\operatorname{Var}(a^T\widetilde Y).\]}
\expandafter\def\csname statadic@proof@STATADIC-LINEAR-006\endcsname{Whitening gives spherical errors and the same inverse-Gram
  coefficient vector. Apply the full-rank finite-law OLS contrast theorem to
  the whitened design and observations.}
\expandafter\def\csname statadic@assumptions@STATADIC-LINEAR-006\endcsname{\begin{itemize}
\item Finite probability law and fixed design
\item Strictly positive deterministic diagonal weights and positive error scale
\item Invertible weighted Gram determinant
\item Zero error means and diagonal pair moments \ensuremath{\sigma}\ensuremath{{}^2}/w\ensuremath{{}_i}
\item Competitor reproduces the contrast on the whitened design
\end{itemize}}
\expandafter\def\csname statadic@diagnostics@STATADIC-LINEAR-006\endcsname{\begin{itemize}
\item Inspect LinearFit.relative\_precision, input\_valid, rank and valid. Record the fixed positive relative-precision weights used for fitting. Frequency weights or data-estimated variances need a different sampling argument.
\item Check the whitened design rank and the claimed diagonal covariance up to common scale. If weights reflect replication or clustering rather than known diagonal precision, do not report the WLS optimality claim.
\end{itemize}}
\expandafter\def\csname statadic@statement@STATADIC-LINEAR-007\endcsname{Let $X$ be a fixed real matrix and let $Y=s+\varepsilon$
  on a finite probability space, where $\mathbb E\varepsilon=0$ and
  $\mathbb E(\varepsilon_i\varepsilon_j)=\mathbf 1_{i=j}\sigma_i^2$.
  With $\hat\beta=(X^TX)^{-1}X^TY$ and $D=\operatorname{diag}(\sigma_i^2)$,
  \[\operatorname{Var}(t^T\hat\beta)=
  t^T(X^TX)^{-1}X^TDX(X^TX)^{-1}t.\]
  Matrix inverses here use Lean's total nonsingular inverse.}
\expandafter\def\csname statadic@proof@STATADIC-LINEAR-007\endcsname{Expand the finite variance of the linear contrast, use the
  diagonal pair moments, then rearrange the resulting sum into the sandwich.
  The deterministic signal cancels under centering.}
\expandafter\def\csname statadic@assumptions@STATADIC-LINEAR-007\endcsname{\begin{itemize}
\item Finite probability law
\item Fixed real design and deterministic signal
\item Zero error means and supplied diagonal pair moments
\end{itemize}}
\expandafter\def\csname statadic@diagnostics@STATADIC-LINEAR-007\endcsname{\begin{itemize}
\item Inspect LinearFit.input\_valid, rank, identified and valid. Require the fitted design to retain all coefficient directions used in the contrast; obtain an estimable contrast or revise the design on rank failure.
\item The variance identity uses true error variances and zero cross-error second moments. Treat sample residual squares in hc\_covariance as plug-in estimates and record the dependence assumptions. Shared units or time clusters may invalidate the diagonal-error comparison.
\end{itemize}}
\expandafter\def\csname statadic@statement@STATADIC-LINEAR-008\endcsname{Let $X$ have invertible $X^TX$ and let $Y=X\beta+\varepsilon$
  on a finite probability space, with $\mathbb E\varepsilon=0$ and
  $\mathbb E(\varepsilon_i\varepsilon_j)=\mathbf 1_{i=j}\sigma^2$.
  Put $\hat\beta=(X^TX)^{-1}X^TY$ and
  $h_i=[X(X^TX)^{-1}X^T]_{ii}$. Then, for each $i$,
  \[\mathbb E\bigl(Y_i-(X\hat\beta)_i\bigr)^2=\sigma^2(1-h_i).\]}
\expandafter\def\csname statadic@proof@STATADIC-LINEAR-008\endcsname{The residual is the corresponding row of $(I-H)\varepsilon$.
  Symmetry and idempotence of $H$ make $I-H$ a symmetric projection, whose
  row sum of squares is its diagonal. Apply spherical second moments.}
\expandafter\def\csname statadic@assumptions@STATADIC-LINEAR-008\endcsname{\begin{itemize}
\item Finite probability law
\item Fixed full-column-rank design
\item Zero error means and spherical second moments
\end{itemize}}
\expandafter\def\csname statadic@diagnostics@STATADIC-LINEAR-008\endcsname{\begin{itemize}
\item Inspect LinearFit.input\_valid, rank, identified and valid. Require the fitted design to retain all coefficient directions used in the contrast; obtain an estimable contrast or revise the design on rank failure.
\item Check df\_resid, the fitted residuals and diagonal leverage when interpreting residual squares. The residual moment is derived under spherical errors; large leverage reduces expected raw residual squares. When leverage approaches one, scrutinize any correction dividing by one minus leverage.
\end{itemize}}
\expandafter\def\csname statadic@statement@STATADIC-ENERGY-003\endcsname{Let $x_1,\ldots,x_n$ lie in a real inner product space and let
  $p_i,q_i$ be nonnegative weights each summing to one. Then
  \[
    2\sum_{i,j}p_iq_j\lVert x_i-x_j\rVert-
    \sum_{i,j}p_ip_j\lVert x_i-x_j\rVert-
    \sum_{i,j}q_iq_j\lVert x_i-x_j\rVert\geq 0.
  \]}
\expandafter\def\csname statadic@proof@STATADIC-ENERGY-003\endcsname{The signed weights $p_i-q_i$ sum to zero. Negative type of the
  norm makes their distance quadratic nonpositive; symmetry identifies its
  negative with the stated energy.}
\expandafter\def\csname statadic@assumptions@STATADIC-ENERGY-003\endcsname{\begin{itemize}
\item A finite family of points in a real inner product space
\item Two finite probability laws over the same finite index family
\end{itemize}}
\expandafter\def\csname statadic@diagnostics@STATADIC-ENERGY-003\endcsname{\begin{itemize}
\item Record the actual ground passed to energy\_distance and verify nonnegative normalized empirical masses on the same support convention. A small computed value alone does not identify a distribution.
\item Use Euclidean distance at exponent one in a real inner-product representation. Verify that the data sent to the kernel use those coordinates and scaling. For other grounds or exponents, supply the applicable negative-type result before asserting nonnegativity.
\end{itemize}}
\expandafter\def\csname statadic@statement@STATADIC-SPLINE-001\endcsname{For neighboring unit cells of a uniform cubic cardinal
  spline, the extensions of their polynomials differ at displacement $h$
  from their join by $(-a+4b-6c+4d-e)h^3/6$, where $a,b,c,d,e$
  are the consecutive control values. Their Taylor coefficients through
  order two therefore agree, including at a cyclic join.}
\expandafter\def\csname statadic@proof@STATADIC-SPLINE-001\endcsname{Evaluate the two cubic polynomials at $1+h$ and $h$
  respectively. Expand and collect powers of $h$. The terms of degrees
  zero, one, and two cancel.}
\expandafter\def\csname statadic@assumptions@STATADIC-SPLINE-001\endcsname{\begin{itemize}
\item Real-valued coefficients and displacement in units of one knot interval
\item Adjacent cells share three consecutive controls, with indices wrapping at the periodic boundary
\end{itemize}}
\expandafter\def\csname statadic@diagnostics@STATADIC-SPLINE-001\endcsname{\begin{itemize}
\item For CyclicSpline.transform, record the declared period, degrees of freedom and phase units. Its constructor requires a positive physical period and enough wrapped neighbours. The theorem applies to adjacent cells of the same uniform cyclic grid. Check the period and grid if a fitted curve jumps at the wrap.
\end{itemize}}
\expandafter\def\csname statadic@statement@STATADIC-DIFFUSION-001\endcsname{On a finite roster, an off-diagonal exponential kernel normalised
  by its row sum has nonnegative entries, zero diagonal and unit row sum for every
  row with at least one peer.}
\expandafter\def\csname statadic@proof@STATADIC-DIFFUSION-001\endcsname{The peer contributes strictly positive exponential mass and every
  term is nonnegative. Dividing by the positive row total preserves signs; summing
  the quotients divides that total by itself.}
\expandafter\def\csname statadic@assumptions@STATADIC-DIFFUSION-001\endcsname{\begin{itemize}
\item Finite real pairwise exponent costs on a common finite roster
\item The chosen row has a distinct peer; the scientific operator uses positive bandwidth
\end{itemize}}
\expandafter\def\csname statadic@diagnostics@STATADIC-DIFFUSION-001\endcsname{\begin{itemize}
\item Before calling rbf\_weights, check a square aligned roster of at least two units, finite off-diagonal costs and a finite positive bandwidth. Inspect zero diagonal, finite weights and unit row sums after normalization; an isolate or nonfinite row invalidates its interpretation as a transition distribution.
\end{itemize}}
\expandafter\def\csname statadic@statement@STATADIC-DIFFUSION-002\endcsname{Multiplying every finite exponent cost and its bandwidth by
  the same positive scalar leaves each normalised diffusion weight unchanged.}
\expandafter\def\csname statadic@proof@STATADIC-DIFFUSION-002\endcsname{Cancel the common factor in the exponent quotient. Both the
  unnormalised affinities and their row sums are unchanged.}
\expandafter\def\csname statadic@assumptions@STATADIC-DIFFUSION-002\endcsname{\begin{itemize}
\item Finite real exponent costs and a common positive rescaling factor
\item Bandwidth is on the exponent-cost scale; the scientific operator uses positive bandwidth
\end{itemize}}
\expandafter\def\csname statadic@diagnostics@STATADIC-DIFFUSION-002\endcsname{\begin{itemize}
\item The scale factor acts on the cost in the exponent. Record how raw distances became costs and rescale the bandwidth by the same positive factor. For squared distances, converting raw units by c rescales costs by c\ensuremath{{}^2}. A row-sum check alone cannot detect a bandwidth-unit mismatch.
\end{itemize}}
\expandafter\def\csname statadic@statement@STATADIC-SPATIAL-001\endcsname{The log absolute determinant of the block diagonal map
  with period blocks $B_t A_t$ equals the sum over periods of the log absolute
  determinants of the lag and error filters, when each is nonsingular.}
\expandafter\def\csname statadic@proof@STATADIC-SPATIAL-001\endcsname{The determinant of a block diagonal matrix is the product of
  the block determinants. Apply multiplicativity to each block and sum logs
  of their nonzero determinants. No commutation of the filters is needed.}
\expandafter\def\csname statadic@assumptions@STATADIC-SPATIAL-001\endcsname{\begin{itemize}
\item Finite period and unit axes on a common fully observed roster
\item The filters A\_t and B\_t are square real matrices with nonzero determinants
\item Each period maps its own outcomes to innovations with linear part B\_t A\_t
\end{itemize}}
\expandafter\def\csname statadic@diagnostics@STATADIC-SPATIAL-001\endcsname{\begin{itemize}
\item For graph\_profile, verify complete outcomes on the common roster, period alignment of lag/error graphs, and nonsingularity of each period’s filters. Inspect finite log determinants and near-singular sensitivity. Use the observed covariance calculation when outcomes are missing.
\end{itemize}}
\expandafter\def\csname statadic@statement@STATADIC-PANEL-001\endcsname{For one occupied group let $M=\sum_i w_i>0$ and
  $m=\sum_i w_i x_i/M$. For every constant $a$,
  $\sum_i w_i(x_i-a)^2=\sum_i w_i(x_i-m)^2+M(m-a)^2$.
  Consequently the group mean is the unique least-squares intercept.}
\expandafter\def\csname statadic@proof@STATADIC-PANEL-001\endcsname{First, $\sum_i w_i(x_i-m)=\sum_i w_i x_i-Mm=0$.
  Substitute $x_i-a=(x_i-m)+(m-a)$ and expand the square.
  The cross term is $2(m-a)\sum_i w_i(x_i-m)=0$.
  The remaining excess $M(m-a)^2$ vanishes only at $a=m$.}
\expandafter\def\csname statadic@assumptions@STATADIC-PANEL-001\endcsname{\begin{itemize}
\item A finite set of rows with one entity or one time label and a scalar data column
\item The group's total weight is strictly positive. In the panel application every input precision is positive.
\end{itemize}}
\expandafter\def\csname statadic@diagnostics@STATADIC-PANEL-001\endcsname{\begin{itemize}
\item Check group IDs, positive finite weights, and positive total weight for every occupied group. panel\_group\_mean leaves unoccupied groups undefined. After one-way absorption verify weighted within-group residual sums near zero. For alternating two-way projections inspect convergence and the remaining within-design rank separately.
\end{itemize}}
\expandafter\def\csname statadic@statement@STATADIC-GLM-001\endcsname{Write $d_i=d\mu_i/d\eta_i\ne0$, $V_i>0$,
  $w_i=d_i^2/V_i$, and $z_i=\eta_i+(y_i-\mu_i)/d_i$ at the current iterate.
  For any design column $x$ and predictor increment $s$,
  $\sum_i x_i w_i(z_i-\eta_i-s_i)=
  \sum_i x_i(y_i-\mu_i)d_i/V_i-\sum_i x_i w_i s_i$.}
\expandafter\def\csname statadic@proof@STATADIC-GLM-001\endcsname{Subtract the old linear predictor from the working response.
  Multiplying $(y_i-\mu_i)/d_i$ by $d_i^2/V_i$ gives
  $(y_i-\mu_i)d_i/V_i$. Subtract the increment contribution $w_i s_i$,
  multiply by the design entry and sum. With $s=X\Delta$ this matches
  each weighted normal equation to the corresponding Fisher scoring equation.}
\expandafter\def\csname statadic@assumptions@STATADIC-GLM-001\endcsname{\begin{itemize}
\item Finite rows with real outcomes, predictors, current means and design entries
\item Nonzero supplied inverse-link derivatives and strictly positive supplied row variances at the current parameter
\item Unit dispersion for the score interpretation, with fixed offsets contributing zero predictor increment
\end{itemize}}
\expandafter\def\csname statadic@diagnostics@STATADIC-GLM-001\endcsname{\begin{itemize}
\item Inspect GLMFit.input\_valid, identified, rank, valid and solver.converged; examine extreme fitted means, near-zero link derivatives or variances and ill-conditioned working information. The working-score identity concerns one iterate. Numerical stationarity on its own cannot establish absence of separation or correctness of the mean/variance model.
\end{itemize}}
\expandafter\def\csname statadic@statement@STATADIC-LINEAR-009\endcsname{Let $X$ be fixed with $X^T X$ invertible, and let
  $Y=X\beta+\varepsilon$ with each $\varepsilon_i$ integrable and of mean zero.
  Under any probability law, $\mathbb E[\hat\beta_j]=\beta_j$ for every
  coefficient of $\hat\beta=(X^T X)^{-1}X^T Y$.}
\expandafter\def\csname statadic@proof@STATADIC-LINEAR-009\endcsname{Write $L=(X^TX)^{-1}X^T$. Invertibility gives $LX=I$,
  so $\hat\beta_j=\beta_j+\sum_i L_{ji}\varepsilon_i$.
  Integrate term by term over the finite row index. The constant integrates
  to $\beta_j$ and each integrable error has zero expectation.}
\expandafter\def\csname statadic@assumptions@STATADIC-LINEAR-009\endcsname{\begin{itemize}
\item Finite fixed design X with invertible Gram determinant and a fixed coefficient vector
\item Integrable, coordinatewise mean-zero errors on an arbitrary probability space
\end{itemize}}
\expandafter\def\csname statadic@diagnostics@STATADIC-LINEAR-009\endcsname{\begin{itemize}
\item Inspect LinearFit.input\_valid, rank, identified and valid. Require the fitted design to retain all coefficient directions used in the contrast; obtain an estimable contrast or revise the design on rank failure.
\item Document why the conditional error mean is zero at the fixed design and why first moments exist. Orthogonality of fitted residuals follows from the solve and cannot validate exogeneity.
\end{itemize}}
\expandafter\def\csname statadic@statement@STATADIC-LINEAR-010\endcsname{Let $X$ be fixed with $X^TX$ invertible. Let the response
  vector $Y$ have mean $X\beta$, square-integrable coordinates and covariance $vI$, $v\ge0$.
  For any contrast $t$ and any weights $a$ with $X^Ta=t$, put $b=X(X^TX)^{-1}t$.
  Then $\operatorname{Var}(a^TY)-\operatorname{Var}(t^T\hat\beta)
  = v\sum_i(a_i-b_i)^2\ge0$. In the linear mean model $a^TY$ and
  $t^T\hat\beta$ estimate the same contrast without bias.}
\expandafter\def\csname statadic@proof@STATADIC-LINEAR-010\endcsname{The OLS contrast equals $b^TY$, and $X^Tb=t$. Taking expectations
  gives $b^TX\beta=a^TX\beta=t^T\beta$.
  Hence $X^T(a-b)=0$ while $b$ lies in the column space of $X$, giving
  $b^T(a-b)=0$. Square integrability permits expansion of each variance as
  a finite double sum of covariances. The spherical covariance leaves only
  $v\sum_i a_i^2$ or $v\sum_i b_i^2$. The Pythagorean identity gives the gap.}
\expandafter\def\csname statadic@assumptions@STATADIC-LINEAR-010\endcsname{\begin{itemize}
\item Finite fixed design and invertible Gram determinant, with response mean X\ensuremath{\beta}
\item Square-integrable observations on a probability space, with covariance vI for a common nonnegative v
\item A scalar contrast t and competitor weights a reproducing it for every coefficient vector, X\ensuremath{{}^T}a=t
\end{itemize}}
\expandafter\def\csname statadic@diagnostics@STATADIC-LINEAR-010\endcsname{\begin{itemize}
\item Inspect LinearFit.input\_valid, rank, identified and valid. Require the fitted design to retain all coefficient directions used in the contrast; obtain an estimable contrast or revise the design on rank failure.
\item For BLUE comparisons, state the common error variance, absence of cross-error covariance and square integrability. Inspect residual patterns and clustering for reasons to revisit these assumptions; a passing rank check certifies only the design condition.
\end{itemize}}
\expandafter\def\csname statadic@statement@STATADIC-LINEAR-011\endcsname{Let positive precisions $w_i$ be fixed and
  $W=\operatorname{diag}(w)$. Assume $X^TWX$ is invertible, $E[Y]=X\beta$
  and $\operatorname{Cov}(Y)=vW^{-1}$, $v\ge0$, with square-integrable rows.
  For any $t$ and deterministic $a$ with $X^Ta=t$, the WLS contrast
  $t^T\hat\beta_W$ has expectation $t^T\beta$ and variance at most that of $a^TY$.}
\expandafter\def\csname statadic@proof@STATADIC-LINEAR-011\endcsname{Put $D=\operatorname{diag}(\sqrt{w_i})$. Transform the design
  and response to $DX,DY$, and transform the competitor to $D^{-1}a$.
  These changes preserve its contrast estimate and reproduction condition.
  The transformed covariance is $vI$. Apply the OLS variance ordering and use
  the weighted Gram identity to identify the whitened OLS fit with WLS.}
\expandafter\def\csname statadic@assumptions@STATADIC-LINEAR-011\endcsname{\begin{itemize}
\item Positive fixed relative precisions, a fixed design and invertible weighted Gram determinant
\item Square-integrable observations on a probability space with mean X\ensuremath{\beta} and covariance v times the inverse diagonal precision matrix, for v \ensuremath{\ge} 0
\item Arbitrary deterministic original-coordinate competitor a satisfying X\ensuremath{{}^T}a=t
\end{itemize}}
\expandafter\def\csname statadic@diagnostics@STATADIC-LINEAR-011\endcsname{\begin{itemize}
\item Inspect LinearFit.relative\_precision, input\_valid, rank and valid. Record the fixed positive relative-precision weights used for fitting. Frequency weights or data-estimated variances need a different sampling argument.
\item Justify the inverse relation between the true error variances and supplied precision weights, up to a common scale, and the absence of cross-row error covariance. A successful weighted solve also occurs under misspecified weights.
\end{itemize}}
\expandafter\def\csname statadic@statement@STATADIC-GMM-001\endcsname{Put $v=Wu$ and $e=u-\lambda v$. For any symmetric
  moment matrix $A$ on $n>0$ units,
  \[
    e^TAe/n=u^TAu/n-\lambda(2v^TAu/n)+\lambda^2 v^TAv/n .
  \]}
\expandafter\def\csname statadic@proof@STATADIC-GMM-001\endcsname{Expand both occurrences of $e$. The mixed products
  $u^TAv$ and $v^TAu$ agree by symmetry. Collect them and divide by $n$.
  This gives the two coefficient signs used in the sample moment system.}
\expandafter\def\csname statadic@assumptions@STATADIC-GMM-001\endcsname{\begin{itemize}
\item A finite, nonempty cross section and a real spatial lag matrix.
\item A symmetric real moment matrix and any candidate spatial-error coefficient.
\end{itemize}}
\expandafter\def\csname statadic@diagnostics@STATADIC-GMM-001\endcsname{\begin{itemize}
\item SpatialGMM.estimate returns moment\_matrices, innovations, moment\_derivative and the lambda used for the coefficient step. With the same residuals and candidate lambda, compare the direct filtered quadratic with its expanded polynomial. Retain the solver and boundary flags; an algebraic match does not identify lambda.
\end{itemize}}
\expandafter\def\csname statadic@statement@STATADIC-GMM-002\endcsname{Let $\varepsilon$ be a real vector on a finite probability
  space. Assume $E[\varepsilon_i\varepsilon_j]=0$ whenever $i\ne j$.
  If $A_{ii}=0$ for every $i$, then
  $E[\sum_{i,j}A_{ij}\varepsilon_i\varepsilon_j]=0$.}
\expandafter\def\csname statadic@proof@STATADIC-GMM-002\endcsname{Exchange the finite expectation with both index sums.
  Each diagonal summand vanishes by the hypothesis on $A$. Every off-diagonal
  summand vanishes by the corresponding raw second moment.}
\expandafter\def\csname statadic@assumptions@STATADIC-GMM-002\endcsname{\begin{itemize}
\item A finite probability law on innovation vectors and finitely many units.
\item Every off-diagonal raw second moment vanishes and every diagonal entry of the moment matrix is zero.
\end{itemize}}
\expandafter\def\csname statadic@diagnostics@STATADIC-GMM-002\endcsname{\begin{itemize}
\item Check which variance\_model was used and inspect the diagonal of the actual moment matrices. In heteroskedastic mode, off-diagonal raw innovation cross moments must vanish at the population parameter. Residual spatial dependence is evidence to investigate that premise; a trace-zero matrix alone cannot replace a zero diagonal under unequal variances.
\end{itemize}}
\expandafter\def\csname statadic@statement@STATADIC-GMM-003\endcsname{Write the weighted affine moments as $a+B\theta$,
  where $a=R\bar g_0$ and $B=R D$ for a fixed weight factor $R$.
  Suppose $B^TB$ is invertible and set
  $b=(B^TB)^{-1}B^T(-a)$. For every $h$,
  \[
    \|a+B(b+h)\|^2=\|a+Bb\|^2+\|Bh\|^2 .
  \]
  In particular the computed coefficient minimizes the quadratic criterion.}
\expandafter\def\csname statadic@proof@STATADIC-GMM-003\endcsname{The normal equations give $B^T(a+Bb)=0$. Expand
  the squared norm of $(a+Bb)+Bh$. Its mixed inner product is
  $h^TB^T(a+Bb)=0$. The remaining sum of squares is nonnegative.}
\expandafter\def\csname statadic@assumptions@STATADIC-GMM-003\endcsname{\begin{itemize}
\item Real weighted moment offset a and fixed weighted derivative matrix B.
\item The Gram matrix B\ensuremath{{}^T}B is invertible.
\end{itemize}}
\expandafter\def\csname statadic@diagnostics@STATADIC-GMM-003\endcsname{\begin{itemize}
\item Inspect weighted\_rank, first\_order\_residual, identified, valid and the retained effective\_weight on the fixed-weight affine fit. Increase scrutiny for nearly singular weighted derivatives. Distinguish a sample objective minimum from population identification or validity of a J reference law.
\end{itemize}}
\expandafter\def\csname statadic@statement@STATADIC-IV-001\endcsname{For instruments $H$ with nonsingular $H^TH$, write
  $P=H(H^TH)^{-1}H^T$ and $D=PX$ for the projected structural regressors.
  If $D^TD$ is nonsingular and $b=(D^TD)^{-1}D^Ty$, then
  $D^T(y-Xb)=0$. Here $X$ may include the endogenous column $Wy$.}
\expandafter\def\csname statadic@proof@STATADIC-IV-001\endcsname{Symmetry and idempotence of $P$ give
  $D^TD=X^TP^TPX=X^TP^T X=D^TX$. Substitute this identity
  into the normal equations $D^TD b=D^Ty$ and subtract.}
\expandafter\def\csname statadic@assumptions@STATADIC-IV-001\endcsname{\begin{itemize}
\item A real instrument matrix H and structural regressor matrix X with compatible finite row axes.
\item Both the instrument Gram matrix and projected-design Gram matrix are invertible.
\end{itemize}}
\expandafter\def\csname statadic@diagnostics@STATADIC-IV-001\endcsname{\begin{itemize}
\item Inspect SpatialIVFit.instrument\_rank, first\_stage\_rank, boundary and valid. Compare projected-regressor residual moments at the fitted coefficients. A full-rank first stage does not rule out weak instruments or establish instrument exogeneity; document both before interpreting an IV coefficient.
\end{itemize}}
\expandafter\def\csname statadic@statement@STATADIC-KKP-001\endcsname{Let $Q$ mix periods, $W_t$ act on units, and $L(u)_t=W_t u_t$.
  Then $Q L(u)=L(Qu)$ for every balanced panel $u$ if and only if
  $Q_{ts}\ne0$ implies $W_t=W_s$. Thus the full temporal mean commutes
  precisely for a common graph across periods.}
\expandafter\def\csname statadic@proof@STATADIC-KKP-001\endcsname{Expand both operators at period $t$ and unit $i$. Their terms
  are $Q_{ts}(W_s)_{ij}u_{sj}$ and $Q_{ts}(W_t)_{ij}u_{sj}$.
  The stated condition makes these terms equal. Conversely put a single unit
  impulse at $(s,j)$. Equality at $(t,i)$ forces equality of both matrix
  entries after cancelling the nonzero temporal coefficient.}
\expandafter\def\csname statadic@assumptions@STATADIC-KKP-001\endcsname{\begin{itemize}
\item Finite time and unit axes on a balanced panel.
\item Linear time mixing Q and a fixed real matrix of spatial weights for each period.
\end{itemize}}
\expandafter\def\csname statadic@diagnostics@STATADIC-KKP-001\endcsname{\begin{itemize}
\item KKP.estimate requires a balanced panel and one common graph. Verify the same unit ordering across periods, a repeated W, and the final fit’s valid flag. If periods mixed by the within/between projection have different W matrices, the commutation argument fails even when each W is row-normalized.
\end{itemize}}
\expandafter\def\csname statadic@statement@STATADIC-BOOTSTRAP-001\endcsname{Fix observations $x_i$ with normalized row masses $p_i$.
  Draw $m>0$ rows independently with replacement. Set
  $\mu=\sum p_ix_i$ and $v=\sum p_i(x_i-\mu)^2$. Then
  $E_*(\bar X^*-\mu)=0$ and
  $E_*[(\bar X^*-\mu)^2]=v/m$. Uniform masses give the empirical law.}
\expandafter\def\csname statadic@proof@STATADIC-BOOTSTRAP-001\endcsname{Each draw contributes $z_i=(x_i-\mu)/m$ to the centered mean.
  Its expectation is zero and its second moment is $v/m^2$.
  Conditional on an accumulated sum $s$, expand $(s+z_i)^2$.
  Averaging removes the mixed term. Induction over the independent draws
  gives $s^2+m v/m^2$ at $s=0$.}
\expandafter\def\csname statadic@assumptions@STATADIC-BOOTSTRAP-001\endcsname{\begin{itemize}
\item Fixed real observations with a normalized finite row law. Uniform row masses represent the ordinary empirical law.
\item A positive number of independent draws with replacement.
\end{itemize}}
\expandafter\def\csname statadic@diagnostics@STATADIC-BOOTSTRAP-001\endcsname{\begin{itemize}
\item For bootstrap\_estimate select the arithmetic mean, verify which data axis indexes sampling units, and document the iid-with-replacement row schedule. Its returned standard error uses finitely many replicates; compare it to the empirical variance divided by resample size for this specific statistic. A dependent time or cluster design needs another resampling law.
\end{itemize}}
\expandafter\def\csname statadic@statement@STATADIC-BOOTSTRAP-002\endcsname{For two independent fair Bernoulli observations, the
  expected ideal empirical bootstrap variance of their mean is $1/16$.
  The actual sampling variance of the mean is $1/8$.}
\expandafter\def\csname statadic@proof@STATADIC-BOOTSTRAP-002\endcsname{The four observed pairs have equal probabilities. Their
  conditional bootstrap variances are $0,1/8,1/8,0$ since two empirical
  draws are made. Their sample means are $0,1/2,1/2,1$, with mean $1/2$
  and variance $1/8$.}
\expandafter\def\csname statadic@assumptions@STATADIC-BOOTSTRAP-002\endcsname{\begin{itemize}
\item Two independent fair Bernoulli observations and the sample mean.
\item The ideal conditional bootstrap variance is computed over independent uniform resamples of size two.
\end{itemize}}
\expandafter\def\csname statadic@diagnostics@STATADIC-BOOTSTRAP-002\endcsname{\begin{itemize}
\item Separate variation over bootstrap resamples conditional on this dataset from variation over independently sampled datasets. More bootstrap replicates reduce simulation error but cannot remove the finite-sample bias shown by the two-row example. State the target when reporting a bootstrap standard error.
\end{itemize}}
\expandafter\def\csname statadic@statement@STATADIC-MULTIPLICITY-001\endcsname{Let $T$ be the true hypotheses among $n$ tests. Suppose
  $0\le\alpha\le1$ and $P(p_i\le u)\le u$ for every $i\in T$ and
  $0\le u\le1$. Reject each $i$ with $p_i\le\alpha/n$. The probability of
  rejecting any member of $T$ is at most $\alpha$, under arbitrary dependence.}
\expandafter\def\csname statadic@proof@STATADIC-MULTIPLICITY-001\endcsname{For nonempty $T$, also $n>0$. Apply the union bound over $T$,
  then marginal validity at $u=\alpha/n$. This bounds the rejection probability
  by $|T|\alpha/n\le\alpha$. The empty union has probability zero.}
\expandafter\def\csname statadic@assumptions@STATADIC-MULTIPLICITY-001\endcsname{\begin{itemize}
\item A finite family, a fixed subset of true nulls, and a common probability law on an arbitrary measurable outcome space.
\item For every true null and every cutoff u in [0,1], the probability of its p-value being at most u is at most u. The nominal level is in [0,1].
\end{itemize}}
\expandafter\def\csname statadic@diagnostics@STATADIC-MULTIPLICITY-001\endcsname{\begin{itemize}
\item Before calling equivalence\_multiplicity\_control, define the family, verify finite p-values in [0,1] and a valid nominal level, and justify marginal super-uniformity for each true null under the same joint law. The Bonferroni threshold does not repair an anticonservative input.
\end{itemize}}
\expandafter\def\csname statadic@statement@STATADIC-MULTIPLICITY-002\endcsname{Under the marginal validity and level assumptions above,
  rank the hypotheses for each outcome. Reject a hypothesis at zero-based
  rank $s$ only if every rank $j\le s$ passes its comparison to
  $\alpha/(n-j)$. The probability of rejecting any true null is at most
  $\alpha$. This includes Holm's ascending order with arbitrary ties.}
\expandafter\def\csname statadic@proof@STATADIC-MULTIPLICITY-002\endcsname{If any true null is rejected, let $r$ be the first rank of
  a true null. The prefix test must have passed at $r$. All $|T|$ true
  nulls occupy distinct ranks in the suffix from $r$ onward, hence
  $|T|\le n-r$. Thus the first true null has $p_i\le\alpha/(n-r)
  \le\alpha/|T|$. Take the union over true nulls, use marginal validity,
  and sum $|T|$ contributions of at most $\alpha/|T|$. If $T$ is empty
  there are no false rejections.}
\expandafter\def\csname statadic@assumptions@STATADIC-MULTIPLICITY-002\endcsname{\begin{itemize}
\item A finite family, a fixed subset of true nulls, and a common probability law on an arbitrary measurable outcome space.
\item For every true null and every cutoff u in [0,1], the probability of its p-value being at most u is at most u. The nominal level is in [0,1].
\end{itemize}}
\expandafter\def\csname statadic@diagnostics@STATADIC-MULTIPLICITY-002\endcsname{\begin{itemize}
\item Check the family definition and that all inputs are finite, marginally valid p-values under their respective nulls. Holm compares the sorted values sequentially and stops at the first failure. Reorder the decisions back to the original hypotheses; avoid an FWER interpretation for unvalidated bootstrap tails.
\end{itemize}}
\expandafter\def\csname statadic@statement@STATADIC-RANDOMIZATION-003\endcsname{Let a finite group $G$ act on a probability space of
  outcomes. Assume all transformed statistics are measurable and the joint
  distribution of $(T(hX))_{h\in G}$ agrees with that of $(T(hgX))_{h\in G}$
  for each $g$. Write $r(x)=\#\{h:T(hx)\ge T(x)\}$. For $\alpha\ge0$,
  the full-group test satisfies $P(r(X)\le\alpha |G|)\le\alpha$.}
\expandafter\def\csname statadic@proof@STATADIC-RANDOMIZATION-003\endcsname{Let $A$ be the inclusive rank rejection event. It is
  measurable as a threshold on a finite sum of comparisons. Joint orbit-law
  invariance gives $P(gX\in A)=P(X\in A)$ for each $g$. For every $x$,
  at most $\alpha |G|$ transformed observations reject, by the upper-rank
  count including ties. Integrate this bound, exchange the finite sum and
  integral, and divide by the positive cardinality $|G|$.}
\expandafter\def\csname statadic@assumptions@STATADIC-RANDOMIZATION-003\endcsname{\begin{itemize}
\item The finite group acts on an arbitrary measurable probability space. Each transformed real statistic is measurable.
\item For every group element g, the joint vector of statistics at hgX has the same distribution as the vector at hX over all h.
\item The level is nonnegative, the entire finite group is used, and the upper rank counts equality.
\end{itemize}}
\expandafter\def\csname statadic@diagnostics@STATADIC-RANDOMIZATION-003\endcsname{\begin{itemize}
\item For full enumeration, verify the group action and that the null makes the joint orbit-statistic vector invariant. For an add-one kernel, ensure the supplied nonidentity schedule enumerates the remaining group elements exactly once. Duplicates, omitted transforms, an incorrect alternative or nonfinite statistics prevent the exact rank correspondence.
\end{itemize}}
\expandafter\def\csname statadic@statement@STATADIC-SPATIAL-002\endcsname{Let $F$ be an invertible real innovation filter and let $P$
  select observations. If $F^T Z=P^T$, then
  $Z^T Z=P(F^{-1}F^{-T})P^T$. The identity holds for any rectangular $P$.}
\expandafter\def\csname statadic@proof@STATADIC-SPATIAL-002\endcsname{Transpose the solve to obtain $Z^T F=P$. Multiply on the right
  by $F^{-1}$, yielding $Z^T=PF^{-1}$. Substitute this in the Gram matrix
  and associate the products.}
\expandafter\def\csname statadic@assumptions@STATADIC-SPATIAL-002\endcsname{\begin{itemize}
\item The full-roster innovation filter F is an invertible finite real matrix.
\item The rectangular solve is exact: F\ensuremath{{}^T}Z = P\ensuremath{{}^T}. For the observed Gaussian interpretation P selects distinct observed rows and innovation covariance is \ensuremath{\sigma}\ensuremath{{}^2}I.
\end{itemize}}
\expandafter\def\csname statadic@diagnostics@STATADIC-SPATIAL-002\endcsname{\begin{itemize}
\item Align each observed row index with the full-period roster and graph. Check the full filters are nonsingular and the selected covariance supports Cholesky factorization before scoring. Dropping missing nodes from W changes paths through them. The solve cannot diagnose whether the mechanism selecting observed outcomes is ignorable.
\end{itemize}}
\expandafter\def\csname statadic@statement@STATADIC-SPATIAL-003\endcsname{For a fixed graph $W$ and multiplier $S$, the marginal effect
  matrix is $D=S(\beta I+\theta W)$. Its mean diagonal is
  $\beta\operatorname{tr}(S)/n+\theta\operatorname{tr}(SW)/n$. Its mean row
  sum is the same linear combination of the mean row sums of $S$ and $SW$.}
\expandafter\def\csname statadic@proof@STATADIC-SPATIAL-003\endcsname{Distribute $S$ over the sum and commute scalar multiplication
  through the matrix product. Apply linearity of the trace and of both finite
  sums, and divide by the roster size.}
\expandafter\def\csname statadic@assumptions@STATADIC-SPATIAL-003\endcsname{\begin{itemize}
\item Finite real matrices on the same nonempty roster and fixed coefficients.
\item For the spatial Durbin interpretation, W is held fixed and the spatial multiplier S equals (I \ensuremath{-} \ensuremath{\rho}W)\ensuremath{{}^{-1}} for a nonsingular filter.
\end{itemize}}
\expandafter\def\csname statadic@diagnostics@STATADIC-SPATIAL-003\endcsname{\begin{itemize}
\item For spatial\_impacts and spatial\_panel\_impacts, verify graph ordering, nonsingular lag filters and which covariates appear in the lagged design. Inspect valid separately from uncertainty\_valid. The derivative holds W fixed; if a covariate change would retrain or redefine W, compute or model that additional response separately.
\end{itemize}}
\expandafter\def\csname statadic@statement@STATADIC-ENERGY-004\endcsname{Suppose $P,Q$ are probability laws, $k(x,y)=\langle\phi(x),\phi(y)\rangle$
  for a real Hilbert-valued feature map, and $d(x,y)=k(x,x)+k(y,y)-2k(x,y)$.
  Assume $\phi$ and $\|\phi\|^2$ are integrable under each law. Writing $m_{AB}f$
  for the pair expectation, put $E=2m_{PQ}d-m_{PP}d-m_{QQ}d$.
  Then $E=2(m_{PP}k+m_{QQ}k-2m_{PQ}k)=2\|\mu_P-\mu_Q\|^2$,
  where $\mu_P$ and $\mu_Q$ are the Bochner feature means. Thus $E\ge0$ and
  $E=0$ exactly when the feature means agree.}
\expandafter\def\csname statadic@proof@STATADIC-ENERGY-004\endcsname{Integrate the generating relation for each pair of laws. Unit
  mass leaves one diagonal moment per draw. The inner-product integrals factor
  into the inner products of the two means. In the energy combination the
  diagonal moments cancel, leaving the expansion of the squared norm.}
\expandafter\def\csname statadic@assumptions@STATADIC-ENERGY-004\endcsname{\begin{itemize}
\item P and Q are probability measures on a common measurable space of arbitrary support.
\item k(x,y) is the inner product of supplied features \ensuremath{\phi}(x), \ensuremath{\phi}(y) in a complete real Hilbert space; d(x,y) = k(x,x) + k(y,y) \ensuremath{-} 2k(x,y).
\item The Hilbert feature and its squared norm are Bochner integrable under both laws.
\end{itemize}}
\expandafter\def\csname statadic@diagnostics@STATADIC-ENERGY-004\endcsname{\begin{itemize}
\item Identify the feature map inducing the selected ground, with integrable norm and squared norm for both population laws. The observed energy uses finite input masses; population integrability and injectivity need separate evidence before using zero distance to identify laws.
\end{itemize}}
\expandafter\def\csname statadic@statement@STATADIC-ENERGY-005\endcsname{Let $P,Q$ be probability measures and let a symmetric real
  array $k$ generate $d(x,y)=k(x,x)+k(y,y)-2k(x,y)$. Assume integrability of
  the diagonal under both laws and of $k$ under the three product measures
  $P\otimes P$, $P\otimes Q$ and $Q\otimes Q$. Writing $m_{AB}$ for the pair
  expectation, $2m_{PQ}d-m_{PP}d-m_{QQ}d$ equals
  $2(m_{PP}k-m_{PQ}k-m_{QP}k+m_{QQ}k)$.}
\expandafter\def\csname statadic@proof@STATADIC-ENERGY-005\endcsname{Integrate the generating relation with respect to each product
  measure. Both diagonal terms reduce to their marginal expectations. Symmetry
  and Fubini equate the two cross expectations. Subtracting within-law terms
  from twice the cross expectation cancels both diagonal moments.}
\expandafter\def\csname statadic@assumptions@STATADIC-ENERGY-005\endcsname{\begin{itemize}
\item P and Q are probability measures on the same measurable space, k is a symmetric real kernel, and d(x,y) = k(x,x) + k(y,y) \ensuremath{-} 2k(x,y).
\item The diagonal is integrable under both laws. The kernel is integrable under the P\ensuremath{\times}P, P\ensuremath{\times}Q and Q\ensuremath{\times}Q product measures.
\end{itemize}}
\expandafter\def\csname statadic@diagnostics@STATADIC-ENERGY-005\endcsname{\begin{itemize}
\item Check symmetry of the selected generating kernel and its diagonal relation to the ground. Record why product-kernel and diagonal terms are integrable under the two stated laws; inspecting the pooled Gram matrix can detect sample nonfiniteness but cannot establish the population tail conditions.
\end{itemize}}
\expandafter\def\csname statadic@statement@STATADIC-EQUIVALENCE-001\endcsname{Let $X$ be fair Bernoulli and $Y=0$, with unrooted energy
  and absolute distance. Their population energy is the null margin $\delta=1/2$.
  Observe two iid rows of $X$. For any nineteen two-row resamples of those rows,
  let $p=(1+\#\{E_b^*\ge\delta\})/20$. Under the null,
  $\Pr(p\le 1/20)\ge1/4$.}
\expandafter\def\csname statadic@proof@STATADIC-EQUIVALENCE-001\endcsname{The population cross mean is $1/2$ and the within-$X$
  distance mean is $1/2$, giving energy $1/2$. With probability $1/4$
  both observed rows are zero. On that event every empirical resample has
  zero energy, so the numerator is one for every row schedule. Counting this
  event alone gives the lower bound.}
\expandafter\def\csname statadic@assumptions@STATADIC-EQUIVALENCE-001\endcsname{\begin{itemize}
\item Absolute-distance unrooted energy between Bernoulli(1/2) and a point mass at zero. The null margin is 1/2.
\item Two iid rows of the first population and nineteen two-row empirical resamples, with any choice of indices.
\end{itemize}}
\expandafter\def\csname statadic@diagnostics@STATADIC-EQUIVALENCE-001\endcsname{\begin{itemize}
\item equivalence\_pvalue\_from\_bootstrap requires finite replicate energies and a finite margin; inspect the number of replicates at or above the margin. The add-one denominator controls resolution, not null calibration. Before reporting a level test or applying a familywise correction, supply a justified null resampling law or calibration argument for this margin.
\end{itemize}}
\expandafter\def\csname statadic@statement@STATADIC-LINEAR-012\endcsname{Write $G_n=X_n^T X_n/N$, $B_n(v)=X_n^T\operatorname{diag}(v^2)X_n/N$
  where $N=n+1$. Suppose $G_n\to A$, $B_n(u)\to V$ for true errors $u$,
  and $A$ is nonsingular. Suppose the maximum absolute difference between
  fitted residuals $r$ and $u$ tends to zero, while the sample means of
  $|x_i x_j||r+u|$ are bounded. Then $G_n^{-1}B_n(r)G_n^{-1}\to A^{-1}VA^{-1}$.}
\expandafter\def\csname statadic@proof@STATADIC-LINEAR-012\endcsname{Factor $r^2-u^2=(r-u)(r+u)$. For each meat entry, its
  replacement error is bounded by the residual difference bound times a
  bounded sample moment. Squeeze each entry to zero, add the oracle limit,
  and use continuity of inversion at $A$ and of matrix multiplication.}
\expandafter\def\csname statadic@assumptions@STATADIC-LINEAR-012\endcsname{\begin{itemize}
\item Finite coefficient dimension. On the chosen sample path the per-observation Gram G\_n converges to nonsingular A and the per-observation true-error score outer product converges to V.
\item The maximum absolute difference between fitted residuals r and true errors u is bounded by \ensuremath{\delta}\_n tending to zero. Bounded designs and coefficient consistency suffice for this difference condition.
\item The sample averages of |X\_i X\_j| |r + u| are uniformly bounded in n for every coefficient pair.
\end{itemize}}
\expandafter\def\csname statadic@diagnostics@STATADIC-LINEAR-012\endcsname{\begin{itemize}
\item Use hc\_covariance in its HC0 branch. Check full design rank, the smallest eigenvalue of the normalized Gram matrix, leverage and large score contributions across the intended sample sequence. Near singularity or concentrated scores warrant re-examining the assumed limits.
\item The proof assumes convergence of the ideal score meat and Gram matrix and bounds residual replacement. A single fitted covariance or finite condition number cannot establish those asymptotic assumptions.
\end{itemize}}
\expandafter\def\csname statadic@statement@STATADIC-MMD-007\endcsname{Let $k(x,y)=\langle\phi(x),\phi(y)\rangle$, where $\phi$ is
  Bochner integrable under every law in a class $C$. Then the within-plus-within
  minus twice cross kernel expectation $M(P,Q)$ is nonnegative on $C$.
  Furthermore, $M(P,Q)=0$ exactly for equal laws on $C$ if and only if
  $P\mapsto\int\phi\,dP$ is injective on $C$.}
\expandafter\def\csname statadic@proof@STATADIC-MMD-007\endcsname{Factor the three pair expectations into inner products of
  feature means and expand their squared norm difference. A squared norm is
  nonnegative and vanishes exactly at zero. Thus a zero discrepancy amounts
  to equality of the two embeddings. This is equivalent to equality of the
  laws throughout the class exactly when the embedding is injective there.}
\expandafter\def\csname statadic@assumptions@STATADIC-MMD-007\endcsname{\begin{itemize}
\item A symmetric positive-semidefinite kernel represented as k(x,y) = \ensuremath{\langle}\ensuremath{\phi}(x),\ensuremath{\phi}(y)\ensuremath{\rangle} in a complete real Hilbert space. The canonical RKHS feature map is one such representation.
\item A class C of probability measures under which \ensuremath{\phi} is Bochner integrable. No finite-support restriction or feature second-moment hypothesis.
\end{itemize}}
\expandafter\def\csname statadic@diagnostics@STATADIC-MMD-007\endcsname{\begin{itemize}
\item For weighted\_mmd\_squared\_cross, identify the feature map, the class of probability laws and integrability of the feature norm. To interpret zero as equality of laws, cite characteristicness on that class. A PSD finite Gram matrix or a small empirical V statistic cannot supply that premise.
\end{itemize}}
\expandafter\def\csname statadic@statement@STATADIC-NEIGHBOUR-001\endcsname{For one anchor, consider the same finite candidate roster before and
  after changing its distances from $d_i$ to $e_i$, with $|e_i-d_i|\le\epsilon$.
  If every member of $S$ is more than $2\epsilon$ closer than every nonmember,
  every perturbed top-$|S|$ set equals $S$. Ties within either side are allowed.}
\expandafter\def\csname statadic@proof@STATADIC-NEIGHBOUR-001\endcsname{If the equally sized sets differ, choose an original member $i$
  displaced by an outsider $j$. The gap gives
  $e_i\le d_i+\epsilon<d_j-\epsilon\le e_j$, whereas top-set membership
  requires $e_j\le e_i$.}
\expandafter\def\csname statadic@assumptions@STATADIC-NEIGHBOUR-001\endcsname{\begin{itemize}
\item The two distance rows use the same candidates and exclude the same anchor. Their top-k sets have equal cardinality.
\item Each distance changes by at most \ensuremath{\varepsilon}. Every original inside/outside pair has a strict gap greater than 2\ensuremath{\varepsilon}.
\end{itemize}}
\expandafter\def\csname statadic@diagnostics@STATADIC-NEIGHBOUR-001\endcsname{\begin{itemize}
\item neighbour\_order validates square, finite, symmetric distances and unique integer IDs; sorting breaks equal-distance ties by ID. Compare matrices on the same object ordering.
\item cross\_encoder\_neighbour\_stability reports mean\_top\_k\_overlap and exact\_nearest\_agreement. For a quantitative certificate, compute the maximum absolute row change and the boundary gap after self-exclusion; the two summaries do not provide these inputs.
\end{itemize}}
\expandafter\def\csname statadic@statement@STATADIC-GROUND-001\endcsname{On nonzero vectors put $u(x)=x/\|x\|$ and
  $d(x,y)=\|u(x)-u(y)\|$. Then $d(x,z)\le d(x,y)+d(y,z)$, and
  $d(x,y)=0$ exactly when $y=cx$ for some $c>0$.}
\expandafter\def\csname statadic@proof@STATADIC-GROUND-001\endcsname{Apply the ambient triangle inequality to the three directions.
  Zero distance gives equality of directions. Multiply by $\|y\|$ to obtain
  $y=(\|y\|/\|x\|)x$. Conversely, for $c>0$ the factor $c$ cancels
  against the same factor in $\|cx\|=c\|x\|$.}
\expandafter\def\csname statadic@assumptions@STATADIC-GROUND-001\endcsname{\begin{itemize}
\item Every input vector has positive finite norm on the declared raw domain.
\item Distances use the norm of the difference of normalized vectors.
\end{itemize}}
\expandafter\def\csname statadic@diagnostics@STATADIC-GROUND-001\endcsname{\begin{itemize}
\item Check for zero or numerically unrepresentable norms before selecting the ground. The low-level kernel divides by row norms and clips dot products; zero rows propagate NaN.
\item quotient\_metric requires a declared fibre witness. Use positive-ray equivalence when interpreting zero distance; distinct raw vectors can share a direction.
\end{itemize}}
\expandafter\def\csname statadic@statement@STATADIC-EMBEDDING-001\endcsname{Let $D$ be symmetric, hollow squared dissimilarities on a nonempty
  finite roster. Form $B_{ij}=-(D_{ij}-\bar D_i-\bar D_j+\bar D)/2$.
  If $B$ is positive semidefinite there is a real matrix $X$ with $B=XX^T$ and
  $D_{ij}=\sum_l (X_{il}-X_{jl})^2$ for every pair. Conversely, any such full
  factor of $B$ certifies positive semidefiniteness.}
\expandafter\def\csname statadic@proof@STATADIC-EMBEDDING-001\endcsname{Factor the positive semidefinite matrix as $XX^T$.
  The row means and grand mean cancel from $B_{ii}+B_{jj}-2B_{ij}$.
  Since $D_{ii}=D_{jj}=0$, the result is $D_{ij}$. Expand the three row
  products and collect squares coordinate by coordinate. A real Gram factor has
  nonnegative quadratic forms.}
\expandafter\def\csname statadic@assumptions@STATADIC-EMBEDDING-001\endcsname{\begin{itemize}
\item Finite, nonempty, common row/column roster and symmetric hollow squared dissimilarities.
\item Exact-real double centering with the same arithmetic mean on both axes.
\end{itemize}}
\expandafter\def\csname statadic@diagnostics@STATADIC-EMBEDDING-001\endcsname{\begin{itemize}
\item PCoA checks its input dissimilarities, representability of squared distances and its tolerance. pcoa\_kernel computes a symmetrized Gram, minimum\_eigenvalue, spectral\_scale, active and valid.
\item A negative eigenvalue below \ensuremath{-}t makes valid false, where t=tolerance \ensuremath{\times} spectral\_scale. A passing result can still have negative eigenvalues above that bound. Keep the spectrum and tolerance with the embedding.
\end{itemize}}
\expandafter\def\csname statadic@statement@STATADIC-EMBEDDING-002\endcsname{A symmetric hollow squared-dissimilarity matrix on a nonempty finite
  roster is realized by Euclidean coordinates exactly when its double-centered Gram is
  positive semidefinite. All axes are retained in this exact criterion.}
\expandafter\def\csname statadic@proof@STATADIC-EMBEDDING-002\endcsname{For the forward direction factor the PSD Gram and recover squared distances.
  For the converse, let $m_l$ be the mean of coordinate $l$. Expanding the squares in
  each row mean and the grand mean identifies the centered Gram with $YY^T$, where
  $Y_{il}=X_{il}-m_l$. Its quadratic form is nonnegative.}
\expandafter\def\csname statadic@assumptions@STATADIC-EMBEDDING-002\endcsname{\begin{itemize}
\item Finite, nonempty, common row/column roster and symmetric hollow squared dissimilarities.
\item Exact-real double centering with the same arithmetic mean on both axes.
\end{itemize}}
\expandafter\def\csname statadic@diagnostics@STATADIC-EMBEDDING-002\endcsname{\begin{itemize}
\item PCoA checks its input dissimilarities, representability of squared distances and its tolerance. pcoa\_kernel computes a symmetrized Gram, minimum\_eigenvalue, spectral\_scale, active and valid.
\item A negative eigenvalue below \ensuremath{-}t makes valid false, where t=tolerance \ensuremath{\times} spectral\_scale. A passing result can still have negative eigenvalues above that bound. Keep the spectrum and tolerance with the embedding.
\end{itemize}}
\expandafter\def\csname statadic@statement@STATADIC-EMBEDDING-003\endcsname{Let $U$ have orthonormal rows and write a reconstructed squared
  distance as $D_{ij}=\sum_l \lambda_l(U_{il}-U_{jl})^2$. Assume all eigenvalues
  are at least $-t$, where $t\ge0$. Retain only eigenvalues above $t$. Then for
  distinct rows the resulting squared distance differs from $D_{ij}$ by at most $2t$.}
\expandafter\def\csname statadic@proof@STATADIC-EMBEDDING-003\endcsname{On a discarded axis $|\lambda_l|\le t$. Thus the absolute
  omitted sum is at most $t\sum_l(U_{il}-U_{jl})^2$. Distinct orthonormal rows
  give this last sum as $1+1-2(0)=2$.}
\expandafter\def\csname statadic@assumptions@STATADIC-EMBEDDING-003\endcsname{\begin{itemize}
\item Two distinct rows of an exact orthonormal eigenvector matrix. The exact squared-distance expression uses their coordinate differences and signed eigenvalues.
\item Nonnegative cutoff t, and all eigenvalues at least \ensuremath{-}t. Only eigenvalues strictly above t are retained.
\end{itemize}}
\expandafter\def\csname statadic@diagnostics@STATADIC-EMBEDDING-003\endcsname{\begin{itemize}
\item PCoA checks its input dissimilarities, representability of squared distances and its tolerance. pcoa\_kernel computes a symmetrized Gram, minimum\_eigenvalue, spectral\_scale, active and valid.
\item A negative eigenvalue below \ensuremath{-}t makes valid false, where t=tolerance \ensuremath{\times} spectral\_scale. A passing result can still have negative eigenvalues above that bound. Keep the spectrum and tolerance with the embedding.
\item Use the same effective tolerance and spectral\_scale as the fit. Additional user-selected rank reduction requires its own discarded-eigenvalue bound.
\end{itemize}}
\expandafter\def\csname statadic@statement@STATADIC-MEDIAN-001\endcsname{Translate the proposed median to zero. Let the noncoincident residuals
  be $r_i\ne0$ with weights $w_i\ge0$, and let $M\ge0$ be the mass at zero.
  Write $R=\sum_i w_i r_i/\|r_i\|$ and $F(h)=M\|h\|+\sum_i w_i\|r_i-h\|$.
  Then $F(h)-F(0)\ge (M-\|R\|)\|h\|$ for every $h$.
  In particular $\|R\|\le M$ certifies global optimality.}
\expandafter\def\csname statadic@proof@STATADIC-MEDIAN-001\endcsname{Each unit vector $u_i=r_i/\|r_i\|$ supports its norm:
  $\|r_i-h\|\ge\langle u_i,r_i-h\rangle=\|r_i\|-\langle u_i,h\rangle$.
  Sum with nonnegative weights, add the coincident contribution $M\|h\|$,
  and use $\langle R,h\rangle\le\|R\|\|h\|$.}
\expandafter\def\csname statadic@assumptions@STATADIC-MEDIAN-001\endcsname{\begin{itemize}
\item Finite real Euclidean or inner-product sample, translated to a candidate at zero.
\item Separate observations at the candidate into their total mass M. All remaining residuals are nonzero and their weights are nonnegative.
\end{itemize}}
\expandafter\def\csname statadic@diagnostics@STATADIC-MEDIAN-001\endcsname{\begin{itemize}
\item Use spatial\_medians for its per-group convergence flag. It centers and rescales the group, classifies coincidence using a dtype-dependent radius, and accepts stationarity with a dtype- and group-size-dependent slack.
\item The theorem uses exact equality at a location and the exact force inequality. Record the group size, tolerance convention, iteration budget and convergence flag. weiszfeld\_medians uses softened distances and a fixed budget.
\end{itemize}}
\expandafter\def\csname statadic@statement@STATADIC-MEDIAN-002\endcsname{A nonempty finite sample in a proper metric space, including every
  finite-dimensional Euclidean group, has a point minimizing its sum of distances.}
\expandafter\def\csname statadic@proof@STATADIC-MEDIAN-002\endcsname{Choose an observed point $x_0$ and let $R$ be its objective value.
  Outside the closed ball about $x_0$ of radius $R$, the term $d(y,x_0)$ alone exceeds
  $R$. Inside, the continuous objective attains its minimum on the compact ball,
  which contains $x_0$.}
\expandafter\def\csname statadic@assumptions@STATADIC-MEDIAN-002\endcsname{\begin{itemize}
\item A nonempty finite unweighted sample and a proper metric space; Euclidean space of each finite dimension is proper.
\end{itemize}}
\expandafter\def\csname statadic@diagnostics@STATADIC-MEDIAN-002\endcsname{\begin{itemize}
\item Verify that each requested group has at least one valid coordinate row. The group-centering step divides by its membership count; an empty group has no certified target.
\end{itemize}}
\expandafter\def\csname statadic@statement@STATADIC-MEDIAN-003\endcsname{For residuals $r_i\ne0$, weights $w_i\ge0$, and coincident mass
  $M\ge0$, zero globally minimizes $F(h)=M\|h\|+\sum_i w_i\|r_i-h\|$
  exactly when $\|R\|\le M$, where $R=\sum_i w_i r_i/\|r_i\|$.}
\expandafter\def\csname statadic@proof@STATADIC-MEDIAN-003\endcsname{The supporting-plane bound proves sufficiency. For necessity,
  $(\|r-h\|-\|r\|)^2\ge0$ yields the upper bound
  $\|r-h\|-\|r\|\le-\langle r/\|r\|,h\rangle+\|h\|^2/(2\|r\|)$.
  Set $h=tR$. If $\|R\|>M$, the linear term in the total objective change
  is strictly negative, while the remainder is quadratic in $t$, giving descent.}
\expandafter\def\csname statadic@assumptions@STATADIC-MEDIAN-003\endcsname{\begin{itemize}
\item Finite real Euclidean or inner-product sample, translated to a candidate at zero.
\item Separate observations at the candidate into their total mass M. All remaining residuals are nonzero and their weights are nonnegative.
\item The coincident mass M is nonnegative.
\end{itemize}}
\expandafter\def\csname statadic@diagnostics@STATADIC-MEDIAN-003\endcsname{\begin{itemize}
\item Use spatial\_medians for its per-group convergence flag. It centers and rescales the group, classifies coincidence using a dtype-dependent radius, and accepts stationarity with a dtype- and group-size-dependent slack.
\item The theorem uses exact equality at a location and the exact force inequality. Record the group size, tolerance convention, iteration budget and convergence flag. weiszfeld\_medians uses softened distances and a fixed budget.
\end{itemize}}
\expandafter\def\csname statadic@statement@STATADIC-BARYCENTRE-001\endcsname{For a fixed support row and fixed matchings let $x_i$ be the
  matched source locations. With nonnegative weights summing to one, put
  $m=\sum_i w_i x_i$. At any $y$, the fixed-assignment squared cost satisfies
  $\sum_i w_i\|x_i-y\|^2=\sum_i w_i\|x_i-m\|^2+\|y-m\|^2$.
  Thus the centroid is the unique minimizing location for this row.}
\expandafter\def\csname statadic@proof@STATADIC-BARYCENTRE-001\endcsname{Write $x_i-y=(x_i-m)-(y-m)$ and expand the squared norms.
  The cross terms sum to zero since $\sum_i w_i(x_i-m)=0$.
  The remaining excess cost is strictly positive unless $y=m$.}
\expandafter\def\csname statadic@assumptions@STATADIC-BARYCENTRE-001\endcsname{\begin{itemize}
\item Squared Euclidean ground cost and normalized source weights. The simplex interpretation also requires nonnegative weights.
\item Hold the assignment from each source fixed. Each target support row has exactly one matched location from each source.
\end{itemize}}
\expandafter\def\csname statadic@diagnostics@STATADIC-BARYCENTRE-001\endcsname{\begin{itemize}
\item The eager solver checks same-cardinality uniform source laws and finite simplex weights. Check the actual metric is Euclidean and p is 2 before invoking the theorem.
\item Inspect converged, iterations, support\_shift, objective and initial\_objective. Stable assignments indicate a fixed point; optimality\_gap is improvement over the initialization and supplies no global lower bound.
\end{itemize}}
\expandafter\def\csname statadic@statement@STATADIC-TRANSPORT-001\endcsname{A nonnegative transport matrix $P$ has row masses $a_i$ and
  column masses $b_j$. If potentials satisfy $u_i+v_j\le C_{ij}$, then
  $\sum P_{ij}C_{ij}-\sum a_iu_i-\sum b_jv_j=
  \sum P_{ij}(C_{ij}-u_i-v_j)\ge0$.
  Equality on the support of $P$ certifies minimality among all feasible plans.}
\expandafter\def\csname statadic@proof@STATADIC-TRANSPORT-001\endcsname{Insert the row and column mass constraints into the potential sums.
  Collect terms by cell. Each reduced-cost contribution is nonnegative and vanishes
  when its mass is zero or its potential constraint is tight. Hence a tight plan
  attains a lower bound shared by every other feasible plan.}
\expandafter\def\csname statadic@assumptions@STATADIC-TRANSPORT-001\endcsname{\begin{itemize}
\item Finite nonnegative transport plan with the required row and column masses.
\item The supplied dual potentials sum to at most the unregularized cost in every cell.
\end{itemize}}
\expandafter\def\csname statadic@diagnostics@STATADIC-TRANSPORT-001\endcsname{\begin{itemize}
\item balanced\_wasserstein\_from\_cost checks a nonempty square, finite, nonnegative cost matrix and finite p\ensuremath{\ge}1. Its public result is a scalar.
\item An independent certificate needs the assignment or coupling and dual potentials. Check both marginals, all reduced costs, and the gap on the same powered cost matrix; the scalar alone supplies none of these witnesses.
\end{itemize}}
\expandafter\def\csname statadic@statement@STATADIC-SAMPLE-001\endcsname{Fix row weights $p_i$, points $x_i$, a linear map $A$ and
  an offset $b$. The first moment of the pushforward is
  $\sum_i p_i(Ax_i+b)=A(\sum_i p_i x_i)+(\sum_i p_i)b$.
  With unit mass the offset appears exactly once. Zero-mass rows have no contribution.}
\expandafter\def\csname statadic@proof@STATADIC-SAMPLE-001\endcsname{Distribute scalar multiplication across the translated point.
  Commute the linear map with each scalar multiplication and the finite sum.
  Collect the constant terms to obtain the total mass times the offset.}
\expandafter\def\csname statadic@assumptions@STATADIC-SAMPLE-001\endcsname{\begin{itemize}
\item Finite rows, unchanged weights and the same affine map A and offset b for every atom.
\item The probability interpretation requires nonnegative normalized mass.
\end{itemize}}
\expandafter\def\csname statadic@diagnostics@STATADIC-SAMPLE-001\endcsname{\begin{itemize}
\item Inspect PackedClouds.valid. It checks count bounds, positive dimension, finite positive-mass coordinates, nonnegative finite masses, zero padding masses and normalization within eight dtype epsilons.
\item For a sensitive offset, inspect the actual mass residual. The exact offset discrepancy is (total mass \ensuremath{-} 1)b; the mean and affine methods themselves do not turn this into an exact normalization check.
\end{itemize}}
\expandafter\def\csname statadic@statement@STATADIC-DISCO-001\endcsname{Write $p_g$ for each empirical group law, $w_g$ for its
  share, and $A_{gh}=\sum_{ij}p_{gi}p_{hj}d_{ij}$. Suppose the symmetric
  distance array has nonpositive quadratic form on zero-sum vectors. Then
  \[
  T-W=\tfrac{n}{4}\sum_{g,h}w_gw_h(2A_{gh}-A_{gg}-A_{hh})\ge0,
  \]
  with $T=(n/2)\sum_{gh}w_gw_h A_{gh}$ and $W=(n/2)\sum_gw_g A_{gg}$.}
\expandafter\def\csname statadic@proof@STATADIC-DISCO-001\endcsname{Expand the weighted pair sum. Summing out either group in the
  two diagonal terms uses $\sum_g w_g=1$ and recovers the within dispersion.
  Each remaining pair energy is minus the distance quadratic evaluated on
  $p_g-p_h$, whose mass is zero. Weight these nonnegative contributions by
  $n w_g w_h/4$.}
\expandafter\def\csname statadic@assumptions@STATADIC-DISCO-001\endcsname{\begin{itemize}
\item Finite normalized group shares and group laws on a common finite support. For empirical DISCO the group shares equal the observed group size divided by pooled size.
\item The dissimilarity is symmetric and its quadratic form is nonpositive on all zero-mass signed vectors on the support.
\end{itemize}}
\expandafter\def\csname statadic@diagnostics@STATADIC-DISCO-001\endcsname{\begin{itemize}
\item Construct codes with group\_design. It rejects empty or noninteger vectors, noncontiguous or inconsistent group IDs, fewer than two groups and designs with no residual degrees of freedom.
\item Document why the selected ground has negative type. Symmetry and nonnegative entries on a particular matrix do not certify it. Inspect positive total and within components before reporting the dispersion index or F ratio.
\end{itemize}}
\expandafter\def\csname statadic@statement@STATADIC-GEL-001\endcsname{For finite nonempty moment rows $g_i$ and positive
  carrier weights $a_i$, suppose $\sum_i a_i g_i=0$. Normalizing by
  $A=\sum_i a_i$ gives positive probabilities $p_i=a_i/A$ with
  $\sum_i p_i g_i=0$. Any projection $b^Tg_i$ that is nonnegative on
  every row must vanish on every row.}
\expandafter\def\csname statadic@proof@STATADIC-GEL-001\endcsname{The total carrier weight is strictly positive. Divide the
  normalization and stationarity sums by it. Sum the stationarity equations
  against $b$ and exchange the two finite sums. The sum of the nonnegative
  numbers $a_i b^Tg_i$ is zero, forcing each projection to vanish.}
\expandafter\def\csname statadic@assumptions@STATADIC-GEL-001\endcsname{\begin{itemize}
\item Nonempty finite rows. Unnormalized carrier weights are strictly positive and their weighted sum of every retained moment coordinate is exactly zero.
\item EL requires 1\ensuremath{-}\ensuremath{\lambda}\ensuremath{{}^T}g\ensuremath{{}_i}>0 for each row. ET has a positive exponential carrier for finite real multipliers.
\end{itemize}}
\expandafter\def\csname statadic@diagnostics@STATADIC-GEL-001\endcsname{\begin{itemize}
\item Check FixedGelFit.converged, criterion, effective\_rank and the rank cutoff. EL candidates must stay inside the log domain. The nonlinear solve checks the retained moment span and tests corrected probability positivity by backward error.
\item Inspect a separating linear combination in the retained coordinates: if every projected row is nonnegative and any is strictly positive, no strictly positive balancing distribution exists. Reassess overlap or the moment choice before interpreting the fit.
\end{itemize}}
\expandafter\def\csname statadic@statement@STATADIC-SURVIVAL-001\endcsname{At exit $t$ write $f,S$ for the failure density and
  survival and $g,G$ for the censor density and survival. At entry $a<t$
  let $S_a,G_a>0$. Conditional independence gives joint contributions
  $fG/(S_aG_a)$ for an event and $Sg/(S_aG_a)$ for censoring.
  With positive evaluated factors, their log splits as
  \[
  (\log(f\text{ or }S)-\log S_a)
  +(\log(G\text{ or }g)-\log G_a).
  \]
  Survival beyond entry has conditional probability $S/S_a\in[0,1]$
  if $S\le S_a$.}
\expandafter\def\csname statadic@proof@STATADIC-SURVIVAL-001\endcsname{Select the numerator according to the observed indicator.
  Split the logarithm of the product and quotient using positivity. The
  censoring term is constant across failure parameters when its model is
  separate. Dividing $0<S\le S_a$ by $S_a>0$ bounds the conditional tail.}
\expandafter\def\csname statadic@assumptions@STATADIC-SURVIVAL-001\endcsname{\begin{itemize}
\item Conditional on covariates, failure and censor times yield the indicated independent joint event/censor factors. Entry is noninformative under the specified observation scheme.
\item Entry precedes exit. All evaluated densities and survival factors in the theorem are positive and the failure survival cannot increase from entry to exit. The censoring model is free of failure-model parameters if dropped from fitting.
\end{itemize}}
\expandafter\def\csname statadic@diagnostics@STATADIC-SURVIVAL-001\endcsname{\begin{itemize}
\item Check SurvivalFit.input\_valid, identified and valid, along with solver.converged. The fit tests finite design and times, positive duration, nonnegative entry strictly before duration, and binary event codes. No-event or degenerate exact fits are guarded.
\item Document how subjects entered the cohort, why censoring is conditionally noninformative, and which covariates support that assumption. Valid row encodings and local curvature do not establish the observation model.
\end{itemize}}
\expandafter\def\csname statadic@statement@STATADIC-SURVIVAL-002\endcsname{Suppose $t>0$, $\sigma>0$, $z=(\log t-\eta)/\sigma$,
  and the standardized log-time law has survival $S_0$ with derivative
  $-p$ at $z$, where $p>0$. The AFT survival has derivative
  $-p/(t\sigma)$ at $t$, so its log event density on original duration units
  is $\log p-\log t-\log\sigma$.}
\expandafter\def\csname statadic@proof@STATADIC-SURVIVAL-002\endcsname{Differentiate the log-time standardization to obtain
  $1/(t\sigma)$. Compose with the baseline survival derivative, then use
  positivity to split the logarithm of the density into three terms.}
\expandafter\def\csname statadic@assumptions@STATADIC-SURVIVAL-002\endcsname{\begin{itemize}
\item Positive observed duration t and log-time scale \ensuremath{\sigma}. The standardized survivor is differentiable with derivative \ensuremath{-}p at (log t\ensuremath{-}\ensuremath{\eta})/\ensuremath{\sigma} and p is positive there.
\end{itemize}}
\expandafter\def\csname statadic@diagnostics@STATADIC-SURVIVAL-002\endcsname{\begin{itemize}
\item Record duration units and the ancillary convention before comparing fits. Lognormal uses log-scale ancillary; Weibull and log-logistic use log shape. Check positivity of the implied scale and duration and inspect nonfinite event contributions in tail data.
\end{itemize}}
\expandafter\def\csname statadic@statement@STATADIC-DYADIC-001\endcsname{In an undirected dyad $(i,j)$ the node-effect design
  has entries $D_{ijk}=1_{i=k}+1_{j=k}$. With any base node $b$,
  \[
    c+u_i+u_j=(c+2u_b)+\sum_{k\ne b}D_{ijk}(u_k-u_b).
  \]
  Consequently changing the base preserves the fitted linear predictor.
  Including every effect and the intercept leaves a common-shift ambiguity.}
\expandafter\def\csname statadic@proof@STATADIC-DYADIC-001\endcsname{Insert the omitted base term, whose contrast is zero.
  The endpoint indicator sums select $u_i-u_b$ and $u_j-u_b$.
  Adding the adjusted intercept recovers the original predictor.}
\expandafter\def\csname statadic@assumptions@STATADIC-DYADIC-001\endcsname{\begin{itemize}
\item A finite common node index, two endpoints per row, an intercept and additive effects from the same coefficient vector at both endpoints.
\end{itemize}}
\expandafter\def\csname statadic@diagnostics@STATADIC-DYADIC-001\endcsname{\begin{itemize}
\item Use the validated canonical node index and record its first node as the reference. The dyadic sample rejects self-pairs. Inspect the retained and dropped column names from build\_dyadic\_design, including dropped regressors, before interpreting any coefficient.
\item A shared endpoint creates a potential dependence channel. Record the sampling unit before using standard errors or a node bootstrap; successful rank cleaning only addresses linear redundancy.
\end{itemize}}
\expandafter\def\csname statadic@statement@STATADIC-SHARPE-001\endcsname{Let $n\ge2$, $\bar r=n^{-1}\sum_i r_i$ and
  $s^2=(n-1)^{-1}\sum_i(r_i-\bar r)^2$. Then $s=0$ exactly when
  every observation equals $\bar r$. For $c>0$,
  $s(cr)=c\,s(r)$ and the ratio $\overline{cr}/s(cr)=\bar r/s(r)$
  whenever the series has positive volatility.}
\expandafter\def\csname statadic@proof@STATADIC-SHARPE-001\endcsname{The divisor $n-1$ is positive. A sum of squares vanishes
  precisely when every deviation vanishes. Rescaling multiplies the
  deviations by $c$, hence the variance by $c^2$ and the standard deviation
  by $c$. Cancel the common factor in the ratio.}
\expandafter\def\csname statadic@assumptions@STATADIC-SHARPE-001\endcsname{\begin{itemize}
\item A fixed series of at least two finite real differential returns, in common units and at the same observation frequency. Positive rescaling.
\end{itemize}}
\expandafter\def\csname statadic@diagnostics@STATADIC-SHARPE-001\endcsname{\begin{itemize}
\item Form excess or differential returns against the intended benchmark before calling sharpe\_ratio. Check finite observations, sample size, peak-to-peak variation and positive sample standard deviation. For constant or insufficient data the API returns a numeric zero; report the degeneracy rather than interpreting a zero risk-adjusted performance estimate.
\end{itemize}}
\expandafter\def\csname statadic@statement@STATADIC-INEQUALITY-001\endcsname{For standardized moments $z_j$, the squared negative-part
  statistic $M(z)=\sum_j\min(z_j,0)^2$ is the minimum of
  $\sum_j(z_j-t_j)^2$ over $t_j\ge0$, attained at $t_j=\max(z_j,0)$.
  It vanishes exactly when every standardized sample moment is nonnegative.}
\expandafter\def\csname statadic@proof@STATADIC-INEQUALITY-001\endcsname{When $z_j\ge0$, choosing $t_j=z_j$ gives zero.
  When $z_j<0$, the closest admissible value is zero, since moving right
  increases its distance to $z_j$. Sum these coordinatewise bounds.
  Vanishing of the sum of squares forces each negative part to vanish.}
\expandafter\def\csname statadic@assumptions@STATADIC-INEQUALITY-001\endcsname{\begin{itemize}
\item A finite vector of real standardized sample moments. Positive coordinate scales are needed to interpret each entry as \ensuremath{\sqrt{n}} times a sample mean divided by its estimated scale. The null moment inequalities have the nonnegative sign convention.
\end{itemize}}
\expandafter\def\csname statadic@diagnostics@STATADIC-INEQUALITY-001\endcsname{\begin{itemize}
\item Check the moment sign convention, selected statistic, coordinate scales and full-rank HAC covariance. The reported n\_selected and kappa\_n describe the selection step; eta and calibration identify the correction actually used.
\item The default calibration is explicitly uncorrected indicator GMS. Record the simulation count, dependence assumptions and bandwidth when using its p-value; only the observed MMM projection is justified by this claim.
\end{itemize}}
\expandafter\def\csname statadic@statement@STATADIC-SHRINKAGE-001\endcsname{For $p>0$, let $F_{ii}=\sigma_i^2$ and
  $F_{ij}=r\sigma_i\sigma_j$ for $i\ne j$. Suppose
  $1-r\ge0$ and $1+(p-1)r\ge0$. For a nonnegative sample quadratic
  form $S$ and $0\le a\le1$, every contrast of $aF+(1-a)S$
  has nonnegative variance.}
\expandafter\def\csname statadic@proof@STATADIC-SHRINKAGE-001\endcsname{Set $z_i=\sigma_i x_i$, $u=\sum z_i^2$ and
  $v=(\sum z_i)^2/p$. The target risk expands to
  $(1-r)(u-v)+(1+(p-1)r)v$. Cauchy--Schwarz gives $0\le v\le u$.
  Both terms and the convex combination with sample risk are nonnegative.}
\expandafter\def\csname statadic@assumptions@STATADIC-SHRINKAGE-001\endcsname{\begin{itemize}
\item At least one asset and a real contrast with nonnegative sample risk.
\item The blend intensity is between zero and one. The common correlation r satisfies 1\ensuremath{-}r \ensuremath{\ge} 0 and 1+(p\ensuremath{-}1)r \ensuremath{\ge} 0.
\end{itemize}}
\expandafter\def\csname statadic@diagnostics@STATADIC-SHRINKAGE-001\endcsname{\begin{itemize}
\item Estimand(Covariance(...), estimator=LedoitWolf()).prepare delegates to the covariance owner's preparation and retains PreparedCovariance. It validates finite rectangular observations and rejects zero marginal sample variances in a multi-asset fit. Native LedoitWolf.estimate and the bare kernel expect prepared inputs. The one-asset branch returns its sample covariance directly.
\item Read CovarianceEstimate.target and shrinkage. For p>1, divide an off-diagonal target entry by its two standard deviations to recover r and check both scalar bounds. Nonnegative portfolio variance does not test out-of-sample risk dominance.
\end{itemize}}
\expandafter\def\csname statadic@statement@STATADIC-GLMCOV-001\endcsname{Let $U$ be the matrix of fitted score rows, $B$ the
  coefficient response map, and $c>0$ the common scale, including the sample
  normalization. For $V=c BU^TUB^T$, a contrast $v$ has
  $v^TVv=c\|UB^Tv\|^2$. In particular, its variance vanishes exactly when
  every fitted score row annihilates $B^Tv$.}
\expandafter\def\csname statadic@proof@STATADIC-GLMCOV-001\endcsname{Move both transposes across the scalar product to write the
  quadratic form as a sum of squared projected scores. A positive common
  multiplier preserves its zeros. A sum of real squares vanishes precisely
  when every term vanishes.}
\expandafter\def\csname statadic@assumptions@STATADIC-GLMCOV-001\endcsname{\begin{itemize}
\item U contains fitted score rows in the same parameter order as the coefficient-response map B.
\item The matrix scale c is positive. At independent-row mean-equation scaling, c includes the common 1/n\ensuremath{{}^2} and finite-sample multiplier.
\end{itemize}}
\expandafter\def\csname statadic@diagnostics@STATADIC-GLMCOV-001\endcsname{\begin{itemize}
\item Inspect GLMUncertainty.valid and status, fit.identified and fit.df\_resid. The producer gates on fit validity, numerical covariance validity and positive diagonal entries. HC1 additionally needs positive residual degrees of freedom.
\item For a proposed contrast, inspect the fitted score projections along B\ensuremath{{}^T}v or their squared norm. To interpret the HC matrix as sampling covariance, justify the estimating equation, contribution-row independence or appropriate uncorrelatedness, and asymptotic regularity separately.
\end{itemize}}
\expandafter\def\csname statadic@statement@STATADIC-PANELCOV-001\endcsname{For influence rows $U$ with unique entity--time keys,
  form the entity and time membership matrices $E,T$. The unadjusted
  two-way covariance $(EU)^T(EU)+(TU)^T(TU)-U^TU$ equals $U^TKU$,
  where $K_{ij}$ is one when rows share an entity or a time, and zero otherwise.}
\expandafter\def\csname statadic@proof@STATADIC-PANELCOV-001\endcsname{The two membership Gram matrices pick the pairs sharing each
  index. Uniqueness makes their intersection the diagonal. Inclusion--exclusion
  therefore gives $E^TE+T^TT-I=K$. Multiply on each side by the influences.}
\expandafter\def\csname statadic@assumptions@STATADIC-PANELCOV-001\endcsname{\begin{itemize}
\item Finite coefficient-influence rows and entity/time labels, with injective joint labels.
\item Unadjusted cluster totals; set group\_debias=False. A common final scalar correction may be applied afterward.
\end{itemize}}
\expandafter\def\csname statadic@diagnostics@STATADIC-PANELCOV-001\endcsname{\begin{itemize}
\item Prepared panel indices reject duplicate entity/time pairs. Check method, group\_debias and the occupied entity/time counts. The returned n\_clusters is the smaller count for two-way clustering.
\item The validity gate checks fit validity, correction denominator, minimum occupied count and numerical covariance evidence. Few or highly unbalanced clusters and dependence outside both labels require a separate judgment.
\end{itemize}}
\expandafter\def\csname statadic@statement@STATADIC-PANELCOV-002\endcsname{Let $z_i$ be centered, scalar influence contributions under
  a finite law. Suppose $E[z_i z_j]=0$ whenever rows differ in both entity and
  time. Then $\operatorname{Var}(\sum_i z_i)=
  E[\sum_{i,j}1\{e_i=e_j\ \mathrm{or}\ t_i=t_j\}\,z_i z_j]$. Shared pairs
  may have arbitrary cross moments.}
\expandafter\def\csname statadic@proof@STATADIC-PANELCOV-002\endcsname{The total has mean zero. Expand its second moment as the
  sum of all pair expectations. The dependence assumption removes every
  pair outside the cluster union. Exchange the finite sums and expectation
  to obtain the claimed masked expression.}
\expandafter\def\csname statadic@assumptions@STATADIC-PANELCOV-002\endcsname{\begin{itemize}
\item Finite probability law and centered scalar row influences.
\item Each pair sharing neither entity nor time has zero raw cross moment. Shared pairs may have arbitrary cross moments.
\end{itemize}}
\expandafter\def\csname statadic@diagnostics@STATADIC-PANELCOV-002\endcsname{\begin{itemize}
\item Document why an observation outside both labels cannot share a shock with the focal observation. Common omitted factors or geographic links across entities and periods can violate this premise.
\item PanelUncertainty.valid and n\_clusters reflect fitting and numerical/occupancy checks. They do not test the cross-moment restriction.
\end{itemize}}
\expandafter\def\csname statadic@statement@STATADIC-HAC-001\endcsname{Let $U$ contain $n>0$ ordered score rows and let
  $b\ge0$ be an integer. Define $K_{ts}=(b+1-|t-s|)_+/(b+1)$ on the
  observed time positions. Then the Bartlett covariance $U^TKU/n$ is
  positive semidefinite. Scores outside the observation range contribute zero.}
\expandafter\def\csname statadic@proof@STATADIC-HAC-001\endcsname{Introduce the matrix $A_{kt}=1$ for $t\le k<t+b+1$
  and zero otherwise. Counting the integer overlap of two windows gives
  $(A^TA)_{ts}=(b+1-|t-s|)_+$. Thus the covariance is the Gram matrix
  of $AU$ divided by $n(b+1)>0$.}
\expandafter\def\csname statadic@assumptions@STATADIC-HAC-001\endcsname{\begin{itemize}
\item Nonempty equally spaced time rows and a nonnegative integer bandwidth b.
\item Real contribution rows in the chosen order and the same sample-size divisor for all lags.
\end{itemize}}
\expandafter\def\csname statadic@diagnostics@STATADIC-HAC-001\endcsname{\begin{itemize}
\item The operator returns a matrix. Check that adjacent rows represent successive equally spaced time units and that the chosen bandwidth is a nonnegative integer. STATADIC uses row positions even across calendar gaps.
\item Record the bandwidth, centering convention and sensitivity to nearby bandwidths. An eigenvalue check addresses numerical PSD; weak dependence, moment conditions and bandwidth growth for consistency require separate justification.
\end{itemize}}
\expandafter\def\csname statadic@statement@STATADIC-HAC-002\endcsname{Let a real time series under a finite law have the same
  raw second moment $\gamma_h=E[z_t z_{t+h}]$ at every observed pair of lag
  $h$. For $n>0$ the expectation of the uncentered lag sum with denominator
  $n$ is $((n-h)_+/n)\gamma_h$. This equals a lag covariance interpretation
  when the process has mean zero.}
\expandafter\def\csname statadic@proof@STATADIC-HAC-002\endcsname{Move expectation inside the sum. Exactly $(n-h)_+$ indices
  have both endpoints in the record. Replace each pair expectation by the
  common lag moment and divide by $n$.}
\expandafter\def\csname statadic@assumptions@STATADIC-HAC-002\endcsname{\begin{itemize}
\item Finite-law real contribution process and a common raw lag-h pair expectation at all observed pairs.
\item For scientific interpretation n>0; the Lean formula uses total real division and truncated natural subtraction.
\end{itemize}}
\expandafter\def\csname statadic@diagnostics@STATADIC-HAC-002\endcsname{\begin{itemize}
\item Compare the bandwidth with sample length and record how many pairs contribute at the largest included lag.
\item The raw constructor does not center or test score means. The normalized row count and a PSD matrix cannot establish second-order stationarity or a dependence rate.
\end{itemize}}
\expandafter\def\csname statadic@statement@STATADIC-BLOCKLENGTH-001\endcsname{Let $n,D,b_0>0$ and $D b_0^3=2nG^2$. Then
  $b_0$ uniquely minimizes $R(b)=G^2/b^2+Db/n$ over $b>0$.}
\expandafter\def\csname statadic@proof@STATADIC-BLOCKLENGTH-001\endcsname{Substitute the balance condition and factor the risk gap as
  $R(b)-R(b_0)=(D/n)(b-b_0)^2(2b+b_0)/(2b^2)$.
  The factor multiplying the square is strictly positive on the stated domain.}
\expandafter\def\csname statadic@assumptions@STATADIC-BLOCKLENGTH-001\endcsname{\begin{itemize}
\item A positive observation count, positive variance constant D and positive candidate b\ensuremath{{}_0} with D b\ensuremath{{}_0}\ensuremath{{}^3} = 2nG\ensuremath{{}^2}.
\item Continuous positive competing block lengths. The balance hypothesis excludes zero G.
\end{itemize}}
\expandafter\def\csname statadic@diagnostics@STATADIC-BLOCKLENGTH-001\endcsname{\begin{itemize}
\item The API checks the scheme name and that at least two observations exist. Record stationary versus circular resampling, max\_lag and the resulting block length. The numeric output alone does not report an interior optimum.
\item Review \ensuremath{\hat G}, D before its 10\ensuremath{{}^{-8}} floor, the length before its floor at one, and the flat-top lag sensitivity if using the interior result. Constant or nearly degenerate series can activate the floors. The selector does not test stationarity, weak dependence or validity for a degenerate functional.
\end{itemize}}
\expandafter\def\csname statadic@statement@STATADIC-SIMPLEX-001\endcsname{Let $f(x)=x^TQx/2+c^Tx$ with $Q$ symmetric and positive semidefinite on zero-sum directions.
  Let $x$ have sum one and lie in $[0,u]^n$. Suppose lower and upper multipliers
  $l,r\ge0$ satisfy $l_ix_i=r_i(u-x_i)=0$ and
  $(Qx+c)_i-l_i+r_i=\nu$ for every coordinate. For every feasible $y$,
  $f(x)\le f(y)$.}
\expandafter\def\csname statadic@proof@STATADIC-SIMPLEX-001\endcsname{Write $h=y-x$. Its coordinates sum to zero, so the common multiplier
  drops out of $\nabla f(x)^Th$. Complementarity turns the remaining terms into
  $\sum_i l_iy_i+r_i(u-y_i)\ge0$. Adding $h^TQh/2\ge0$ gives the objective gap.}
\expandafter\def\csname statadic@assumptions@STATADIC-SIMPLEX-001\endcsname{\begin{itemize}
\item Finite real quadratic with a symmetric Hessian that is positive semidefinite on zero-sum directions.
\item The candidate sums to one, respects its box, has nonnegative inequality multipliers, and satisfies exact tangent stationarity and complementary slackness.
\end{itemize}}
\expandafter\def\csname statadic@diagnostics@STATADIC-SIMPLEX-001\endcsname{\begin{itemize}
\item Inspect convex, primal\_residual, dual\_residual and complementarity\_residual on SimplexQPSolution at the actual solver.value. A negative eigenvalue off the equality tangent need not invalidate this certificate.
\item Report the maximum weight and tolerance. If tangent curvature or feasibility fails, the global convex certificate is unavailable.
\end{itemize}}
\expandafter\def\csname statadic@statement@STATADIC-SIMPLEX-002\endcsname{Suppose $x$ is feasible for a quadratic with $Q$ symmetric and positive
  semidefinite on zero-sum directions. Let $L$ be a lower bound on $g^Ty$ over its feasible
  set, where $g=Qx+c$. For every feasible $y$, $f(x)-f(y)\le g^Tx-L$.
  This applies to the capped-simplex linear minimization gap at the returned iterate.}
\expandafter\def\csname statadic@proof@STATADIC-SIMPLEX-002\endcsname{Expand $f(y)$ about $x$ and discard the nonnegative quadratic
  remainder. Substitute $g^Ty\ge L$ into the resulting linear bound.}
\expandafter\def\csname statadic@assumptions@STATADIC-SIMPLEX-002\endcsname{\begin{itemize}
\item A symmetric quadratic with nonnegative curvature on zero-sum directions.
\item Candidate and competitor have the same total weight. L is a lower bound on the competitor’s linearized objective.
\end{itemize}}
\expandafter\def\csname statadic@diagnostics@STATADIC-SIMPLEX-002\endcsname{\begin{itemize}
\item Read optimality\_gap together with primal feasibility and convex. The quantity bounds objective suboptimality for feasible points in exact arithmetic. A small gap alone does not bound weight error in a flat direction.
\end{itemize}}
\expandafter\def\csname statadic@statement@STATADIC-SIMPLEX-003\endcsname{For $n>0$, there exists a nonnegative vector of total mass
  one with each entry at most $u$ if and only if $nu\ge1$.}
\expandafter\def\csname statadic@proof@STATADIC-SIMPLEX-003\endcsname{Sum the $n$ upper bounds to obtain necessity.
  Conversely $1/n\le u$ makes the uniform vector feasible.}
\expandafter\def\csname statadic@assumptions@STATADIC-SIMPLEX-003\endcsname{\begin{itemize}
\item There are n>0 coordinates, a common upper bound u, nonnegative weights and total mass one.
\end{itemize}}
\expandafter\def\csname statadic@diagnostics@STATADIC-SIMPLEX-003\endcsname{\begin{itemize}
\item Check maximum\_weight against 1/K before calling simplex\_qp\_result. Below this bound the cap makes the problem infeasible regardless of iteration budget.
\end{itemize}}
\expandafter\def\csname statadic@statement@STATADIC-PSD-001\endcsname{In an orthonormal eigenbasis $U$, replace $d_i$ by
  $c_i=\max(\delta,d_i)$. The change of trace is exactly
  $\sum_i\max(\delta-d_i,0)$. In particular it is zero precisely when every
  original eigenvalue meets the floor.}
\expandafter\def\csname statadic@proof@STATADIC-PSD-001\endcsname{Move $U^TU=I$ across the trace. Subtract the two diagonal
  sums term by term. Every summand is a nonnegative shortfall, so a zero sum
  forces every shortfall to vanish.}
\expandafter\def\csname statadic@assumptions@STATADIC-PSD-001\endcsname{\begin{itemize}
\item A finite real symmetric input expressed in an orthonormal eigenbasis.
\item Each eigenvalue below the chosen floor is raised to that floor.
\end{itemize}}
\expandafter\def\csname statadic@diagnostics@STATADIC-PSD-001\endcsname{\begin{itemize}
\item Compare total\_var\_after - total\_var\_before with the positive shortfalls computed from the symmetrized input using the same eps. Check offdiag\_distortion separately when off-diagonal structure matters.
\end{itemize}}
\expandafter\def\csname statadic@statement@STATADIC-PSD-002\endcsname{In an eigenbasis of the input, let $D=\mathrm{diag}(d)$
  and $C=\mathrm{diag}(\max(d_i,\delta))$. For every real matrix $B$ with
  $B_{ii}\ge\delta$, $\|B-D\|_F^2\ge\|C-D\|_F^2+\|B-C\|_F^2$.
  Every symmetric competitor above the spectral floor has these diagonal bounds.}
\expandafter\def\csname statadic@proof@STATADIC-PSD-002\endcsname{Off-diagonal terms agree on the two sides since $D$ and $C$
  are diagonal. When $d_i$ already meets the floor its correction is zero.
  Otherwise the cross term $2(B_{ii}-\delta)(\delta-d_i)$ is nonnegative.
  Sum the entrywise inequalities.}
\expandafter\def\csname statadic@assumptions@STATADIC-PSD-002\endcsname{\begin{itemize}
\item Work in a real eigenbasis of the input. Every competing diagonal entry meets the floor.
\end{itemize}}
\expandafter\def\csname statadic@diagnostics@STATADIC-PSD-002\endcsname{\begin{itemize}
\item Compare Frobenius errors using the same symmetric target and absolute eigenvalue floor. If preserving input variances is required, eig\_clip can change them.
\end{itemize}}
\expandafter\def\csname statadic@statement@STATADIC-PSD-003\endcsname{Suppose $A=U\operatorname{diag}(d)U^T$ with orthogonal $U$.
  Replacing $d_i$ by $c_i=\max(d_i,\delta)$ gives a symmetric matrix $C$
  with $z^TCz\ge\delta z^Tz$ for every $z$. Every competing symmetric $B$
  obeying that same floor satisfies $\|C-A\|_F^2\le\|B-A\|_F^2$.}
\expandafter\def\csname statadic@proof@STATADIC-PSD-003\endcsname{Put $v=U^Tz$. Orthogonality preserves length, and the
  clipped quadratic form is $\sum_i c_i v_i^2\ge\delta\sum_i v_i^2$.
  For the competitor, test its floor on each unit column of $U$ to bound the
  diagonal of $U^TBU$. Apply the entrywise floor projection inequality in that
  basis. Cyclic invariance of the trace preserves squared Frobenius error
  when returning to the original coordinates.}
\expandafter\def\csname statadic@assumptions@STATADIC-PSD-003\endcsname{\begin{itemize}
\item An exact orthogonal eigenbasis and real eigenvalues for the symmetrized target.
\item The competitor bounds every directional quadratic form below by the same floor times squared norm.
\end{itemize}}
\expandafter\def\csname statadic@diagnostics@STATADIC-PSD-003\endcsname{\begin{itemize}
\item Use the same eps on the input and competitor, and compare min\_eig\_after with it on the working scale. repair\_diagnostics.is\_psd\_after uses its own tolerance and can pass a matrix below a requested positive floor.
\end{itemize}}
\expandafter\def\csname statadic@statement@STATADIC-LSTSQ-001\endcsname{For $r=y-Xb$, suppose $X^Tr=0$. For every alternative
  $b+h$, the squared residual increases by exactly $\|Xh\|^2$.
  The coefficient is unique precisely when $Xh=0$ implies $h=0$.}
\expandafter\def\csname statadic@proof@STATADIC-LSTSQ-001\endcsname{Expand the squared norm of $r-Xh$. Its mixed term vanishes
  because $r^TXh=h^TX^Tr=0$. A competing point attains equality precisely
  when its fitted-vector increment has zero squared norm.}
\expandafter\def\csname statadic@assumptions@STATADIC-LSTSQ-001\endcsname{\begin{itemize}
\item A finite real design and a coefficient whose residual is orthogonal to its columns.
\end{itemize}}
\expandafter\def\csname statadic@diagnostics@STATADIC-LSTSQ-001\endcsname{\begin{itemize}
\item Use least\_squares\_result.rank and compare with the number of columns. If rank is deficient, fitted values can agree while coefficients differ. Independently evaluate X.T @ (y - X @ b) at the reported value.
\end{itemize}}
\expandafter\def\csname statadic@statement@STATADIC-FWL-001\endcsname{Let $Z$ be a nuisance design with invertible Gram matrix,
  $z=x-Z\hat\gamma$ and $t=y-Z\hat\delta$ its least-squares residuals.
  If $e=y-Za-bx$ is orthogonal to $Z$ and $x$, and $z^Tz>0$, then
  $b=z^Tt/(z^Tz)$.}
\expandafter\def\csname statadic@proof@STATADIC-FWL-001\endcsname{The first residualization gives $Z^Tz=0$, hence
  $z^Tt=z^Ty$ and $z^Tz=z^Tx$. Orthogonality of $e$ to $x$ and to
  every nuisance column gives $z^Te=0$. Substitute $y=Za+bx+e$ and divide.}
\expandafter\def\csname statadic@assumptions@STATADIC-FWL-001\endcsname{\begin{itemize}
\item The same nuisance regressors and aligned observations are used for each residualization. The nuisance Gram matrix is invertible.
\item The full-fit residual is orthogonal to nuisance and focal columns, and the focal residual has positive squared norm.
\end{itemize}}
\expandafter\def\csname statadic@diagnostics@STATADIC-FWL-001\endcsname{\begin{itemize}
\item Inspect the focal residual sum of squares relative to the original column scale, and check the nuisance design rank independently. STATADIC returns NaN for a nonpositive denominator; a tiny positive denominator still permits an unstable estimate.
\end{itemize}}
\expandafter\def\csname statadic@statement@STATADIC-INTEGRATION-001\endcsname{If $n$ simulation weights are nonnegative and
  $S_2=\sum w_i^2>0$, then $1\le (\sum w_i)^2/S_2\le n$.
  Equality at $n$ requires all weights to coincide.}
\expandafter\def\csname statadic@proof@STATADIC-INTEGRATION-001\endcsname{Each $w_i$ is at most the total weight, so
  $\sum w_i^2\le(\sum w_i)^2$. Cauchy--Schwarz against the constant vector
  gives the upper bound. At equality the sum of squared deviations about
  the mean vanishes, hence each deviation is zero.}
\expandafter\def\csname statadic@assumptions@STATADIC-INTEGRATION-001\endcsname{\begin{itemize}
\item Nonnegative finite weights on a nonempty replicate and a positive sum of squared weights.
\end{itemize}}
\expandafter\def\csname statadic@diagnostics@STATADIC-INTEGRATION-001\endcsname{\begin{itemize}
\item Inspect effective\_counts separately for each replicate. A small value flags concentration; a value near draws does not test independence, quadrature bias or rare-event tail coverage.
\end{itemize}}
\expandafter\def\csname statadic@statement@STATADIC-INTEGRATION-002\endcsname{Suppose $R>0$ probability-scale replicate estimates have common
  expectation $p$, common finite variance $v$, and zero covariance across
  distinct replicates. Their average $A$ has mean $p$, variance $v/R$, and
  $\Pr(|A-p|\ge t)\le v/(Rt^2)$ for $t>0$.}
\expandafter\def\csname statadic@proof@STATADIC-INTEGRATION-002\endcsname{Linearity preserves the common expectation. Expand the variance
  of the sum as a double sum of covariances. Only the $R$ diagonal terms survive.
  Divide by $R^2$ and apply Chebyshev's inequality to the resulting average.}
\expandafter\def\csname statadic@assumptions@STATADIC-INTEGRATION-002\endcsname{\begin{itemize}
\item Square-integrable probability-scale replicates have the same target mean, the same finite variance and zero pairwise covariance.
\item Positive replicate count and positive absolute error tolerance.
\end{itemize}}
\expandafter\def\csname statadic@diagnostics@STATADIC-INTEGRATION-002\endcsname{\begin{itemize}
\item Read relative\_standard\_error together with replicate\_log\_probabilities, design.replicates and valid. One replicate yields NaN for this uncertainty estimate. Fixed or repeated point sets do not justify the independent-replicate interpretation.
\item Antithetic draws within a replicate need not be independent. A sample variance from a few replicates is not itself an upper bound on the population variance in the theorem.
\end{itemize}}
\expandafter\def\csname statadic@statement@STATADIC-POISSON-001\endcsname{In an intercept-only Poisson model with positive exposures, write
  $Y=\sum y_i$, $E=\sum e_i>0$ and $\ell(b)=Yb-Ee^b$ up to an additive constant.
  For $Y>0$, $b_*=\log(Y/E)$ is the unique maximizer. If $Y=0$, lowering any
  finite intercept by one strictly increases the likelihood.}
\expandafter\def\csname statadic@proof@STATADIC-POISSON-001\endcsname{The log offset converts $\exp(b+\log e_i)$ to $e_i\exp b$.
  At $b_*$ the likelihood difference for a change $h$ is
  $Y(h+1-\exp h)$, strictly negative for $h\ne0$.
  With no observed counts the remaining term $-E\exp b$ is strictly decreasing.}
\expandafter\def\csname statadic@assumptions@STATADIC-POISSON-001\endcsname{\begin{itemize}
\item Intercept-only Poisson log-link likelihood with known positive row exposures and nonnegative counts.
\item The total exposure is positive. In the interior case the total observed count is positive.
\end{itemize}}
\expandafter\def\csname statadic@diagnostics@STATADIC-POISSON-001\endcsname{\begin{itemize}
\item Check every exposure is positive and finite and count the events before interpreting a fitted rate. An all-zero stratum with its own unconstrained intercept has no finite interior MLE. A small numerical score can occur on the way to this boundary.
\end{itemize}}
\expandafter\def\csname statadic@statement@STATADIC-QUALIFICATION-001\endcsname{Of $n>0$ attempted simulation trials let $V$ be the valid trials
  and $R\subseteq V$ the observed rejections. If $A$ completes the unobserved outcomes
  with $A\cap V=R$, the overall rejection rate satisfies
  $|R|/n\le |A|/n\le(|R|+n-|V|)/n$. Both endpoints are attainable.}
\expandafter\def\csname statadic@proof@STATADIC-QUALIFICATION-001\endcsname{Every completed set contains $R$ and lies in the disjoint union
  $R\cup V^c$. Compare the three counts, divide by the same positive $n$, then
  take the completions $A=R$ and $A=R\cup V^c$ to attain the endpoints.}
\expandafter\def\csname statadic@assumptions@STATADIC-QUALIFICATION-001\endcsname{\begin{itemize}
\item A nonempty finite set of attempted trials, a subset with valid outcomes, and observed rejections only among valid trials.
\end{itemize}}
\expandafter\def\csname statadic@diagnostics@STATADIC-QUALIFICATION-001\endcsname{\begin{itemize}
\item Report requested, valid, failed and rejected counts and use the attempted denominator for the conservative rate. Rejections per valid run can understate operating error when failures are selective. Keep the failure-budget and Monte Carlo uncertainty checks in the qualification report.
\end{itemize}}
\expandafter\def\csname statadic@statement@STATADIC-ROOTENERGY-001\endcsname{Let $d(x,y)=\|\phi(x)-\phi(y)\|^2$, with both first and second
  feature moments integrable for every law in a class $C$. If the feature mean
  $m(P)$ is injective on $C$, the root of unrooted energy satisfies the triangle
  inequality and is zero exactly for equal laws in $C$.}
\expandafter\def\csname statadic@proof@STATADIC-ROOTENERGY-001\endcsname{The population energy identity gives
  $E(P,Q)=2\|m(P)-m(Q)\|^2$. Its square root is $\sqrt 2$ times
  the Hilbert distance. Apply the triangle inequality and injectivity of $m$.}
\expandafter\def\csname statadic@assumptions@STATADIC-ROOTENERGY-001\endcsname{\begin{itemize}
\item The ground is squared Hilbert feature distance; first and second feature moments are integrable on the chosen class.
\item The feature mean is injective on that class of probability measures.
\end{itemize}}
\expandafter\def\csname statadic@diagnostics@STATADIC-ROOTENERGY-001\endcsname{\begin{itemize}
\item Record the ground, scheme and law class. Justify finite moments and separation for that class before treating zero distance as equality. A signed U-statistic or an arbitrary nonnegative square matrix cannot justify this interpretation.
\end{itemize}}
\expandafter\def\csname statadic@statement@STATADIC-PSD-004\endcsname{Let $A+\operatorname{diag}(y)=U\operatorname{diag}(d)U^T$ with $U$
  orthogonal. If its positive part $C=U\operatorname{diag}(\max(0,d))U^T$
  has unit diagonal, then $C$ is positive semidefinite and minimizes
  $\|B-A\|_F^2$ among positive semidefinite matrices $B$ with unit diagonal.}
\expandafter\def\csname statadic@proof@STATADIC-PSD-004\endcsname{Clipping at zero yields the PSD Frobenius projection of the shifted
  target. A diagonal shift changes the squared loss by the same quantity for all
  matrices on the unit-diagonal slice. Subtract the losses to cancel that quantity.}
\expandafter\def\csname statadic@assumptions@STATADIC-PSD-004\endcsname{\begin{itemize}
\item A real input with a diagonal shift admitting the supplied orthogonal eigendecomposition.
\item The positive-part projection at that shift has diagonal exactly equal to one. Competitors are real PSD matrices with unit diagonal.
\end{itemize}}
\expandafter\def\csname statadic@diagnostics@STATADIC-PSD-004\endcsname{\begin{itemize}
\item NearestCorrelationResult.residual measures the Euclidean norm of the positive-part diagonal minus one at the returned dual. converged compares this residual with the configured tolerance and an internal ceiling.
\item An exact zero residual certifies the raw PSD projection. The exposed matrix also undergoes flooring and congruence scaling; unit diagonal of this exposed matrix alone cannot certify the raw optimum.
\end{itemize}}
\expandafter\def\csname statadic@statement@STATADIC-PSD-005\endcsname{Let $F$ be positive definite with $F_{ii}>0$ and desired variances
  $v_i>0$. Define $s_i=\sqrt{v_i}/\sqrt{F_{ii}}$ and $B_{ij}=s_iF_{ij}s_j$.
  Then $B$ is positive definite and $B_{ii}=v_i$. Choosing $v_i=1$ gives the
  final correlation normalization.}
\expandafter\def\csname statadic@proof@STATADIC-PSD-005\endcsname{The diagonal scaling matrix $S$ has no zero diagonal entry, so
  $Sz\ne0$ for nonzero $z$. Hence $z^T SFSz=(Sz)^TF(Sz)>0$. On the diagonal,
  substitute the square roots to obtain $s_i^2F_{ii}=v_i$.}
\expandafter\def\csname statadic@assumptions@STATADIC-PSD-005\endcsname{\begin{itemize}
\item The floored real matrix is positive definite and has positive diagonal. Every requested variance is strictly positive.
\item Use scale factors sqrt(v\ensuremath{{}_i})/sqrt(F\ensuremath{{}_i}\ensuremath{{}_i}).
\end{itemize}}
\expandafter\def\csname statadic@diagnostics@STATADIC-PSD-005\endcsname{\begin{itemize}
\item Check strictly positive finite original variances before calling nearest\_covariance; square roots and standardization otherwise need not produce a finite repair.
\item The floored reconstruction uses a fraction of the largest positive-part eigenvalue. A positive absolute floor and an exact orthogonal eigendecomposition imply positive definiteness before scaling. Check finiteness, requested diagonal and a Cholesky factorization of the returned covariance separately from the pre-floor residual.
\end{itemize}}
\expandafter\def\csname statadic@statement@STATADIC-RANDOMIZATION-004\endcsname{Let a finite group $G$ act by probability-preserving transformations
  under the null. Independently of the data, draw $B$ iid uniform transformations with
  replacement. With an inclusive upper tail and $p=(1+b)/(B+1)$,
  $\Pr(p\le\alpha)\le\alpha$ for every $\alpha\ge0$, including ties.}
\expandafter\def\csname statadic@proof@STATADIC-RANDOMIZATION-004\endcsname{For fixed data $x$, choose a uniform origin $h$ on its orbit.
  Apply the finite sampled-rank counting theorem to the statistic $g\mapsto T(gx)$
  on the group itself. It bounds the average over origins and ordered schedules.
  Integrate the bound over $x$. Each origin preserves the null law, so the resulting
  average equals the rejection probability at the observed origin.}
\expandafter\def\csname statadic@assumptions@STATADIC-RANDOMIZATION-004\endcsname{\begin{itemize}
\item A finite group preserves the measurable null law. All B transformations are iid uniform with replacement and independent of the data.
\end{itemize}}
\expandafter\def\csname statadic@diagnostics@STATADIC-RANDOMIZATION-004\endcsname{\begin{itemize}
\item Record the exchangeability group, sampling unit, tie rule, B and generator design. Repeated transformations are permitted.
\end{itemize}}
\expandafter\def\csname statadic@statement@STATADIC-RANDOMIZATION-005\endcsname{For a constant statistic all transformations tie, and the
  inclusive add-one p-value equals one for every schedule and every draw count.}
\expandafter\def\csname statadic@proof@STATADIC-RANDOMIZATION-005\endcsname{Every transformed statistic equals the observed value. All $B$
  indicators are one, giving $(1+B)/(B+1)=1$.}
\expandafter\def\csname statadic@assumptions@STATADIC-RANDOMIZATION-005\endcsname{\begin{itemize}
\item A constant statistic, any observed point and any schedule of B finite-group transformations.
\end{itemize}}
\expandafter\def\csname statadic@diagnostics@STATADIC-RANDOMIZATION-005\endcsname{\begin{itemize}
\item Report tie frequency and the attainable grid before interpreting p-value histograms.
\end{itemize}}
\expandafter\def\csname statadic@statement@STATADIC-RANDOMIZATION-006\endcsname{On an invariant null probability space let all transformed
  statistics be distinct almost surely. For every integer $0\le k\le |G|$,
  the inclusive full-group p-value satisfies $P(p\le k/|G|)=k/|G|$.}
\expandafter\def\csname statadic@proof@STATADIC-RANDOMIZATION-006\endcsname{Sort the transformed scores at each tie-free observation. The upper
  ranks along that ordering are $|G|,\ldots,1$, so exactly $k$ transformed origins
  reject. Integrate this count, exchange the finite sum with the integral and use
  invariance of the null law to give every origin the same rejection probability.}
\expandafter\def\csname statadic@assumptions@STATADIC-RANDOMIZATION-006\endcsname{\begin{itemize}
\item A finite measure-preserving group action and a measurable statistic with distinct values across transformations almost surely.
\end{itemize}}
\expandafter\def\csname statadic@diagnostics@STATADIC-RANDOMIZATION-006\endcsname{\begin{itemize}
\item Check the realized group order, repeated values and stabilizers before using an equality instead of a level bound.
\end{itemize}}
\expandafter\def\csname statadic@statement@STATADIC-GLM-002\endcsname{For a grouped count $y$ on $m$ trials let
  $p=\exp(\eta)/(1+\exp(\eta))$. For the implemented log likelihood
  $C-y\log(1+e^{-\eta})-(m-y)\log(1+e^\eta)$ the score is $y-mp$
  and its derivative is $-mp(1-p)$.}
\expandafter\def\csname statadic@proof@STATADIC-GLM-002\endcsname{Replace the negative-predictor softplus by the positive one
  minus the predictor. The likelihood becomes $C+y\eta-m\log(1+e^\eta)$.
  Differentiate it, then apply the quotient rule to $p$.}
\expandafter\def\csname statadic@assumptions@STATADIC-GLM-002\endcsname{\begin{itemize}
\item Finite predictor and parameter-independent log-binomial coefficient. For statistical interpretation the grouped count lies between zero and its trial total.
\end{itemize}}
\expandafter\def\csname statadic@diagnostics@STATADIC-GLM-002\endcsname{\begin{itemize}
\item Inspect trial totals, counts, saturation, separation and design rank when interpreting curvature.
\end{itemize}}
\expandafter\def\csname statadic@statement@STATADIC-GLM-003\endcsname{For $\ell(\eta)=y\eta-\exp(\eta)-C$ the score is
  $y-\exp(\eta)$ and its derivative is $-\exp(\eta)$.
  An offset held fixed changes the mean through $\eta$ and has no coefficient derivative.}
\expandafter\def\csname statadic@proof@STATADIC-GLM-003\endcsname{Differentiate the linear count term and the exponential
  cumulant. The log-factorial term is constant in the predictor.}
\expandafter\def\csname statadic@assumptions@STATADIC-GLM-003\endcsname{\begin{itemize}
\item Finite predictor including any fixed log-exposure offset and a parameter-independent log-factorial. For the Poisson model, counts are nonnegative integers and exposures positive.
\end{itemize}}
\expandafter\def\csname statadic@diagnostics@STATADIC-GLM-003\endcsname{\begin{itemize}
\item Check exposure units and positivity, and examine the fit for zero-event strata and boundary estimates.
\end{itemize}}
\expandafter\def\csname statadic@statement@STATADIC-GLMCOV-002\endcsname{On a fixed design, let $a_i=x_i^Tu$ be a contrast direction
  and $S=\sum_i a_i(Y_i-\mu_i)$. If the row residuals have mean zero, are
  square integrable and uncorrelated across rows, with actual variances $v_i$,
  then $E S=0$ and $\mathrm{Var}(S)=\sum_i a_i^2 v_i$.
  A working variance $w_i$ leaves discrepancy $\sum_i a_i^2(v_i-w_i)$.}
\expandafter\def\csname statadic@proof@STATADIC-GLMCOV-002\endcsname{Integrate the finite weighted sum. Expand its variance as a
  double sum of covariances, take each scalar factor outside and discard
  off-diagonal terms. Subtract the working information in the same direction.}
\expandafter\def\csname statadic@assumptions@STATADIC-GLMCOV-002\endcsname{\begin{itemize}
\item Fixed contrast loadings, centered square-integrable residuals and vanishing covariance between distinct rows.
\end{itemize}}
\expandafter\def\csname statadic@diagnostics@STATADIC-GLMCOV-002\endcsname{\begin{itemize}
\item Compare dispersion and residual variances to the working variance. Inspect row clustering before interpreting a row-wise sandwich.
\end{itemize}}
\expandafter\def\csname statadic@statement@STATADIC-GLMCOV-003\endcsname{Let $S_n,R_n$ be the mean scores at the true and fitted
  parameters, and $H_n$ the secant sensitivity between them. Assume full-rank
  sample sensitivities, $H_n\to H$ in probability for nonsingular $H$,
  $\sqrt N S_n\Rightarrow Z$, and $\sqrt N R_n\to0$ in probability.
  Then $\sqrt N(\hat\beta-\beta_0)\Rightarrow H^{-1}Z$.}
\expandafter\def\csname statadic@proof@STATADIC-GLMCOV-003\endcsname{Rowwise secants give $H_n(\hat\beta-\beta_0)=S_n-R_n$.
  Slutsky removes the negligible fitted-score error. Inversion is continuous
  at the nonsingular sensitivity limit. Apply Slutsky again to the inverse
  sensitivity and the scaled score difference.}
\expandafter\def\csname statadic@assumptions@STATADIC-GLMCOV-003\endcsname{\begin{itemize}
\item Scaled true score has a weak limit. Random secant sensitivity converges in probability to an invertible matrix, is invertible at each sampled fit almost surely, and the scaled fitted-score residual vanishes.
\end{itemize}}
\expandafter\def\csname statadic@diagnostics@STATADIC-GLMCOV-003\endcsname{\begin{itemize}
\item Inspect separation, rank, smallest information eigenvalue, scaled gradient and numerical convergence, and justify the score CLT using the sampling units.
\end{itemize}}
\expandafter\def\csname statadic@statement@STATADIC-GLMCOV-004\endcsname{Suppose the reported sensitivity $A_n$, fitted-score meat $B_n$
  and multiplier $c_n$ converge jointly in probability to $A,B,1$ with nonsingular $A$.
  Then $c_n A_n^{-1} B_n A_n^{-T}\to A^{-1}BA^{-T}$ in probability.
  For the model-based choice $B=A=A^T$ the limit reduces to $A^{-1}$.}
\expandafter\def\csname statadic@proof@STATADIC-GLMCOV-004\endcsname{Matrix inversion is continuous at $A$; multiplication,
  transposition and scalar scaling are continuous. Apply continuous mapping
  at the joint limit. For model information use symmetry and cancel the
  inverse with $A$.}
\expandafter\def\csname statadic@assumptions@STATADIC-GLMCOV-004\endcsname{\begin{itemize}
\item Joint convergence in probability of normalized fitted sensitivity, fitted meat and the HC factor to A, B and one. A is nonsingular and the estimators are measurable.
\end{itemize}}
\expandafter\def\csname statadic@diagnostics@STATADIC-GLMCOV-004\endcsname{\begin{itemize}
\item Check the bread rank, working variance model, independence of score units, and the ratio of parameters to observations for HC1.
\end{itemize}}
\expandafter\def\csname statadic@statement@STATADIC-GLMCOV-005\endcsname{For a Gaussian limiting score $Z$ with covariance $B$ and
  nonsingular sensitivity $A$, the coefficient limit $A^{-1}Z$ is Gaussian
  with covariance $A^{-1} B A^{-T}$. Together with the fitted sandwich limit
  this justifies the covariance used for asymptotic contrasts.}
\expandafter\def\csname statadic@proof@STATADIC-GLMCOV-005\endcsname{Apply the Gaussian linear-response identity with response
  matrix $-A^{-1}$. The sign in the score response cancels, and the two minus
  signs in its covariance cancel as well.}
\expandafter\def\csname statadic@assumptions@STATADIC-GLMCOV-005\endcsname{\begin{itemize}
\item A Gaussian score limit with known coordinate covariances B and finite second moments. The preceding estimating-equation limit identifies the inverse of a nonsingular sensitivity A as its coefficient response.
\end{itemize}}
\expandafter\def\csname statadic@diagnostics@STATADIC-GLMCOV-005\endcsname{\begin{itemize}
\item Justify the score central limit theorem and identify the covariance for the sampling units before choosing Gaussian critical values.
\end{itemize}}
\expandafter\def\csname statadic@statement@STATADIC-SURVIVAL-003\endcsname{Let $T,C$ have independent laws $F,G$. Enter at fixed $a$
  only when $T,C\ge a$, with positive probability $H=F[a,\infty)G[a,\infty)$.
  For a measurable time window $B$, the conditional event probability is
  $H^{-1}\int_{B\cap[a,\infty)}G[t,\infty)\,dF(t)$.
  The conditional censor probability is
  $H^{-1}\int_{B\cap[a,\infty)}F(c,\infty)\,dG(c)$. Ties are assigned to events.}
\expandafter\def\csname statadic@proof@STATADIC-SURVIVAL-003\endcsname{Independence gives a product measure. Slice the event set
  at fixed failure time and the censor set at fixed censor time. The surviving
  opposite coordinate lies in the corresponding upper interval; integrate each
  slice by Tonelli. Both events imply entry. Divide by the product probability
  of the two entry conditions.}
\expandafter\def\csname statadic@assumptions@STATADIC-SURVIVAL-003\endcsname{\begin{itemize}
\item Independent failure and censoring times at fixed covariates, a fixed entry threshold and positive probability of both times surviving to entry. Ties count as events.
\end{itemize}}
\expandafter\def\csname statadic@diagnostics@STATADIC-SURVIVAL-003\endcsname{\begin{itemize}
\item Check entry precedes recorded exit. State the covariates conditioning independence, the time origin, the treatment of ties and whether entry or censoring is informative.
\end{itemize}}
\expandafter\def\csname statadic@statement@STATADIC-SPATIAL-004\endcsname{Suppose the population error moments at $\lambda_0$ vanish.
  Write their change at $r$ as $a_j(r-\lambda_0)+b_j(r^2-\lambda_0^2)$,
  absorbing the sign convention into the coefficients. On an admissible
  set $S$ containing $\lambda_0$, the parameter is uniquely identified exactly
  when every other $r\in S$ has some $a_j+b_j(r+\lambda_0)\ne0$.}
\expandafter\def\csname statadic@proof@STATADIC-SPATIAL-004\endcsname{Factor each moment change into
  $(r-\lambda_0)[a_j+b_j(r+\lambda_0)]$. At a distinct candidate the
  first factor is nonzero, so all moments vanish precisely when every second
  factor vanishes.}
\expandafter\def\csname statadic@assumptions@STATADIC-SPATIAL-004\endcsname{\begin{itemize}
\item Population quadratic moments vanish at the true admissible parameter. Their changes have linear and quadratic coefficients a\ensuremath{{}_j} and b\ensuremath{{}_j}.
\end{itemize}}
\expandafter\def\csname statadic@diagnostics@STATADIC-SPATIAL-004\endcsname{\begin{itemize}
\item Search the entire admissible interval for common roots. Compare the moment map at competing minima, alongside local Jacobian and rank checks.
\end{itemize}}
\expandafter\def\csname statadic@statement@STATADIC-SPATIAL-005\endcsname{On $(-1,1)$ the quadratic moment $m(r)=r^2-1/4$
  has two roots, $1/2$ and $-1/2$, although its derivative at $1/2$ is nonzero.}
\expandafter\def\csname statadic@proof@STATADIC-SPATIAL-005\endcsname{Substitute the two candidates and differentiate $r^2-1/4$.
  Both roots are strictly interior, and the derivative at the positive one is one.}
\expandafter\def\csname statadic@assumptions@STATADIC-SPATIAL-005\endcsname{\begin{itemize}
\item One quadratic moment r\ensuremath{{}^2}\ensuremath{-}1/4, with the admitted interval (\ensuremath{-}1,1).
\end{itemize}}
\expandafter\def\csname statadic@diagnostics@STATADIC-SPATIAL-005\endcsname{\begin{itemize}
\item Inspect alternative minima across the interval even when the selected root is interior and locally regular.
\end{itemize}}
\expandafter\def\csname statadic@statement@STATADIC-SPATIAL-006\endcsname{With the same spatial matrix $W$, let $A(r)=I-rW$.
  Then $A(r)A(s)=A(s)A(r)$. In a zero-mean SARAR system with the same spherical
  innovation law, swapping lag and error coefficients leaves the composed
  innovation filter and hence the implied covariance unchanged whenever the filters are invertible.}
\expandafter\def\csname statadic@proof@STATADIC-SPATIAL-006\endcsname{Multiply out the filters. The coefficients of $W$ and $W^2$
  are $-(r+s)$ and $rs$, invariant to the swap. Inverting the identical
  composed filter gives the same covariance factor.}
\expandafter\def\csname statadic@assumptions@STATADIC-SPATIAL-006\endcsname{\begin{itemize}
\item Common fixed graph for lag and error, no distinguishing mean signal, and invertible filters for the covariance interpretation with unit spherical innovation covariance.
\end{itemize}}
\expandafter\def\csname statadic@diagnostics@STATADIC-SPATIAL-006\endcsname{\begin{itemize}
\item Inspect graph equality, exogenous mean variation and swapped parameter pairs on a profiled objective.
\end{itemize}}
\expandafter\def\csname statadic@statement@STATADIC-SELECTION-001\endcsname{Let training and evaluation data have joint law $P\otimes Q$.
  Let $N$ be the measurable set of training outcomes whose selected hypothesis
  is true. If the evaluation rejection probability for almost every $t\in N$
  is at most $\alpha$, then the probability of reporting a false rejection is
  at most $\alpha P(N)$. The selection rule may be continuous-valued.}
\expandafter\def\csname statadic@proof@STATADIC-SELECTION-001\endcsname{Restrict the rejection set to training outcomes in $N$.
  Tonelli expresses its probability as the integral over $N$ of the held-out
  rejection probabilities. Apply their almost-sure bound and integrate
  the constant.}
\expandafter\def\csname statadic@assumptions@STATADIC-SELECTION-001\endcsname{\begin{itemize}
\item The joint law factors into training and evaluation laws. The selected-null set and rejection event are measurable, and almost every true selected null has held-out rejection probability at most \ensuremath{\alpha}.
\end{itemize}}
\expandafter\def\csname statadic@diagnostics@STATADIC-SELECTION-001\endcsname{\begin{itemize}
\item Document how independent units entered the split and justify the fixed-weight held-out test for its null. Splitting a dependent time series into contiguous halves does not establish independence.
\end{itemize}}
\expandafter\def\csname statadic@statement@STATADIC-SELECTION-002\endcsname{For two independent p-values uniform on
  $\{1/20,\ldots,1\}$, either one individually rejects at $1/20$ in
  exactly $20$ of $400$ equally likely pairs. Reporting the smaller p-value
  rejects in $39$ of those pairs.}
\expandafter\def\csname statadic@proof@STATADIC-SELECTION-002\endcsname{In the $20$ by $20$ grid the first row or the first column
  rejects. Their intersection consists of the single pair of minimum ranks.}
\expandafter\def\csname statadic@assumptions@STATADIC-SELECTION-002\endcsname{\begin{itemize}
\item Two independent uniform null p-values on the twenty-point grid. The analyst reports the smaller and tests at 1/20.
\end{itemize}}
\expandafter\def\csname statadic@diagnostics@STATADIC-SELECTION-002\endcsname{\begin{itemize}
\item Record whether weighting, selection or tuning inspected the inference outcomes before reporting the test.
\end{itemize}}
\expandafter\def\csname statadic@statement@STATADIC-WALD-001\endcsname{Let $X_n$ be a scaled contrast error and $D_n$ its positive
  estimated scale in the same units. If $X_n\Rightarrow\sigma Z$, $D_n\to\sigma$
  in probability, and $\sigma>0$, then $X_n/D_n\Rightarrow Z$.
  For every measurable acceptance region $K$ whose boundary has zero limiting
  probability, $P(X_n/D_n\in K)\to P(Z\in K)$. Taking $K=[-c,c]$ gives
  the limiting coverage of the corresponding two-sided interval.}
\expandafter\def\csname statadic@proof@STATADIC-WALD-001\endcsname{The reciprocal is continuous at the positive scale limit.
  Slutsky gives convergence of the product of the contrast error and the
  reciprocal scale. Apply the continuity-set implication of the portmanteau
  theorem to the laws of these studentized errors.}
\expandafter\def\csname statadic@assumptions@STATADIC-WALD-001\endcsname{\begin{itemize}
\item A scaled contrast error converges in distribution to \ensuremath{\sigma}Z and its estimated scale converges in probability to \ensuremath{\sigma}>0. The acceptance set is measurable and its boundary has zero limiting mass.
\end{itemize}}
\expandafter\def\csname statadic@diagnostics@STATADIC-WALD-001\endcsname{\begin{itemize}
\item Check positive scale, adequate independent units, justified coefficient CLT, consistent variance and the reported reference quantile.
\end{itemize}}
\expandafter\def\csname statadic@statement@STATADIC-GMM-004\endcsname{Write population moments as $g(b)=a+Db$. Suppose
  $g(b_0)=0$ and $D^T W D$ is nonsingular. Then $g(b)=0$ exactly when
  $b=b_0$. For any $a$ and $b_0$, the coefficient defined by the weighted
  normal equations obeys $\hat b-b_0=-L(a+Db_0)$ with
  $L=(D^T W D)^{-1}D^T W$.}
\expandafter\def\csname statadic@proof@STATADIC-GMM-004\endcsname{Invertibility gives $LD=I$. Expand $-L(a+Db_0)$ to
  obtain the error identity. Two moment zeros therefore coincide. For spatial
  IV this calculation needs the population instrument-error moment to vanish
  as well as rank of the projected population regressors.}
\expandafter\def\csname statadic@assumptions@STATADIC-GMM-004\endcsname{\begin{itemize}
\item Fixed finite real moment and coefficient axes, affine population mean moments a+Db, and an invertible D\ensuremath{{}^T}WD.
\item For identification of the structural coefficient, a+Db\ensuremath{{}_0}=0. For a convex minimizer interpretation, use the symmetric positive definite weight accepted by the fit.
\end{itemize}}
\expandafter\def\csname statadic@diagnostics@STATADIC-GMM-004\endcsname{\begin{itemize}
\item Inspect fit.rank and fit.weighted\_rank against the parameter count. Both are numerical ranks on the observed design; a nearly singular population derivative can pass both tests.
\item State why the instrument and structural error are orthogonal in the population. Neither a small sample first\_order\_residual nor strong sample rank can certify this identifying assumption.
\end{itemize}}
\expandafter\def\csname statadic@statement@STATADIC-GMM-005\endcsname{If $D_n\to D$, $W_n\to W$ and the estimated moment covariance
  $S_n\to S$ on a sample path, and $D^TWD$ is nonsingular, then
  $L_nS_nL_n^T\to LSL^T$ for the affine GMM response map $L_n$.}
\expandafter\def\csname statadic@proof@STATADIC-GMM-005\endcsname{Finite matrix products and transposes preserve limits. Inversion
  is continuous at the nonsingular limiting Gram matrix, giving $L_n\to L$.
  Apply continuity again to the two outer products.}
\expandafter\def\csname statadic@assumptions@STATADIC-GMM-005\endcsname{\begin{itemize}
\item The sample derivative, actual moment weight and moment-covariance estimator converge to fixed matrices along the same sample path.
\item The limiting weighted Gram is nonsingular.
\end{itemize}}
\expandafter\def\csname statadic@diagnostics@STATADIC-GMM-005\endcsname{\begin{itemize}
\item Record the actual retained effective\_weight, its provenance and the variability scheme. A finite or converged coefficient solve does not establish correct moments, a valid HAC bandwidth or a consistently estimated efficient weight.
\item The theorem works on the square-root-sample-size coefficient scale. For a coefficient standard error, divide the limiting covariance by the observation count and account separately for any finite-sample correction.
\end{itemize}}
\expandafter\def\csname statadic@statement@STATADIC-GMM-006\endcsname{Let $\bar g_n(b)=a_n+D_n b$. Suppose the finite arrays
  $(a_n,D_n,W_n)$ converge jointly in probability to $(a,D,W)$, the
  population moment satisfies $a+Db_0=0$, and $D^TWD$ is invertible.
  Then the measurable affine GMM coefficient $-L_n a_n$ converges in probability to $b_0$.
  With a symmetric positive definite weight this is the identified weighted normal-equation solution.}
\expandafter\def\csname statadic@proof@STATADIC-GMM-006\endcsname{Matrix products and transposes are continuous, and inversion
  is continuous at the nonsingular limiting Gram matrix. The continuous
  mapping principle at the constant joint limit gives convergence to $-La$.
  The identity $LD=I$ and the population moment identify this limit as $b_0$.}
\expandafter\def\csname statadic@assumptions@STATADIC-GMM-006\endcsname{\begin{itemize}
\item The sample mean intercept, derivative and retained weight converge jointly in probability to fixed limits. Coefficient estimates are measurable.
\item Correct population moments and invertible limiting weighted Gram. Positive definite weights give the native quadratic-minimization interpretation.
\end{itemize}}
\expandafter\def\csname statadic@diagnostics@STATADIC-GMM-006\endcsname{\begin{itemize}
\item Inspect fit.rank and fit.weighted\_rank against the parameter count. Both are numerical ranks on the observed design; a nearly singular population derivative can pass both tests.
\item Record the actual retained effective\_weight, its provenance and the variability scheme. A finite or converged coefficient solve does not establish correct moments, a valid HAC bandwidth or a consistently estimated efficient weight.
\item State why the instrument and structural error are orthogonal in the population. Neither a small sample first\_order\_residual nor strong sample rank can certify this identifying assumption.
\end{itemize}}
\expandafter\def\csname statadic@statement@STATADIC-GMM-007\endcsname{Suppose $\sqrt{N}\,\bar g_n(b_0)$ converges in law to $Z$
  and $(D_n,W_n)$ converges in probability to $(D,W)$ with nonsingular
  $D^TWD$. Assume measurable response maps and almost surely full-rank
  sample weighted Grams. Then $\sqrt{N}(\hat b_n-b_0)$ converges in law
  to $-LZ$, $L=(D^TWD)^{-1}D^TW$, with $N=n+1$.}
\expandafter\def\csname statadic@proof@STATADIC-GMM-007\endcsname{With probability one the exact normal equations give
  $\hat b_n-b_0=-L_n\bar g_n(b_0)$. The matrix LLN and continuity at the
  population Gram imply $L_n\to L$ in probability. Slutsky's theorem and
  continuity of matrix-vector multiplication transfer the assumed score CLT.
  A Gaussian score limit consequently yields a Gaussian coefficient limit.}
\expandafter\def\csname statadic@assumptions@STATADIC-GMM-007\endcsname{\begin{itemize}
\item Normalized true-parameter sample moments converge in distribution, and sample derivatives and weights converge jointly in probability to fixed limits.
\item The limiting weighted Gram is invertible; the weighted Gram of each admitted sampling problem is nonsingular almost surely. The response maps are measurable.
\end{itemize}}
\expandafter\def\csname statadic@diagnostics@STATADIC-GMM-007\endcsname{\begin{itemize}
\item Record the actual retained effective\_weight, its provenance and the variability scheme. A finite or converged coefficient solve does not establish correct moments, a valid HAC bandwidth or a consistently estimated efficient weight.
\item State why the instrument and structural error are orthogonal in the population. Neither a small sample first\_order\_residual nor strong sample rank can certify this identifying assumption.
\item Specify a sampling or dependence argument for a centered score CLT. A numerical Wald statistic does not test that CLT.
\end{itemize}}
\expandafter\def\csname statadic@statement@STATADIC-GMM-008\endcsname{If $Z$ has a Gaussian law with square-integrable coordinates
  and covariance $S$, the law of $-LZ$ is Gaussian. Its coordinate covariance
  is $LSL^T$, including when the limit is singular.}
\expandafter\def\csname statadic@proof@STATADIC-GMM-008\endcsname{Gaussian laws are closed under continuous linear maps.
  Expand each output coordinate as a finite linear combination and use bilinearity
  of covariance; the two negative signs cancel.}
\expandafter\def\csname statadic@assumptions@STATADIC-GMM-008\endcsname{\begin{itemize}
\item A fixed real response matrix and a multivariate Gaussian score law with square-integrable coordinates and specified covariance S.
\end{itemize}}
\expandafter\def\csname statadic@diagnostics@STATADIC-GMM-008\endcsname{\begin{itemize}
\item The theorem works on the square-root-sample-size coefficient scale. For a coefficient standard error, divide the limiting covariance by the observation count and account separately for any finite-sample correction.
\item Inspect covariance validity and the rank of the tested contrast covariance separately from design rank. A singular Gaussian limit still has the displayed covariance.
\end{itemize}}
\expandafter\def\csname statadic@statement@STATADIC-GMM-009\endcsname{Let the moment covariance $S$ be symmetric positive definite,
  with $D^T S^{-1}D$ nonsingular. Put $E=(D^T S^{-1}D)^{-1}D^TS^{-1}$.
  For every alternative response $A$ satisfying $AD=I$,
  \[
    ASA^T-ESE^T=(A-E)S(A-E)^T .
  \]
  Thus each contrast variance under $E$ is at most its variance under $A$.}
\expandafter\def\csname statadic@proof@STATADIC-GMM-009\endcsname{Both $A$ and $E$ reproduce the coefficients, hence
  $(A-E)D=0$. Since $ES=(D^TS^{-1}D)^{-1}D^T$, the two mixed terms
  in the expansion vanish. The remaining quadratic form is nonnegative
  from the moment covariance assumption.}
\expandafter\def\csname statadic@assumptions@STATADIC-GMM-009\endcsname{\begin{itemize}
\item The moment covariance S is symmetric, nonsingular and positive semidefinite. D\ensuremath{{}^T}S\ensuremath{{}^{-1}}D is nonsingular.
\item The alternative response A reproduces coefficients, AD=I. The same population moments are used for both estimates.
\end{itemize}}
\expandafter\def\csname statadic@diagnostics@STATADIC-GMM-009\endcsname{\begin{itemize}
\item Record weight\_fixed\_point\_error and solver.converged at the retained coefficient. A small fixed-point error assesses the empirical relation W Omega \ensuremath{\approx} I and does not establish long-run covariance consistency.
\end{itemize}}
\expandafter\def\csname statadic@statement@STATADIC-GMM-010\endcsname{Each observation supplies the intercept $a_i$ and derivative
  $D_i$ of its affine moments. Assume the joint rows have a common law, are
  pairwise independent across observations, and have a finite first moment.
  Let $(a,D)=E[(a_i,D_i)]$, $a+Db_0=0$, and $D^T W D$ nonsingular.
  For any weight sequence $W_n\to W$ almost surely, the GMM coefficients of
  the row averages converge to $b_0$ almost surely.}
\expandafter\def\csname statadic@proof@STATADIC-GMM-010\endcsname{Apply the strong law to the joint row vector, which allows
  dependence between its entries. The mean intercept and mean derivative
  therefore converge on the same probability-one set. On its intersection
  with weight convergence, use continuity of the response map and the
  population moment equation.}
\expandafter\def\csname statadic@assumptions@STATADIC-GMM-010\endcsname{\begin{itemize}
\item The joint observation rows of affine intercepts and derivatives have a common law, are pairwise independent across rows and have a finite first moment. Entries within a row may depend on each other.
\item The retained weight converges almost surely to a deterministic matrix. The population mean moment vanishes and the limiting weighted Gram is invertible.
\end{itemize}}
\expandafter\def\csname statadic@diagnostics@STATADIC-GMM-010\endcsname{\begin{itemize}
\item State why the instrument and structural error are orthogonal in the population. Neither a small sample first\_order\_residual nor strong sample rank can certify this identifying assumption.
\item Document the independent sampling unit and whether any observations share shocks. Clustered or time-series rows may instead use the general convergence-in-probability theorem with a justified dependent-data LLN.
\end{itemize}}
\expandafter\def\csname statadic@statement@STATADIC-IV-002\endcsname{For $g(b)=a+Db$ with $D^TWD$ invertible, suppose at a
  structural value $b_0$ the moment violation equals $D h$ for some $h\ne0$.
  The weighted coefficient equals $b_0-h\ne b_0$, and every population
  moment vanishes there, even when there are more moments than coefficients.}
\expandafter\def\csname statadic@proof@STATADIC-IV-002\endcsname{Substitute $g(b_0)=Dh$ into $\hat b-b_0=-Lg(b_0)$.
  Since $LD=I$, the error is $-h$. Subtract $Dh$ from $g(b_0)$ to see
  that the misspecified coefficient fits all moment equations.}
\expandafter\def\csname statadic@assumptions@STATADIC-IV-002\endcsname{\begin{itemize}
\item The weighted derivative Gram is invertible and the moment violation at the structural value is Dh for a nonzero h.
\end{itemize}}
\expandafter\def\csname statadic@diagnostics@STATADIC-IV-002\endcsname{\begin{itemize}
\item State why the instrument and structural error are orthogonal in the population. Neither a small sample first\_order\_residual nor strong sample rank can certify this identifying assumption.
\item Even several rank-valid instruments may have compatible violations. Report the exclusion argument and sensitivity to instrument sets in addition to sample tests.
\end{itemize}}
\expandafter\def\csname statadic@statement@STATADIC-IV-003\endcsname{Let $g(b)=a+db$ with $d\ne0$. Perturbing $a$ by $e$
  changes its zero by $-e/d$. In particular, for any tolerance $t>0$ and
  requested coefficient change $M>0$, one can choose a nonzero $d$ and a
  perturbation $|e|\le t$ for which the change has magnitude at least $M$.}
\expandafter\def\csname statadic@proof@STATADIC-IV-003\endcsname{Subtract the two scalar solutions. Take $e=t$ and $d=t/M$,
  which meets the rank condition and attains magnitude $M$.}
\expandafter\def\csname statadic@assumptions@STATADIC-IV-003\endcsname{\begin{itemize}
\item A scalar affine equation with nonzero derivative for the identification interpretation. The perturbation and desired coefficient displacement may be arbitrarily small and large, respectively.
\end{itemize}}
\expandafter\def\csname statadic@diagnostics@STATADIC-IV-003\endcsname{\begin{itemize}
\item Inspect fit.rank and fit.weighted\_rank against the parameter count. Both are numerical ranks on the observed design; a nearly singular population derivative can pass both tests.
\item For SpatialIV inspect first\_stage\_rank and instrument\_rank and examine scaled first-stage sensitivity. A boundary or rank failure precludes the standard interpretation.
\end{itemize}}
\expandafter\def\csname statadic@statement@STATADIC-JTEST-001\endcsname{Let $P_n=I-D_n L_n$ and assume almost-sure sample rank.
  Suppose the normalized true-parameter moments $\sqrt{N}\bar g_n(b_0)$
  converge in law to $Z$ and $(D_n,W_n)$ converge in probability to $(D,W)$,
  with $D^TWD$ nonsingular. With measurable residual weight maps, the
  statistic $N\bar g_n(\hat b_n)^TW_n\bar g_n(\hat b_n)$ converges in law
  to $Z^T P^T W P Z$, where $P=I-DL$.}
\expandafter\def\csname statadic@proof@STATADIC-JTEST-001\endcsname{The identity $L_n D_n=I$ gives
  $\bar g_n(\hat b_n)=P_n\bar g_n(b_0)$. The mapping from
  $(D,W)$ to $P^T W P$ is continuous at any nonsingular weighted Gram.
  Slutsky's theorem then transfers the score limit through the quadratic form.}
\expandafter\def\csname statadic@assumptions@STATADIC-JTEST-001\endcsname{\begin{itemize}
\item The normalized true-parameter moments converge in law, and the sample derivatives and weights converge in probability.
\item The limiting weighted Gram is invertible, sample weighted Grams are almost surely nonsingular, and the residual-weight random matrices are measurable.
\end{itemize}}
\expandafter\def\csname statadic@diagnostics@STATADIC-JTEST-001\endcsname{\begin{itemize}
\item Inspect the explicit correct-specification assumption, positive degrees of freedom, weight provenance, weight\_fixed\_point\_error, covariance rank and solver status. An unavailable test cannot be interpreted as a nonrejection.
\end{itemize}}
\expandafter\def\csname statadic@statement@STATADIC-JTEST-002\endcsname{In a finite-dimensional real inner product space $E$, let
  $K$ be the fitted tangent subspace. For a standard Gaussian $G$, the law of
  $\|\operatorname{proj}_{K^\perp}G\|^2$ is the law of a sum of
  $\dim E-\dim K$ independent squared standard normals. This is the
  chi-square reference for an efficiently whitened residual projection.}
\expandafter\def\csname statadic@proof@STATADIC-JTEST-002\endcsname{The Gaussian characteristic function at a vector in
  $K^\perp$ is unchanged by orthogonal projection, so the projected law is
  standard Gaussian on that subspace. Express it in any orthonormal basis:
  the coordinates are independent standard normals and the squared norm is
  their sum of squares. Orthogonal dimensions add to the ambient dimension.}
\expandafter\def\csname statadic@assumptions@STATADIC-JTEST-002\endcsname{\begin{itemize}
\item A centered, whitened standard Gaussian score limit in a finite-dimensional real inner product space. The fitted directions form a linear subspace K.
\end{itemize}}
\expandafter\def\csname statadic@diagnostics@STATADIC-JTEST-002\endcsname{\begin{itemize}
\item Record the number of moments and parameters. In the full-column-rank case their difference is the residual dimension; rank loss changes this argument. The sample covariance and efficient-weight checks are additional gates.
\end{itemize}}
\expandafter\def\csname statadic@statement@STATADIC-JTEST-003\endcsname{Suppose a matrix $C$ whitens the limiting score law
  ($CZ$ is standard Gaussian), the actual limiting weight is $W=C^T C$,
  and $D^T W D$ is nonsingular. For $P=I-DL$, the law of
  $Z^T P^T W P Z$ is chi-square with $m-p$ degrees of freedom.}
\expandafter\def\csname statadic@proof@STATADIC-JTEST-003\endcsname{The whitened residual $CPz$ is orthogonal to the columns of
  $CD$ by the weighted normal equations. The fitted component $Cz-CPz$
  lies in their span, so $CPz$ is the orthogonal residual projection of
  $Cz$. The weighted Gram gives a left inverse for $CD$, hence this span
  has dimension $p$. Apply the standard Gaussian projection law.}
\expandafter\def\csname statadic@assumptions@STATADIC-JTEST-003\endcsname{\begin{itemize}
\item The limiting weight has the form C\ensuremath{{}^T}C and the same C maps the centered limiting score law to the standard multivariate Gaussian.
\item The population weighted Gram D\ensuremath{{}^T}C\ensuremath{{}^T}CD is nonsingular.
\end{itemize}}
\expandafter\def\csname statadic@diagnostics@STATADIC-JTEST-003\endcsname{\begin{itemize}
\item Check the correct-specification assumption, positive moment-minus-parameter count, full covariance rank, efficient weighting provenance and fixed-point error. An accepted J statistic cannot identify which instrument is exogenous.
\end{itemize}}
\expandafter\def\csname statadic@statement@STATADIC-GMM-011\endcsname{Suppose finite-dimensional true-parameter moment vectors
  $g_i$ are iid, have finite second moment and mean zero. Let $S$ be their
  covariance, expressed as $\operatorname{Var}(t^Tg_i)=t^T S t$ for all
  contrasts $t$. Then $n^{-1/2}\sum_{i=1}^n g_i$ converges in law to
  $N(0,S)$, including singular $S$.}
\expandafter\def\csname statadic@proof@STATADIC-GMM-011\endcsname{Project every row onto a fixed contrast. The scalar projected
  variables remain iid, square integrable and centered, with the specified
  variance. Apply the scalar iid CLT and identify the contrast law of the
  centered multivariate Gaussian by its covariance form. Cramér–Wold gives
  the vector limit.}
\expandafter\def\csname statadic@assumptions@STATADIC-GMM-011\endcsname{\begin{itemize}
\item Finite-dimensional iid true-parameter moment vectors have mean zero and finite second moment.
\item S is their positive semidefinite covariance, specified by the variance of every linear contrast.
\end{itemize}}
\expandafter\def\csname statadic@diagnostics@STATADIC-GMM-011\endcsname{\begin{itemize}
\item Identify the independent sampling unit, check finite second moments, and justify population instrument-error orthogonality. A sample condition number or solver success does not establish this CLT.
\end{itemize}}
\expandafter\def\csname statadic@statement@STATADIC-SURVIVAL-004\endcsname{Let $t>0$, log scale $\eta$, log shape $a$, and
  $z=e^a(\log t-\eta)$. A unit-exponential cumulative-hazard threshold
  gives $S(t)=1-F_{\rm Exp(1)}(e^z)=\exp(-e^z)$.
  The derivative of this survivor is the negative of the density with
  log value $a-\log t+z-e^z$.}
\expandafter\def\csname statadic@proof@STATADIC-SURVIVAL-004\endcsname{Evaluate the unit-exponential survival at its positive
  cumulative hazard $e^z$. Differentiate through the hazard and log time,
  using $z'(t)=e^a/t$, then collect the exponential factors.}
\expandafter\def\csname statadic@assumptions@STATADIC-SURVIVAL-004\endcsname{\begin{itemize}
\item A positive duration, finite log scale and log shape, and a unit exponential threshold for the cumulative hazard.
\end{itemize}}
\expandafter\def\csname statadic@diagnostics@STATADIC-SURVIVAL-004\endcsname{\begin{itemize}
\item Check positive durations and ordered entry times. Record the distribution and ancillary log parameter. Compare fitted survival with observed survival over the supported range, and justify the censoring/entry assumptions independently.
\end{itemize}}
\expandafter\def\csname statadic@statement@STATADIC-SURVIVAL-005\endcsname{For standard normal $Z$, write $T=\exp(\eta-e^a Z)$,
  which has lognormal location $\eta$ and log standard deviation $a$.
  At $t>0$ let $z=(\log t-\eta)e^{-a}$. Then
  $P(T\ge t)=\Phi(-z)$, and $-S'(t)$ is positive with log value
  $-\log t-a-z^2/2-\log(2\pi)/2$.}
\expandafter\def\csname statadic@proof@STATADIC-SURVIVAL-005\endcsname{The event $T\ge t$ is $Z\le -z$ by monotonicity of the
  exponential. Differentiate the normal CDF by the fundamental theorem of
  calculus, then differentiate its log-time argument. Insert the standard
  Gaussian density and split the logarithms of its normalizing constant.}
\expandafter\def\csname statadic@assumptions@STATADIC-SURVIVAL-005\endcsname{\begin{itemize}
\item A positive duration and a log-time location minus a positive scale times a standard Gaussian variable.
\end{itemize}}
\expandafter\def\csname statadic@diagnostics@STATADIC-SURVIVAL-005\endcsname{\begin{itemize}
\item Check positive durations and ordered entry times. Record the distribution and ancillary log parameter. Compare fitted survival with observed survival over the supported range, and justify the censoring/entry assumptions independently.
\end{itemize}}
\expandafter\def\csname statadic@statement@STATADIC-HAC-003\endcsname{Write $\gamma_0$ for the contemporaneous score moment and
  $\gamma_h$ for the sum of its two orientations at positive lag. If
  $\sum_h |\gamma_h|<\infty$ and $b_n\to\infty$, then
  $\sum_h (1-h/(b_n+1))_+(1-h/(n+1))_+\gamma_h\to\sum_h\gamma_h$.
  The second factor accounts for the incomplete pairs at the sample edges.}
\expandafter\def\csname statadic@proof@STATADIC-HAC-003\endcsname{Both tapers lie in $[0,1]$ and tend to one at every fixed lag.
  Dominate the summands by the absolutely summable lag moments and pass the
  limit through the series.}
\expandafter\def\csname statadic@assumptions@STATADIC-HAC-003\endcsname{\begin{itemize}
\item The contemporaneous moment and the sums of the two oriented positive-lag moments are absolutely summable.
\item The integer bandwidth tends to infinity and the number of chronological rows grows.
\end{itemize}}
\expandafter\def\csname statadic@diagnostics@STATADIC-HAC-003\endcsname{\begin{itemize}
\item Record the score definition, chronological order, common row count and bandwidth. Plain Newey–West uses raw products, so its covariance interpretation requires zero-mean scores. STATADIC uses the divisor N at every lag and weight 1-h/(b+1).
\item Recompute with plausible bandwidths and inspect the contribution of long lags. A finite, PSD sample matrix does not show that the omitted population tail is negligible.
\end{itemize}}
\expandafter\def\csname statadic@statement@STATADIC-HAC-004\endcsname{Write $X_{h,t}$ for one raw lag product at lag zero and
  the sum of the two oriented score products at positive lag. Let $N=n+1$.
  Assume square integrability and, for all $h,k,t,m$,
  $\sum_{s<m}|\operatorname{Cov}(X_{h,t},X_{k,s})|\le C$.
  The variance of the sample Bartlett entry with bandwidth $b$ and
  common denominator $N$ is at most $C(b+1)^2/N$.}
\expandafter\def\csname statadic@proof@STATADIC-HAC-004\endcsname{Expand the variance into covariances of time-indexed lag
  products. For every pair of lags, the sum of absolute covariances is at
  most $NC$. Multiplying by two taper weights in $[0,1]$ preserves the upper
  bound. Sum over $(b+1)^2$ lag pairs and divide by $N^2$.}
\expandafter\def\csname statadic@assumptions@STATADIC-HAC-004\endcsname{\begin{itemize}
\item The lag product X(h,t) is square-integrable; at positive lags it includes both orientations.
\item Uniformly in h,k,t,m, the absolute covariances of X(h,t) with X(k,s), summed over s<m, are bounded by C.
\end{itemize}}
\expandafter\def\csname statadic@diagnostics@STATADIC-HAC-004\endcsname{\begin{itemize}
\item Record the score definition, chronological order, common row count and bandwidth. Plain Newey–West uses raw products, so its covariance interpretation requires zero-mean scores. STATADIC uses the divisor N at every lag and weight 1-h/(b+1).
\item Check estimand-relevant score dependence rather than inspecting only raw outcome autocorrelations. A short-looking score autocorrelation plot cannot establish a uniform fourth-order envelope.
\end{itemize}}
\expandafter\def\csname statadic@statement@STATADIC-HAC-005\endcsname{With the same bounded lag-product covariance sums, suppose
  $E X_{h,t}=\gamma_h$ at every time, $\sum_h|\gamma_h|<\infty$,
  $b_n\to\infty$, and $(b_n+1)^2/(n+1)\to0$. The raw Bartlett
  estimate converges in probability to $\sum_h\gamma_h$.}
\expandafter\def\csname statadic@proof@STATADIC-HAC-005\endcsname{There are $(N-h)_+$ available time pairs. Taking expectations
  gives the Bartlett taper times the finite-record edge factor. Dominated
  convergence proves the bias vanishes. The fourth-order dependence bound
  forces the estimator variance to zero. Chebyshev's inequality at half the
  requested tolerance and the triangle inequality finish the proof.}
\expandafter\def\csname statadic@assumptions@STATADIC-HAC-005\endcsname{\begin{itemize}
\item Second moments are stationary at each lag, with absolutely summable oriented lag moments.
\item Lag products are square-integrable and satisfy the uniform absolute covariance-sum bound.
\item Bandwidth b increases, and (b+1)\ensuremath{{}^2}/N tends to zero. For a covariance interpretation, the underlying scores are centered in population.
\end{itemize}}
\expandafter\def\csname statadic@diagnostics@STATADIC-HAC-005\endcsname{\begin{itemize}
\item Record the score definition, chronological order, common row count and bandwidth. Plain Newey–West uses raw products, so its covariance interpretation requires zero-mean scores. STATADIC uses the divisor N at every lag and weight 1-h/(b+1).
\item The default floor(4(N/100)\textasciicircum{}(2/9)) has the stated growth rate. A permanently fixed positive bandwidth targets a finite taper in general. Report any fixed override and justify the dependence/score assumptions independently.
\end{itemize}}
\expandafter\def\csname statadic@statement@STATADIC-HAC-006\endcsname{Suppose every observed scalar score satisfies $|z_t|\le M$.
  Let $H$ denote the raw Bartlett estimate with bandwidth $b$ and let $H_c$
  use $z_t-\bar z$. Then $|H_c-H|\le 6M(b+1)|\bar z|$.}
\expandafter\def\csname statadic@proof@STATADIC-HAC-006\endcsname{The sample mean has absolute value at most $M$. Expanding one
  centered product changes it by at most $3M|\bar z|$, and positive lags have
  two orientations. Each normalized lag contains at most $N$ products and the
  taper never exceeds one, so summation gives the bound.}
\expandafter\def\csname statadic@assumptions@STATADIC-HAC-006\endcsname{\begin{itemize}
\item N chronological rows of a scalar contrast are bounded in absolute value by K.
\item Centered and raw estimates use the same rows, same bandwidth and same fixed-N denominator.
\end{itemize}}
\expandafter\def\csname statadic@diagnostics@STATADIC-HAC-006\endcsname{\begin{itemize}
\item Record the centering scheme and compare the contrast sample mean and maximum score magnitude. A large raw-to-centered difference warrants checking misspecification or nonzero moment means, not switching silently between reference procedures.
\end{itemize}}
\expandafter\def\csname statadic@statement@STATADIC-HAC-007\endcsname{In addition to the raw Bartlett entry conditions, suppose
  $z_t$ is bounded by $M$, $E z_t=0$ and
  $\sum_{s<m}|\operatorname{Cov}(z_t,z_s)|\le D$ uniformly in $t,m$.
  The Bartlett estimate of $z_t-\bar z$ converges in probability to the
  same long-run variance as the raw estimate.}
\expandafter\def\csname statadic@proof@STATADIC-HAC-007\endcsname{Expanding the variance of $\bar z$ gives a bound $D/N$.
  The rate $(b+1)^2/N\to0$ therefore sends $(b+1)\bar z$ to zero in
  probability. The finite centering bound and the union bound transfer the
  raw-score limit to the sample-centered version.}
\expandafter\def\csname statadic@assumptions@STATADIC-HAC-007\endcsname{\begin{itemize}
\item Apply the theorem to true population-centered, uniformly bounded scores. Centering the original data by its sample mean gives the same centered scores.
\item Stationary summable lag moments and the growing-window fourth-order covariance bound hold; score covariance absolute row sums are also uniformly bounded.
\item The bandwidth grows with (b+1)\ensuremath{{}^2}/N tending to zero.
\end{itemize}}
\expandafter\def\csname statadic@diagnostics@STATADIC-HAC-007\endcsname{\begin{itemize}
\item Retain the actual moment matrix, centering flag, row order and bandwidth. Effective sample size, weak serial dependence and the correct sampling unit require a research-design argument.
\end{itemize}}
\expandafter\def\csname statadic@statement@STATADIC-HC-001\endcsname{Assume zero-mean, square-integrable errors with zero
  pair covariances and variances $v_j\le V$. For a projection row $a$ with
  $\sum_j a_j^2=h$, the removed fitted error $d=\sum_j a_j u_j$ satisfies
  $E d^2=\sum_j a_j^2 v_j\le Vh$. For a symmetric idempotent hat matrix
  the squared row norm is its diagonal leverage.}
\expandafter\def\csname statadic@proof@STATADIC-HC-001\endcsname{The fitted error has mean zero. Its second moment is its
  variance; expand that variance and use zero off-diagonal covariances.
  Bound each remaining variance and sum.}
\expandafter\def\csname statadic@assumptions@STATADIC-HC-001\endcsname{\begin{itemize}
\item Fixed (possibly whitened) design, uncorrelated mean-zero, square-integrable errors with variance bounded by V.
\item A fitted-error row has deterministic coefficients whose squared norm is its leverage h.
\end{itemize}}
\expandafter\def\csname statadic@diagnostics@STATADIC-HC-001\endcsname{\begin{itemize}
\item Inspect fit.leverage, including its maximum. A high-leverage observation can have a small fitted residual even with sizable error; residual magnitude alone cannot establish small sampling variance.
\end{itemize}}
\expandafter\def\csname statadic@statement@STATADIC-HC-002\endcsname{For squared influence weights $0\le w_i\le a$ of
  total mass at most $W$, assume $|u_i|\le K$, the squared errors have
  zero cross-covariances and variances at most $F$. Suppose the fitted
  adjustment $d_i$ has $E d_i^2\le q$. Then the expected absolute error
  of $\sum_i w_i(u_i-d_i)^2$ around $\sum_i w_iE u_i^2$ is at most
  $\sqrt{FaW}+W(2K\sqrt q+q)$.}
\expandafter\def\csname statadic@proof@STATADIC-HC-002\endcsname{The variance of the oracle weighted sum of squared errors
  is at most $F\sum w_i^2\le FaW$; Cauchy--Schwarz bounds its mean absolute
  deviation. For each row, expansion changes the square by at most
  $2K|d_i|+d_i^2$. Use $E|d_i|\le\sqrt{E d_i^2}$ and sum.}
\expandafter\def\csname statadic@assumptions@STATADIC-HC-002\endcsname{\begin{itemize}
\item For the sampling-variance interpretation, errors are mean zero and uncorrelated conditional on the fixed design. Squared contrast weights are bounded by a and have total at most W. Errors are bounded by K; squared errors have bounded variance F and zero cross-covariances.
\item The fitted prediction adjustment in each row has mean square bounded by q, derived from leverage for an exact linear projection.
\end{itemize}}
\expandafter\def\csname statadic@diagnostics@STATADIC-HC-002\endcsname{\begin{itemize}
\item Check rank and whether one observation dominates the squared influence weights for the contrast. State why true errors are independent (or have the required pair moments); this property is not determined by the residual plot.
\end{itemize}}
\expandafter\def\csname statadic@statement@STATADIC-HC-003\endcsname{Let $s_i\ge0$ be squared fitted influence contrasts and
  $0\le h_i\le H<1$. For $k=0,1,2$, the adjusted sum
  $A_k=\sum_i s_i/(1-h_i)^k$ lies between $A_0$ and $A_0/(1-H)^k$.
  HC2 and HC3 use $k=1,2$; HC1 uses $k=1$ with constant $h_i=p/N$.}
\expandafter\def\csname statadic@proof@STATADIC-HC-003\endcsname{Bound each inverse leverage factor between one and
  $(1-H)^{-k}$. Multiply by the nonnegative squared contribution and sum.
  If $h_i$ approaches one the upper bound can grow without limit.}
\expandafter\def\csname statadic@assumptions@STATADIC-HC-003\endcsname{\begin{itemize}
\item Squared fitted contrast influences are nonnegative. Every relevant hat value lies in [0,H] with H<1.
\item For HC1 use constant h=p/N and exponent 1, for HC2 the actual h and exponent 1, and for HC3 exponent 2.
\end{itemize}}
\expandafter\def\csname statadic@diagnostics@STATADIC-HC-003\endcsname{\begin{itemize}
\item Check the residual degrees of freedom for HC1 and every 1-h\_i for HC2/HC3. Report extra finite\_sample\_correction. Near-unit leverages amplify noise; validity checks at the boundary do not address that instability.
\end{itemize}}
\expandafter\def\csname statadic@statement@STATADIC-HC-004\endcsname{For row weights $0\le w_{ni}\le a_n\to0$ with bounded
  sum, suppose $|u_{ni}|\le K$, the squared errors have zero cross-covariance
  and bounded variance, and $E d_{ni}^2\le q_n\to0$ for the removed
  fitted errors. If $0\le h_{ni}\le H_n<1$ and $H_n\to0$, then every
  fixed $k=0,1,2$ variance estimate converges in mean absolute error to
  its sampling target $\tau_n=\sum_i w_{ni}E u_{ni}^2$.
  Hence its estimation error tends to zero in probability.}
\expandafter\def\csname statadic@proof@STATADIC-HC-004\endcsname{Diffuse weights concentrate the oracle squared-error sum,
  and the prediction-error bound makes residual replacement negligible.
  The finite leverage inequality places the corrected sum between HC0 and
  $(1-H_n)^{-k}$ times HC0. The multiplier tends to one and $\tau_n$
  remains bounded. Markov's inequality converts the mean absolute bound
  into convergence in probability.}
\expandafter\def\csname statadic@assumptions@STATADIC-HC-004\endcsname{\begin{itemize}
\item Fixed-design independent, mean-zero heteroskedastic errors satisfying the sufficient bounded-error, bounded fourth-variance and zero squared-error cross-covariance conditions.
\item Total squared contrast weight stays bounded, maximum squared weight tends to zero, and the uniform fitted prediction mean-square bound vanishes. The projection identity supplies that bound for OLS.
\item The maximum leverage tends to zero for HC2/3; residual degrees of freedom remain positive with p/N tending to zero for HC1.
\end{itemize}}
\expandafter\def\csname statadic@diagnostics@STATADIC-HC-004\endcsname{\begin{itemize}
\item Record the selected HC variant, leverage distribution, effective rank, residual degrees of freedom, contrast coefficient-map concentration and external scale multiplier. If independence is implausible, choose and justify a dependence model.
\end{itemize}}
\expandafter\def\csname statadic@statement@STATADIC-CLUSTER-001\endcsname{Fix the retained ordered row pairs $E$, e.g. all pairs
  sharing an entity or a time label. Put $Y_e=z_i z_j$ for $e=(i,j)$.
  If each pair product covaries with at most $d$ retained products and all
  these covariances have absolute value at most $C$, then
  $\operatorname{Var}(N^{-1}\sum_{e\in E}Y_e)\le C|E|d/N^2$.
  For zero-mean scores with covariance zero off $E$, its expectation equals
  $\operatorname{Var}(\sum_i z_i)/N$.}
\expandafter\def\csname statadic@proof@STATADIC-CLUSTER-001\endcsname{Off-mask covariances vanish, so the pair expectation recovers
  the full score-sum variance. Expand the sampling variance as a double
  covariance sum over pair products. Remove all off-neighborhood terms,
  count at most $d$ terms for each of $|E|$ choices, then divide by $N^2$.}
\expandafter\def\csname statadic@assumptions@STATADIC-CLUSTER-001\endcsname{\begin{itemize}
\item Mean-zero square-integrable scores, with zero covariance outside the retained ordered row-pair mask E. Pair products have finite second moments.
\item Within the retained pairs, a given pair product covaries with at most d others; all pair-product covariances are bounded in magnitude by C.
\end{itemize}}
\expandafter\def\csname statadic@diagnostics@STATADIC-CLUSTER-001\endcsname{\begin{itemize}
\item Inspect occupied cluster counts, sizes and two-way intersection multiplicities, including after transformations or missing-row filtering. An apparently sparse label graph does not establish independence of shocks across labels.
\end{itemize}}
\expandafter\def\csname statadic@statement@STATADIC-CLUSTER-002\endcsname{Under the cluster moment conditions, if
  $|E_n|d_n/N^2\to0$ and the pair-covariance bound is uniform, the
  difference between the raw cluster meat and the true variance of the
  normalized score sum tends to zero in probability. The pair mask may
  change with $N$ and can include arbitrary within-cluster covariance.}
\expandafter\def\csname statadic@proof@STATADIC-CLUSTER-002\endcsname{The raw cluster meat is unbiased for the finite-record
  variance target. Chebyshev's inequality bounds its error probability at
  tolerance $t>0$ by $C|E_n|d_n/(N^2t^2)$, which tends to zero.}
\expandafter\def\csname statadic@assumptions@STATADIC-CLUSTER-002\endcsname{\begin{itemize}
\item The population exclusions and bounded pair-product covariance conditions in the finite pair-mask theorem hold at each N.
\item The pair-count times maximal product-dependence degree divided by N\ensuremath{{}^2} tends to zero.
\end{itemize}}
\expandafter\def\csname statadic@diagnostics@STATADIC-CLUSTER-002\endcsname{\begin{itemize}
\item Report cluster sizes and counts on each active dimension, and justify the product-dependence neighborhood. Record group\_debias and slope/effect degree-of-freedom choices; consistency at the unadjusted target requires asymptotically neutral multipliers.
\end{itemize}}
\expandafter\def\csname statadic@statement@STATADIC-PANELCOV-003\endcsname{Let $T_{ai}=1$ if the label of observed row $i$ is $a$.
  For any score matrix $U$ and label kernel $K$,
  $(TU)^TK(TU)=U^T\widetilde KU$, where
  $\widetilde K_{ij}=K_{\ell(i),\ell(j)}$.
  For calendar-time pooling this retains cross-entity pairs; using
  entity-time labels with a block-diagonal kernel suppresses them.}
\expandafter\def\csname statadic@proof@STATADIC-PANELCOV-003\endcsname{Compute each entry of $T^TKT$. Each membership column selects
  exactly one label on either side of the kernel. Then associate the matrix
  products. Missing entity-time cells contribute no score.}
\expandafter\def\csname statadic@assumptions@STATADIC-PANELCOV-003\endcsname{\begin{itemize}
\item Finite actual retained row labels and a declared kernel on those labels; all kernel products use the same transformed score rows.
\end{itemize}}
\expandafter\def\csname statadic@diagnostics@STATADIC-PANELCOV-003\endcsname{\begin{itemize}
\item Record the actual calendar IDs, selected kernel, bandwidth and whether cross-entity time pairs are retained. Between estimation has no within-entity time covariance; first differences use transformed rows.
\end{itemize}}
\expandafter\def\csname statadic@statement@STATADIC-COVBRIDGE-001\endcsname{Let $E$ contain the retained ordered observation pairs,
  with $|k_{ij}|\le1$, and let the true score contrast satisfy $|z_i|\le K$.
  Suppose each fitted-score adjustment has $E d_i^2\le q$. Replacing $z_i$
  by $z_i-d_i$ in the normalized weighted pair sum changes it by an expected
  absolute amount at most $(|E|/N)(2K\sqrt q+q)$.}
\expandafter\def\csname statadic@proof@STATADIC-COVBRIDGE-001\endcsname{Expand each pair product. The two cross terms are bounded by
  $K(|d_i|+|d_j|)$ and $|d_i d_j|\le(d_i^2+d_j^2)/2$. Take expectations
  and use Cauchy--Schwarz for each $E|d_i|$. Triangle inequalities and
  $|k_{ij}|\le1$ allow summing these bounds over retained pairs.}
\expandafter\def\csname statadic@assumptions@STATADIC-COVBRIDGE-001\endcsname{\begin{itemize}
\item A fixed set E of retained ordered pairs with deterministic weights bounded in modulus by one.
\item True scores are bounded by K and square-integrable. Replacing any row with its fitted counterpart has mean-square error at most q.
\end{itemize}}
\expandafter\def\csname statadic@diagnostics@STATADIC-COVBRIDGE-001\endcsname{\begin{itemize}
\item Compute the size of the retained mask relative to N and check how aggressively it grows with bandwidth or cluster size. Use an estimator-specific rate for fitted-score error and retain the same Jacobian/normalization when constructing the final sandwich.
\end{itemize}}
\expandafter\def\csname statadic@statement@STATADIC-MEDIAN-004\endcsname{For a nonempty finite unweighted sample in a proper real inner
  product space, if the observations are not all on one affine line, their
  sum of Euclidean distances has a unique minimizer.}
\expandafter\def\csname statadic@proof@STATADIC-MEDIAN-004\endcsname{Existence follows by minimizing on a compact ball. Suppose two
  different points $a,b$ minimize. At their midpoint the triangle inequality
  bounds each distance by its endpoint average. At least one inequality is strict:
  equality for every observation would put $a-x_i$ and $b-x_i$ on the same ray
  for every $i$, hence each $x_i$ on the affine line through $a,b$. Summing gives
  a value below the common minimum, a contradiction.}
\expandafter\def\csname statadic@assumptions@STATADIC-MEDIAN-004\endcsname{\begin{itemize}
\item Nonempty finite unweighted sample in a proper real inner-product space, including finite-dimensional Euclidean space.
\item The sample is not contained in one affine line.
\end{itemize}}
\expandafter\def\csname statadic@diagnostics@STATADIC-MEDIAN-004\endcsname{\begin{itemize}
\item Check the affine span after removing coincident observations: for Euclidean coordinates, center the rows and inspect whether the second singular value is meaningfully positive on the data scale.
\item The group convergence flag tests numerical stationarity. On or near collinearity, retain that design information and avoid inferring uniqueness or stability from the flag.
\end{itemize}}
\expandafter\def\csname statadic@statement@STATADIC-MEDIAN-005\endcsname{For the two observations $-1$ and $1$ on the real line,
  the complete set of minimizers of $|t+1|+|t-1|$ is $[-1,1]$.}
\expandafter\def\csname statadic@proof@STATADIC-MEDIAN-005\endcsname{On the interval the two absolute deviations sum to $2$.
  Below $-1$ they sum to $-2t>2$, and above $1$ to $2t>2$.
  A stationarity certificate therefore need not imply uniqueness on collinear data.}
\expandafter\def\csname statadic@assumptions@STATADIC-MEDIAN-005\endcsname{\begin{itemize}
\item Two equally weighted observations at \ensuremath{-}1 and 1 on the real line.
\end{itemize}}
\expandafter\def\csname statadic@diagnostics@STATADIC-MEDIAN-005\endcsname{\begin{itemize}
\item For tied or collinear inputs, compare objective values and report any solver tie convention. A converged center can be one among many minimizers.
\end{itemize}}
\expandafter\def\csname statadic@statement@STATADIC-TRANSPORT-002\endcsname{For a finite nonempty square cost matrix $C$ there exist a permutation
  $\sigma$ and potentials $u,v$ with $u_i+v_j\le C_{ij}$. The assignment minimizes cost
  among all doubly stochastic couplings and attains $\sum_i u_i+\sum_j v_j$.
  Dividing by the common number of rows gives the uniform probability-mass objective.}
\expandafter\def\csname statadic@proof@STATADIC-TRANSPORT-002\endcsname{Minimize over the finite permutations and relabel the columns by an
  optimum. Subtract the diagonal cost from each row to form residual costs $D$.
  Every nonnegative circulation has nonnegative residual cost: pad its diagonal to
  equal row and column totals, normalize, and decompose the resulting doubly stochastic
  matrix into permutation matrices. From a fixed origin let $v_j$ be the infimum
  cost of sending one unit to $j$. A direct edge makes this set nonempty; adding
  a return edge bounds it below. Extending a flow by one edge gives
  $v_j-v_i\le D_{ij}$. Undo the relabeling and row shift; the matched edges are tight.
  The resulting feasible dual lower bound is the assignment value.}
\expandafter\def\csname statadic@assumptions@STATADIC-TRANSPORT-002\endcsname{\begin{itemize}
\item Finite nonempty equal-size supports and uniform mass. Costs are finite real numbers; use nonnegative powered ground costs for the Wasserstein interpretation.
\end{itemize}}
\expandafter\def\csname statadic@diagnostics@STATADIC-TRANSPORT-002\endcsname{\begin{itemize}
\item Inspect input\_valid, finite and power\_underflow. These flags cannot certify a chosen permutation or a dual gap.
\item For an independent witness solve on the same powered cost and masses. Verify reduced costs are nonnegative, complementary slackness on transported mass, and primal versus dual objective before rooting.
\end{itemize}}
\expandafter\def\csname statadic@statement@STATADIC-TRANSPORT-003\endcsname{For sorted real observations $x_i,y_i$ on the same finite nonempty
  rank axis, matching equal ranks minimizes $\sum P_{ij}(x_i-y_j)^2$ over all
  doubly stochastic $P$. Dividing by the sample size gives squared uniform $W_2$.}
\expandafter\def\csname statadic@proof@STATADIC-TRANSPORT-003\endcsname{In each permutation cost the two squared marginal sums are
  invariant. The remaining cross product is maximized by the rank matching via
  the rearrangement inequality. Birkhoff decomposition extends that optimum
  from permutations to every doubly stochastic matrix.}
\expandafter\def\csname statadic@assumptions@STATADIC-TRANSPORT-003\endcsname{\begin{itemize}
\item Real observations sorted in a common order, equal atom counts and uniform weights.
\item Quadratic ground cost and unit row and column totals before normalization.
\end{itemize}}
\expandafter\def\csname statadic@diagnostics@STATADIC-TRANSPORT-003\endcsname{\begin{itemize}
\item Verify equal normalized row masses before comparing with sorted index matching. For unequal masses or support sizes use the quantile integration path and compare against a weighted transportation LP.
\end{itemize}}
\expandafter\def\csname statadic@statement@STATADIC-TRANSPORT-004\endcsname{For two uniform equal-size samples in a real inner product space,
  let the directions $\theta_l$ have norm at most one and weights $q_l\ge0$
  summing to one. At $p\ge1$, their weighted mean of optimal one-dimensional
  $p$-power costs is at most the full Euclidean $p$-power assignment cost.}
\expandafter\def\csname statadic@proof@STATADIC-TRANSPORT-004\endcsname{Use the same full-space matching as a candidate for every
  projected problem. Cauchy--Schwarz bounds each projected paired distance by
  the ambient distance. Raise to the power $p$, average over rows, and then
  over directions. Optimal projection costs can only be smaller.}
\expandafter\def\csname statadic@assumptions@STATADIC-TRANSPORT-004\endcsname{\begin{itemize}
\item Equal-sized finite uniform samples in a real inner-product space; normalized nonnegative directional weights, direction norms at most one, and real p at least one.
\item Each supplied projected cost is bounded above by the cost of each projected permutation, as for a projected optimum.
\end{itemize}}
\expandafter\def\csname statadic@diagnostics@STATADIC-TRANSPORT-004\endcsname{\begin{itemize}
\item Retain the same ProjectionDesign, finite direction count, seed or directions, ground coordinate preprocessing and order across compared pairs.
\item A small sliced value gives a lower bound on full transport and does not cap its error. A finite design can miss differences perpendicular to all its directions.
\end{itemize}}
\expandafter\def\csname statadic@statement@STATADIC-TRANSPORT-005\endcsname{On finite nonempty supports let $a_i>0$, $b_j\ge0$ be
  two unit-mass vectors. For every finite real cost array $C$, there exist
  a coupling $P$ and potentials $u_i+v_j\le C_{ij}$ with
  $\langle P,C\rangle=a^Tu+b^Tv$. Thus the primal minimum and dual maximum
  agree and both are attained. Zero source masses can be omitted from the support.}
\expandafter\def\csname statadic@proof@STATADIC-TRANSPORT-005\endcsname{The feasible set is closed in a compact cube and contains the
  independent coupling, so choose a minimizer. For each source row $i$, define
  the cost of sending mass to row $k$ as the smallest $C_{kj}-C_{ij}$ among
  its occupied cells. Every permutation of rows has nonnegative total such cost:
  move a common small mass from the corresponding occupied cells according
  to the permutation, which preserves both marginals. The shortest-flow
  potential argument gives row potentials bounded by these differences.
  Set $v_j=\min_i(C_{ij}-u_i)$. Every occupied cell is tight, and the
  reduced-cost identity gives primal-dual equality.}
\expandafter\def\csname statadic@assumptions@STATADIC-TRANSPORT-005\endcsname{\begin{itemize}
\item Finite nonempty rectangular atom rosters. Source masses are positive, target masses nonnegative, and both sum to one. Remove zero-mass source padding first.
\item Finite real cost array. For W\_p apply the result to the prescribed p-th powers and normalize before rooting.
\end{itemize}}
\expandafter\def\csname statadic@diagnostics@STATADIC-TRANSPORT-005\endcsname{\begin{itemize}
\item Check both empirical mass vectors after normalization and mask padding. Compare the native powered result to a rectangular LP on the same costs before taking roots.
\item The scalar numerical result does not expose optimizing potentials. To audit dual feasibility and complementary slackness retain independent LP primal and dual arrays.
\end{itemize}}
\expandafter\def\csname statadic@statement@STATADIC-ENERGY-006\endcsname{Suppose a symmetric hollow ground $d$ has negative type on every
  finite set. If $d$ is integrable under $P\otimes P$, $Q\otimes Q$ and both mixed
  products, write $D_{AB}=\int d(x,y)\,dA(x)\,dB(y)$. Then
  $2D_{PQ}-D_{PP}-D_{QQ}\ge0$ for arbitrary probability laws $P,Q$.}
\expandafter\def\csname statadic@proof@STATADIC-ENERGY-006\endcsname{Take independent samples of size $n$ from each law. Apply finite
  negative type to the pooled sample with $n$ weights $+1$ and $n$ weights $-1$.
  Integrate the finite inequality. Within each sample there are $n(n-1)$
  off-diagonal terms; diagonal costs vanish. Each mixed block contributes
  $n^2D_{PQ}$ by independence and symmetry. Divide by $n$ and let $n$ grow:
  a strictly negative population energy would contradict the resulting bound
  for a sufficiently large finite $n$.}
\expandafter\def\csname statadic@assumptions@STATADIC-ENERGY-006\endcsname{\begin{itemize}
\item The ground is symmetric, hollow and conditionally negative on every finite signed zero-mass configuration.
\item Both inputs are probability measures, and the ground cost is integrable under each within- and cross-law product.
\end{itemize}}
\expandafter\def\csname statadic@diagnostics@STATADIC-ENERGY-006\endcsname{\begin{itemize}
\item Specify the ground and justify its negative type on the scientific domain. A finite Gram or distance check can detect a violation on the observed support but cannot establish the universal ground law.
\item Verify moment/tail assumptions for the target laws. A finite empirical value alone cannot establish integrability of the population costs.
\end{itemize}}
\expandafter\def\csname statadic@statement@STATADIC-ENERGY-007\endcsname{On a second-countable Borel pseudometric space of negative type,
  probability laws $P,Q$ with integrable distances to one common origin have nonnegative
  population energy. This includes discrete, continuous and mixed laws.}
\expandafter\def\csname statadic@proof@STATADIC-ENERGY-007\endcsname{The triangle inequality bounds the pair distance by the sum of the
  two distances to the origin. Each is integrable on the product probability
  space, so all four pair integrals exist. Apply the finite-sample averaging
  theorem with the original finite-negative-type inequality.}
\expandafter\def\csname statadic@assumptions@STATADIC-ENERGY-007\endcsname{\begin{itemize}
\item A second-countable Borel pseudometric space of negative type.
\item The two laws have finite expected distance to a common origin.
\end{itemize}}
\expandafter\def\csname statadic@diagnostics@STATADIC-ENERGY-007\endcsname{\begin{itemize}
\item Record the distance, domain and sampling population, and support the finite first-moment condition. Bounded ground distances suffice; for unbounded data explain the tail assumptions.
\end{itemize}}
\expandafter\def\csname statadic@statement@STATADIC-ENERGY-008\endcsname{Let $k$ be symmetric with positive-semidefinite Gram matrix
  on every finite tuple. For $d(x,y)=k(x,x)+k(y,y)-2k(x,y)$ suppose the
  four population pair costs are integrable. The energy generated by $d$ is nonnegative.}
\expandafter\def\csname statadic@proof@STATADIC-ENERGY-008\endcsname{For finite signed weights summing to zero the diagonal
  kernel terms cancel. The signed distance quadratic equals minus twice
  the kernel quadratic, which is nonpositive by Gram positivity.
  The generated ground is hollow and symmetric, so the population
  negative-type theorem applies.}
\expandafter\def\csname statadic@assumptions@STATADIC-ENERGY-008\endcsname{\begin{itemize}
\item A real symmetric kernel is PSD on every finite tuple. The ground is k(x,x)+k(y,y)\ensuremath{-}2k(x,y).
\item The resulting dissimilarity is integrable under each marginal product.
\end{itemize}}
\expandafter\def\csname statadic@diagnostics@STATADIC-ENERGY-008\endcsname{\begin{itemize}
\item Verify the exact relation between kernel and ground, including scale and sign. Compare a numerical Gram spectrum only after checking symmetry and stating its tolerance.
\item A PSD Gram on one sample is a finite diagnostic, not a proof that a custom kernel is positive on the entire population domain.
\end{itemize}}
\expandafter\def\csname statadic@statement@STATADIC-BARYCENTRE-002\endcsname{For equal-cardinality uniform finite clouds in Euclidean space
  with fixed nonnegative source weights summing to one, the equal-cardinality
  free-support $W_2^2$ barycentre problem has a global minimizer. One minimizer
  is the rowwise weighted centroid of a choice of source permutations: its joint
  squared cost is no larger than the cost of any support and any matching pattern.}
\expandafter\def\csname statadic@proof@STATADIC-BARYCENTRE-002\endcsname{For each tuple of permutations, compute the matched row centroids.
  Their conditional cost is minimal by the exact squared-distance excess formula.
  The set of tuples of permutations is finite and nonempty. Minimize the conditional
  minima over that set; the selected centroids then beat every competing support
  and permutation tuple.}
\expandafter\def\csname statadic@assumptions@STATADIC-BARYCENTRE-002\endcsname{\begin{itemize}
\item Finite equal-cardinality uniform Euclidean source clouds and candidate supports. Source weights are nonnegative and sum to one.
\item Squared Euclidean ground cost before taking any root.
\end{itemize}}
\expandafter\def\csname statadic@diagnostics@STATADIC-BARYCENTRE-002\endcsname{\begin{itemize}
\item Record coordinate preprocessing, metric, p, source weights and initialization. On small examples compare with exhaustive permutation tuples after fixing one relabeling.
\item Stable assignments certify a coordinatewise fixed point. The reported initial-to-final objective decrease is not an upper bound on the remaining global error.
\end{itemize}}
\expandafter\def\csname statadic@statement@STATADIC-BARYCENTRE-003\endcsname{On the candidate probability simplex, let
  $F(s)=t+s^TQs/2+c^Ts$, where $Q$ is symmetric positive semidefinite on
  zero-sum directions. The empirical squared-MMD mixture objective has this
  form with $Q=2G$. A minimizing mixture exists. At any feasible mixture $s$
  with gradient $g=Qs+c$, $0\le F(s)-\min F\le g^Ts-\min_i g_i$.
  A zero Frank--Wolfe gap therefore certifies a global optimum.}
\expandafter\def\csname statadic@proof@STATADIC-BARYCENTRE-003\endcsname{The simplex is a nonempty closed subset of a compact coordinate
  cube and the quadratic is continuous. Every feasible mixture averages its
  gradient coordinates and so has linearized cost at least their minimum.
  Apply the tangent-space quadratic gap bound at any attained minimizer.}
\expandafter\def\csname statadic@assumptions@STATADIC-BARYCENTRE-003\endcsname{\begin{itemize}
\item A nonempty finite roster of candidate laws and variable nonnegative mixture weights summing to one.
\item The empirical squared-MMD or energy objective is the declared quadratic plus a target constant, with symmetric Hessian PSD on all zero-sum directions.
\end{itemize}}
\expandafter\def\csname statadic@diagnostics@STATADIC-BARYCENTRE-003\endcsname{\begin{itemize}
\item Retain simplex\_sum\_error, min\_weight, frank\_wolfe\_gap, projected\_gradient\_norm and the reported tangent curvature. For the energy route verify the negative-type sign before interpreting the gap.
\item Converged compares the computed gap against a scaled numerical tolerance. Degenerate curvature may allow multiple optimal weight vectors.
\end{itemize}}
\expandafter\def\csname statadic@statement@STATADIC-BARYCENTRE-004\endcsname{In any metric space, if $m$ is at equal half-distance from
  endpoints $a,b$, it globally minimizes the equal-weight sum of $p$-powered
  distances for every real $p\ge1$.}
\expandafter\def\csname statadic@proof@STATADIC-BARYCENTRE-004\endcsname{For any candidate $z$, the average of its two endpoint distances
  bounds $d(a,b)/2$ by the triangle inequality. Apply the increasing $p$th power
  and then convexity to bound this by the average of the two powers.
  The midpoint attains both inequalities.}
\expandafter\def\csname statadic@assumptions@STATADIC-BARYCENTRE-004\endcsname{\begin{itemize}
\item Real order p at least one and exactly equal source weights.
\item The candidate is at half the endpoint distance from each source in the declared pseudometric.
\end{itemize}}
\expandafter\def\csname statadic@diagnostics@STATADIC-BARYCENTRE-004\endcsname{\begin{itemize}
\item Check the order and source weights. The runtime uses a tolerance for equality, whereas the theorem requires exact equality.
\item Linear interpolation must realize a geodesic midpoint for the chosen ground. A straight coordinate average need not do so for normalized-angular data.
\end{itemize}}
\expandafter\def\csname statadic@statement@STATADIC-BARYCENTRE-005\endcsname{Suppose $\pi$ minimizes squared assignment cost between two equal-size
  uniform Euclidean clouds $a,b$. With source weights $1-t,t$ in $[0,1]$, the support
  $z_i=(1-t)a_i+t b_{\pi(i)}$ achieves the global minimum: the weighted squared cost
  at every competing support and pair of matchings is at least
  $t(1-t)\sum_i\|a_i-b_{\pi(i)}\|^2$, attained at the stated support.}
\expandafter\def\csname statadic@proof@STATADIC-BARYCENTRE-005\endcsname{Complete the square for each pair of matched atoms. The remainder
  is a nonnegative squared deviation from the paired centroid. Relabel the
  competitor's two matchings to obtain a single permutation between the sources;
  its source-to-source cost is no smaller than the optimal one. The interpolated
  support makes all squared deviations vanish.}
\expandafter\def\csname statadic@assumptions@STATADIC-BARYCENTRE-005\endcsname{\begin{itemize}
\item Two Euclidean clouds of equal nonzero cardinality and uniform atom masses. Source weights 1\ensuremath{-}t and t are nonnegative and sum to one.
\item The supplied permutation minimizes the squared source-to-source cost.
\end{itemize}}
\expandafter\def\csname statadic@diagnostics@STATADIC-BARYCENTRE-005\endcsname{\begin{itemize}
\item Verify assignments against the powered cost or an independent LP, especially if ties or extreme cost scales occur. Compare the achieved mean objective to t(1\ensuremath{-}t) times the source squared-Wasserstein cost.
\end{itemize}}
\expandafter\def\csname statadic@statement@STATADIC-BARYCENTRE-006\endcsname{For two uniform Euclidean clouds with equal counts and exactly equal
  source weights, match atoms with a minimum $p$-powered assignment, $p\ge1$. The midpoints
  of those pairs globally minimize the equal-weight two-source powered objective over
  all competing equal-count supports and their assignments.}
\expandafter\def\csname statadic@proof@STATADIC-BARYCENTRE-006\endcsname{The average powered distance to a pair is minimized at its Euclidean
  midpoint, at value $2^{-p}d(a,b)^p$. Sum that inequality over competing pairs.
  Their paired endpoint cost is at least the optimum after relabeling one assignment.
  The proposed midpoint support attains the resulting lower bound.}
\expandafter\def\csname statadic@assumptions@STATADIC-BARYCENTRE-006\endcsname{\begin{itemize}
\item Equal source weights, equal-size uniform Euclidean clouds and finite real order p at least one.
\item An optimal source-to-source permutation for distance to the p-th power.
\end{itemize}}
\expandafter\def\csname statadic@diagnostics@STATADIC-BARYCENTRE-006\endcsname{\begin{itemize}
\item Match using d raised to the declared order. Verify the returned objective against the matched source powered cost divided by 2 to that order.
\item At p different from two, the API applies a tolerance for equal weights. Treat this exact certificate as applying to actual equal weights.
\end{itemize}}
\expandafter\def\csname statadic@statement@STATADIC-BLOCKLENGTH-002\endcsname{For $N,D>0$ and $G\ne 0$, the displayed positive
  length $b_*=(2G^2/D)^{1/3}N^{1/3}$ uniquely minimizes
  $G^2/b^2+Db/N$ over all positive real lengths.}
\expandafter\def\csname statadic@proof@STATADIC-BLOCKLENGTH-002\endcsname{Cube the two positive factors of $b_*$ and simplify.
  The resulting balance $Db_*^3=2NG^2$ gives uniqueness by the
  factored risk-gap identity.}
\expandafter\def\csname statadic@assumptions@STATADIC-BLOCKLENGTH-002\endcsname{\begin{itemize}
\item The leading AMSE is G\ensuremath{{}^2}/b\ensuremath{{}^2} + Db/N, with D and N positive and G nonzero. Candidate lengths are continuous and positive.
\end{itemize}}
\expandafter\def\csname statadic@diagnostics@STATADIC-BLOCKLENGTH-002\endcsname{\begin{itemize}
\item Choose the resampling scheme for the statistic. Record the zero-frequency and lag-weighted population quantities needed for that scheme, and the subsequent integer policy for circular resampling.
\end{itemize}}
\expandafter\def\csname statadic@statement@STATADIC-BLOCKLENGTH-003\endcsname{Suppose the joint lag pilots $(\hat G_N,\hat S_N)$
  converge in probability to $(G,S)$ with $G\ne0$, $D=cS^2>0$.
  Use the denominator floor $\epsilon>0$ and length floor one. The ratio
  of selected length to the interior oracle converges in probability
  to $(2G^2/\max(D,\epsilon))^{1/3}/(2G^2/D)^{1/3}$. In particular it
  converges to one for $D\ge\epsilon$.}
\expandafter\def\csname statadic@proof@STATADIC-BLOCKLENGTH-003\endcsname{Divide the selected length by $N^{1/3}$; its lower
  bound becomes $N^{-1/3}\to0$. The floored denominator is positive
  and the cube-root scale is continuous there. On almost surely
  convergent subsubsequences of the pilots apply continuity and divide
  by the positive oracle scale.}
\expandafter\def\csname statadic@assumptions@STATADIC-BLOCKLENGTH-003\endcsname{\begin{itemize}
\item The lag-weighted and zero-frequency pilots converge jointly in probability to G and S. They are measurable.
\item The population G is nonzero, D=cS\ensuremath{{}^2} is positive, and the denominator floor \ensuremath{\varepsilon} is positive.
\end{itemize}}
\expandafter\def\csname statadic@diagnostics@STATADIC-BLOCKLENGTH-003\endcsname{\begin{itemize}
\item Record the denominator before its 10\ensuremath{{}^{-8}} floor and examine the size of G. If D is below the floor even in the population, the oracle ratio changes. Report the retained lag cap and assess how the result changes with bandwidth.
\end{itemize}}
\expandafter\def\csname statadic@statement@STATADIC-BLOCKLENGTH-004\endcsname{Fix a finite set of retained signed lags. Suppose that
  each empirical autocorrelation converges to its population value and the
  taper stays between zero and one. On lags with nonzero population
  autocorrelation assume the taper converges to one. Both tapered pilot sums,
  with multipliers $|k|$ and $1$, converge to their untapered population
  sums. These are the full pilots when the retained lags contain the
  population support.}
\expandafter\def\csname statadic@proof@STATADIC-BLOCKLENGTH-004\endcsname{At nonzero population lags multiply the two limits.
  At zero population lags bound the absolute tapered estimate by the
  absolute sample autocorrelation; the taper may continue to fluctuate.
  Add the limits over the finite retained lag set.}
\expandafter\def\csname statadic@assumptions@STATADIC-BLOCKLENGTH-004\endcsname{\begin{itemize}
\item A fixed finite retained signed-lag set and convergent empirical autocorrelations for each retained lag.
\item Taper weights stay in [0,1] and tend to one on each retained lag with nonzero population autocorrelation. To obtain the full population sums, the set must contain the population support.
\end{itemize}}
\expandafter\def\csname statadic@diagnostics@STATADIC-BLOCKLENGTH-004\endcsname{\begin{itemize}
\item Inspect autocorrelation estimates outside the selected flat part. A cutoff preceding later nonzero lags, or an unduly small user cap, invalidates the support condition. The native search uses a five-lag window and a 2/\ensuremath{\sqrt{N}} threshold.
\end{itemize}}
\expandafter\def\csname statadic@statement@STATADIC-BLOCKLENGTH-005\endcsname{For zero lag-weighted pilot $\hat G=0$, the selector
  returns one for every sample size. In particular, a search restricted
  to lag zero can never yield an interior $N^{1/3}$ length.}
\expandafter\def\csname statadic@proof@STATADIC-BLOCKLENGTH-005\endcsname{At lag zero the multiplier $|k|$ vanishes. The cube-root
  term is zero, and taking the maximum with one gives the stated result.}
\expandafter\def\csname statadic@assumptions@STATADIC-BLOCKLENGTH-005\endcsname{\begin{itemize}
\item The real-arithmetic implementation expression is evaluated at \ensuremath{\hat G}=0. A retained grid consisting only of lag zero satisfies this condition.
\end{itemize}}
\expandafter\def\csname statadic@diagnostics@STATADIC-BLOCKLENGTH-005\endcsname{\begin{itemize}
\item Record max\_lag, whether the series is constant, and the unfloored length before interpreting a reported length of one.
\end{itemize}}
\expandafter\def\csname statadic@statement@STATADIC-PANEL-002\endcsname{On a finite roster with positive row weights, suppose
  $x_i-r_i=u_{e_i}+v_{t_i}$ and the weighted residual totals vanish within
  every entity and every period. For every competing additive fit $a,b$,
  its weighted squared error exceeds that of $r$ by exactly
  $\sum_i w_i(u_{e_i}+v_{t_i}-a_{e_i}-b_{t_i})^2$. The competitor attains
  the minimum exactly when it gives the same fitted values on every row.}
\expandafter\def\csname statadic@proof@STATADIC-PANEL-002\endcsname{Group the cross term first by entity, then by period. Both
  parts vanish by the weighted normal equations. The resulting excess is a
  sum of squares with positive weights, which vanishes exactly when every
  fitted row agrees. This argument permits missing cells and disconnected
  components of the entity--period graph.}
\expandafter\def\csname statadic@assumptions@STATADIC-PANEL-002\endcsname{\begin{itemize}
\item Finite observed panel with aligned entity and time labels and positive row weights. The argument is per response or design column.
\item The removed component is additive in entity and time. Weighted residual sums vanish for each entity and each period.
\end{itemize}}
\expandafter\def\csname statadic@diagnostics@STATADIC-PANEL-002\endcsname{\begin{itemize}
\item Use identical retained rows, labels and relative precisions for the sweep and reference design. Check weighted residual sums within both margins and whether the removed component fits the joint dummy design.
\item The independent test test\_weighted\_unbalanced\_two\_way\_absorption\_matches\_dummy\_projection compares the residual to a NumPy least-squares fit of the square-root-weighted dummy matrix.
\end{itemize}}
\expandafter\def\csname statadic@statement@STATADIC-PANEL-003\endcsname{On a rectangular panel with weights $w_{et}=p_eq_t$, where each
  marginal vector is normalized, entity and time demeaning commute. Both orders
  give $x_{et}-\bar x_{e\cdot}-\bar x_{\cdot t}+\bar x_{\cdot\cdot}$.
  Thus one sweep in each direction removes both sets of weighted group means.}
\expandafter\def\csname statadic@proof@STATADIC-PANEL-003\endcsname{Linearity moves a group average through a subtraction.
  The two orders of the grand weighted average agree by interchange of finite
  sums. Normalization preserves constants, so centering a group twice does
  not change it.}
\expandafter\def\csname statadic@assumptions@STATADIC-PANEL-003\endcsname{\begin{itemize}
\item Complete entity by time roster. For the weighted projection interpretation, row weights factor as positive normalized entity and time weights.
\end{itemize}}
\expandafter\def\csname statadic@diagnostics@STATADIC-PANEL-003\endcsname{\begin{itemize}
\item Check that every intended entity-period pair is present exactly once. If weights vary, verify their array factors into entity and period vectors before using a one-pass formula.
\end{itemize}}
\expandafter\def\csname statadic@statement@STATADIC-PANEL-004\endcsname{For positive row weights and occupied groups, suppose a full
  entity-then-time sweep leaves an array $x$ exactly unchanged. Its entity and
  period weighted totals both vanish. Thus checking an exact fixed point, in
  addition to the additive form of the removed effects, suffices for the
  weighted two-way projection certificate.}
\expandafter\def\csname statadic@proof@STATADIC-PANEL-004\endcsname{Write $y=x-A$ after the entity step and $x=y-B$ after the
  time step, so $B=-A$ rowwise. The intermediate $y$ is orthogonal to every
  entity effect and the final $x$ to every time effect. Hence the sum of
  $w_i A_i^2=w_i(x_i-y_i)A_i$ vanishes. Positivity forces $A_i=0$ in
  every row, so $y=x$ and the entity equations also hold.}
\expandafter\def\csname statadic@assumptions@STATADIC-PANEL-004\endcsname{\begin{itemize}
\item Finite observed panel with aligned entity and time labels and positive row weights. The argument is per response or design column.
\item Every label in the declared roster has positive total weight. An entity-then-time sweep gives exactly its input.
\end{itemize}}
\expandafter\def\csname statadic@diagnostics@STATADIC-PANEL-004\endcsname{\begin{itemize}
\item PanelProjection.error measures maximum change after the complete sweep. Recompute both margins directly when checking an approximate answer and retain scale, tolerance and iteration count.
\end{itemize}}
\expandafter\def\csname statadic@statement@STATADIC-PANEL-005\endcsname{With unit weights on the three observed cells
  $(0,0),(0,1),(1,0)$ and responses $(0,2,0)$, the entity-then-time
  residual in the first cell is $-1/2$, whereas time-then-entity gives zero.}
\expandafter\def\csname statadic@proof@STATADIC-PANEL-005\endcsname{The entity-first intermediate is $(-1,1,0)$; the time sweep
  gives $(-1/2,0,1/2)$. In the other order the first time sweep already
  removes the entire response, which depends only on time.}
\expandafter\def\csname statadic@assumptions@STATADIC-PANEL-005\endcsname{\begin{itemize}
\item Observe (0,0), (0,1) and (1,0) with unit weights and responses 0, 2 and 0.
\end{itemize}}
\expandafter\def\csname statadic@diagnostics@STATADIC-PANEL-005\endcsname{\begin{itemize}
\item On an incomplete roster compare the first-pass residual to the dense joint dummy projection. Positive row weights and at least one observation per label alone cannot justify one pass.
\end{itemize}}
\expandafter\def\csname statadic@statement@STATADIC-PANEL-006\endcsname{Starting from the observed column, repeat entity-then-time
  weighted demeaning any finite number of times. If the resulting $r$ is an
  exact fixed point, it gives the unique least-squares residual on the roster:
  for every additive fit $a_e+b_t$, its weighted loss equals the loss of $r$
  plus $\sum_i w_i[x_i-r_i-a_{e_i}-b_{t_i}]^2$.}
\expandafter\def\csname statadic@proof@STATADIC-PANEL-006\endcsname{By induction, every sweep subtracts only an entity vector and
  a time vector, so $x-r$ belongs to the additive design space. The exact
  fixed-point theorem gives both sets of normal equations. The weighted
  projection certificate now gives the squared-error identity.}
\expandafter\def\csname statadic@assumptions@STATADIC-PANEL-006\endcsname{\begin{itemize}
\item Finite observed panel with aligned entity and time labels and positive row weights. The argument is per response or design column.
\item Groups are occupied, the sweep starts at the input column, and an exact fixed point is reached after any finite number of entity-then-time sweeps.
\end{itemize}}
\expandafter\def\csname statadic@diagnostics@STATADIC-PANEL-006\endcsname{\begin{itemize}
\item Retain original sample, weights and panel labels. Compare the achieved projection and both within-group totals to a dense least-squares baseline on tractable fits.
\item Check returned convergence, error, iterations and the weighted within-design rank independently.
\end{itemize}}
\expandafter\def\csname statadic@statement@STATADIC-DW-001\endcsname{For residuals $e_0,\ldots,e_m$ with
  $S=\sum_{i=0}^{m}e_i^2>0$, the Durbin--Watson ratio
  $\sum_{i=0}^{m-1}(e_{i+1}-e_i)^2/S$ lies in $[0,4]$.}
\expandafter\def\csname statadic@proof@STATADIC-DW-001\endcsname{Each term satisfies $(a-b)^2\le 2a^2+2b^2$ because their
  difference is $(a+b)^2$. Summing, the shifted and unshifted squares are each
  partial sums of $S$, so the numerator is at most $4S$.}
\expandafter\def\csname statadic@assumptions@STATADIC-DW-001\endcsname{\begin{itemize}
\item A finite residual sequence e\_0, ..., e\_m with a positive sum of squares.
\end{itemize}}
\expandafter\def\csname statadic@diagnostics@STATADIC-DW-001\endcsname{\begin{itemize}
\item DurbinWatson() uses the retained row order as time order. Sort observations by time before estimating, and check that the fit is valid and has a non-zero residual sum of squares; otherwise the statistic is unavailable.
\end{itemize}}
\expandafter\def\csname statadic@statement@STATADIC-JB-001\endcsname{For $n>0$, $JB=\tfrac n6\bigl(S^2+(K-3)^2/4\bigr)\ge0$,
  with equality exactly when the skewness $S$ is $0$ and the kurtosis $K$
  is $3$.}
\expandafter\def\csname statadic@proof@STATADIC-JB-001\endcsname{The factor $n/6$ is positive and the bracket is a sum of
  nonnegative squares, so the product is zero exactly when both squares
  are zero.}
\expandafter\def\csname statadic@assumptions@STATADIC-JB-001\endcsname{\begin{itemize}
\item A positive sample size n and real sample skewness S and kurtosis K.
\end{itemize}}
\expandafter\def\csname statadic@diagnostics@STATADIC-JB-001\endcsname{\begin{itemize}
\item JarqueBera() reports the biased residual skewness and kurtosis beside the statistic. Inspect them, and a residual Q-Q plot, to see which moment drives a rejection.
\end{itemize}}
\expandafter\def\csname statadic@statement@STATADIC-INFLUENCE-001\endcsname{Let $m>1$ be the residual degrees of freedom, $s^2>0$ the
  residual variance, $h<1$ a leverage and $e$ the residual. Write
  $r^2=e^2/(s^2(1-h))$ and $d=(ms^2-e^2/(1-h))/(m-1)$ for the deleted-row
  variance. If $d>0$, then $e^2/(d(1-h))=r^2(m-1)/(m-r^2)$.}
\expandafter\def\csname statadic@proof@STATADIC-INFLUENCE-001\endcsname{Multiply the positivity of $d$ by $m-1$ to obtain
  $ms^2(1-h)>e^2$, hence $m>r^2$. Then $d=s^2(m-r^2)/(m-1)$, and clearing the
  positive denominators reduces both sides to the same rational expression.}
\expandafter\def\csname statadic@assumptions@STATADIC-INFLUENCE-001\endcsname{\begin{itemize}
\item Residual degrees of freedom m > 1, positive residual variance s\textasciicircum{}2 and leverage h < 1.
\item The deleted-row variance d = (m s\textasciicircum{}2 - e\textasciicircum{}2/(1-h))/(m-1) is positive.
\end{itemize}}
\expandafter\def\csname statadic@diagnostics@STATADIC-INFLUENCE-001\endcsname{\begin{itemize}
\item Influence() is unavailable when any leverage is within sqrt(machine epsilon) of one or when fewer than p + 2 rows remain. Inspect leverage and Cook's distance together; a large value marks sensitivity, not an invalid row.
\end{itemize}}
\newcommand{\statadicstatement}[1]{\ifcsname statadic@statement@#1\endcsname\csname statadic@statement@#1\endcsname\else\PackageError{statadic}{Unknown statement #1}{Regenerate certified fragments}\fi}
\newcommand{\statadicproof}[1]{\ifcsname statadic@proof@#1\endcsname\csname statadic@proof@#1\endcsname\else\PackageError{statadic}{Unknown proof #1}{Regenerate certified fragments}\fi}
\newcommand{\statadicassumptions}[1]{\ifcsname statadic@assumptions@#1\endcsname\csname statadic@assumptions@#1\endcsname\else\PackageError{statadic}{Unknown assumptions #1}{Regenerate certified fragments}\fi}
\newcommand{\statadicdiagnostics}[1]{\ifcsname statadic@diagnostics@#1\endcsname\csname statadic@diagnostics@#1\endcsname\else\PackageError{statadic}{Unknown diagnostics #1}{Regenerate certified fragments}\fi}
